% Consolidation des acquis de la dissertation philosophique — Philosophie Tle Tech — Cours signé OnBuch+
% Programme officiel MINESEC. Compiler avec : tectonic N02-consolidation-acquis-dissertation-philosophique.tex
\newcommand{\DOCMATIERE}{Philosophie}
\newcommand{\DOCNIVEAU}{Tle Tech}
\newcommand{\DOCMODULE}{Philosophie – classe de Terminale technique}
\newcommand{\DOCLECON}{N02}
\newcommand{\DOCTITRE}{Consolidation des acquis de la dissertation philosophique}
% OnBuch+ — préambule « cours signés » ENRICHI (compilé avec tectonic / XeLaTeX)
% Système visuel « Pop » d'OnBuch+ : fond crème (jamais blanc), bordures encre
% épaisses, ombres « dures » décalées, titres Archivo Black en capitales.
% Version MATHÉMATIQUES : graphes (pgfplots/tikz), boîtes pédagogiques
% (définition, propriété, méthode, attention, le savais-tu, exercice résolu,
% exercice, TP, bilan), page de garde et signature OnBuch+ ancrée en bas.
%
% Macros attendues dans main.tex AVANT \input{preamble} :
%   \DOCMATIERE  \DOCNIVEAU  \DOCMODULE  \DOCLECON (numéro)  \DOCTITRE
\documentclass[11pt]{article}

\usepackage{fontspec}
\usepackage[french]{babel}
\usepackage[a4paper,margin=2cm,top=3.4cm,bottom=2.8cm,headheight=1.4cm,footskip=1.3cm]{geometry}
\usepackage{amsmath,amssymb,amsfonts}
\usepackage{mathtools} % \xmapsto, \coloneqq... souvent utilisés par les modèles
\usepackage{cancel} % simplifications barrées (bilans de forces)
% \abs{z} (module, valeur absolue) et \norm{u} (norme), utilisés par les modèles
\providecommand{\abs}{}\let\abs\relax\DeclarePairedDelimiter{\abs}{\lvert}{\rvert}
\providecommand{\norm}{}\let\norm\relax\DeclarePairedDelimiter{\norm}{\lVert}{\rVert}
\usepackage[table]{xcolor} % [table] : fournit aussi \rowcolors (alternance de lignes)
\usepackage{pagecolor}
\usepackage{tikz}
\usetikzlibrary{arrows.meta,positioning,calc,shapes.geometric,shapes.misc,shapes.arrows,decorations.pathmorphing,decorations.pathreplacing,decorations.markings,patterns,shadows,mindmap,trees,fit,backgrounds,matrix,intersections,angles,quotes}
\usepackage{pgfplots}
\usepackage[version=4]{mhchem} % équations nucléaires \ce{^{235}_{92}U}
\usepackage[european,siunitx]{circuitikz} % schémas électriques
\pgfplotsset{compat=1.18}
\usepgfplotslibrary{fillbetween,groupplots} % aires entre courbes (calcul intégral)
\usepackage{enumitem}
\usepackage{fancyhdr}
% [most] charge toutes les bibliothèques tcolorbox (listings, minted, raster,
% documentation...) et gonfle inutilement la mémoire TeX sur un document long
% (~300 boîtes) : on ne charge que ce qui est réellement utilisé.
\usepackage[skins,breakable]{tcolorbox}
\usepackage{graphicx}
\usepackage{colortbl}
\usepackage{booktabs}
\usepackage{tabularx}
\usepackage{multirow}
\usepackage{diagbox} % case d'en-tête diagonale des tableaux à double entrée
% Colonnes extensibles centrée / gauche / droite (C, L, R), souvent utilisées par les modèles
\newcolumntype{C}{>{\centering\arraybackslash}X}
\newcolumntype{L}{>{\raggedright\arraybackslash}X}
\newcolumntype{R}{>{\raggedleft\arraybackslash}X}
\usepackage{array}
\usepackage{multicol}
\usepackage{siunitx}
% parse-numbers=false : accepte aussi \SI{2,5 \times 10^{-3}}{...}
\sisetup{locale=FR,output-decimal-marker={,},group-minimum-digits=4,parse-numbers=false}
\DeclareSIUnit\ohm{\ensuremath{\symup{\Omega}}} % U+2126 absent de la police
\DeclareSIUnit\division{div} % réglages d'oscilloscope (ms/div, V/div)
\DeclareSIUnit\tr{tr} % tours (vitesse de rotation, ex. \SI{1500}{\tr\per\minute})
\DeclareSIUnit\molar{M} % concentration molaire (ex. \SI{3}{\molar}, \SI{1}{\micro\molar})
\DeclareSIUnit\an{an} % année (le modèle confond parfois avec \year, un compteur TeX) — ex. \SI{5}{\kilo\meter\per\an}
\DeclareSIUnit\FCFA{FCFA} % franc CFA (ex. \SI{500}{\FCFA\per\kilo\gram})
\DeclareSIUnit\degreeBrix{\textdegree Brix} % teneur en sucre d'un sirop/jus (réfractométrie)
\DeclareSIUnit\month{mois} % le modèle utilise parfois \month, un compteur TeX, comme unité de durée
\DeclareSIUnit\microliter{\micro\liter} % le modèle écrit parfois \microliter en un seul token
\DeclareSIUnit\million{millions} % dénombrement (ex. \SI{15}{\million\per\mL} spermatozoïdes)
\usepackage{qrcode}
% \degree : commande fréquemment utilisée par les modèles pour un angle ou une
% température en dehors de \SI{}{} (non fournie par les paquets chargés).
\providecommand{\degree}{^{\circ}}
% \qed (fin de démonstration), utilisé par les modèles sans amsthm
\providecommand{\qed}{\hfill$\square$}
% \barre : double barre des valeurs interdites dans un tableau de variations (tkz-tab)
\providecommand{\barre}{\ensuremath{\|}}
% environnement proof (amsthm non chargé) utilisé par les modèles
\ifdefined\proof\else\newenvironment{proof}[1][Démonstration]{\par\noindent\textit{#1.}\ }{\qed\par}\fi
\usepackage{lastpage}
\usepackage{hyperref}
\hypersetup{hidelinks}

% ============================================================
% Palette « Pop » OnBuch+ — mêmes hex que la classe P (plus_ui.dart)
% ============================================================
\definecolor{poporange}{HTML}{F59321}
\definecolor{popdark}{HTML}{E07A0C}
\definecolor{popcream}{HTML}{FFF6E8}
\definecolor{popink}{HTML}{1C1714}
\definecolor{poppurple}{HTML}{7A5AE0}
\definecolor{popgreen}{HTML}{2E9E6A}
\definecolor{poppink}{HTML}{E0457B}
\definecolor{popblue}{HTML}{2F7DE1}
\definecolor{popgold}{HTML}{C98A12}
\definecolor{popmuted}{HTML}{8A7F76}
\definecolor{popline}{HTML}{EADFD2}
\definecolor{popgray}{HTML}{8A7F76}
\colorlet{popgrey}{popgray}
% Alias tolérants (le modèle nomme parfois les couleurs différemment)
\colorlet{popwhite}{white}
\colorlet{popblack}{popink}
\colorlet{popred}{poppink}
\colorlet{popyellow}{popgold}
% Teintes claires (fonds de tableaux/encadrés) — le modèle nomme parfois
% les couleurs avec un suffixe L (ex. popblueL) sans les avoir définies.
\colorlet{poporangeL}{poporange!20}
\colorlet{popdarkL}{popdark!20}
\colorlet{poppurpleL}{poppurple!20}
\colorlet{popgreenL}{popgreen!20}
\colorlet{poppinkL}{poppink!20}
\colorlet{popblueL}{popblue!20}
\colorlet{popgoldL}{popgold!20}
\colorlet{popredL}{popred!20}
\colorlet{popgrayL}{popgray!20}
% Teintes claires pour fonds de tableaux / zones
\colorlet{popblueL}{popblue!10!white}
\colorlet{poporangeL}{poporange!14!white}
\colorlet{popgreenL}{popgreen!12!white}
\colorlet{poppurpleL}{poppurple!12!white}
\colorlet{poppinkL}{poppink!12!white}
\colorlet{popgoldL}{popgold!14!white}

\pagecolor{popcream}

% ============================================================
% Typographie
% ============================================================
\defaultfontfeatures{Ligatures=TeX}
\setmainfont{PlusJakartaSans-Medium.ttf}[Path=../../../fonts/,BoldFont=PlusJakartaSans-Bold.ttf]
% Maths en sans-serif (Fira Math) avec chiffres et lettres droites de Plus
% Jakarta, pour que les formules se fondent dans le texte.
\usepackage[math-style=ISO,bold-style=ISO]{unicode-math}
\setmathfont{FiraMath-Regular.otf}[Path=../../../fonts/]
\setmathfont{PlusJakartaSans-Medium.ttf}[Path=../../../fonts/,range={up/num,up/{Latin,latin},\mathup}]
\usepackage{icomma} % virgule décimale sans espace en mode math
% Caractères mathématiques Unicode absents des polices de texte (Plus Jakarta,
% Archivo Black) : rendus par leur équivalent mathématique, en texte comme en
% formule — sinon ils s'affichent en carré vide (ex. « Arithmétique dans ℤ »).
\usepackage{newunicodechar}
\newunicodechar{ℝ}{\ensuremath{\mathbb{R}}}
\newunicodechar{ℤ}{\ensuremath{\mathbb{Z}}}
\newunicodechar{ℂ}{\ensuremath{\mathbb{C}}}
\newunicodechar{ℕ}{\ensuremath{\mathbb{N}}}
\newunicodechar{ℚ}{\ensuremath{\mathbb{Q}}}
\newunicodechar{𝔻}{\ensuremath{\mathbb{D}}}
\newunicodechar{α}{\ensuremath{\alpha}}
\newunicodechar{β}{\ensuremath{\beta}}
\newunicodechar{γ}{\ensuremath{\gamma}}
\newunicodechar{δ}{\ensuremath{\delta}}
\newunicodechar{ε}{\ensuremath{\varepsilon}}
\newunicodechar{θ}{\ensuremath{\theta}}
\newunicodechar{λ}{\ensuremath{\lambda}}
\newunicodechar{μ}{\ensuremath{\mu}}
\newunicodechar{σ}{\ensuremath{\sigma}}
\newunicodechar{τ}{\ensuremath{\tau}}
\newunicodechar{φ}{\ensuremath{\varphi}}
\newunicodechar{ψ}{\ensuremath{\psi}}
\newunicodechar{ω}{\ensuremath{\omega}}
\newunicodechar{Σ}{\ensuremath{\Sigma}}
% majuscules produites par \MakeUppercase dans les titres (α-aminés -> Α)
\newunicodechar{Α}{\ensuremath{\alpha}}
\newunicodechar{Β}{\ensuremath{\beta}}
\newunicodechar{Γ}{\ensuremath{\Gamma}}
\newunicodechar{Δ}{\ensuremath{\Delta}}
\newunicodechar{Ε}{\ensuremath{\varepsilon}}
\newunicodechar{Θ}{\ensuremath{\Theta}}
\newunicodechar{Λ}{\ensuremath{\Lambda}}
\newunicodechar{Μ}{\ensuremath{\mu}}
\newunicodechar{Τ}{\ensuremath{\tau}}
\newunicodechar{Φ}{\ensuremath{\Phi}}
\newunicodechar{Ψ}{\ensuremath{\Psi}}
\newunicodechar{Ω}{\ensuremath{\Omega}}
\newunicodechar{↦}{\ensuremath{\mapsto}}
\newunicodechar{∈}{\ensuremath{\in}}
\newunicodechar{∉}{\ensuremath{\notin}}
\newunicodechar{∧}{\ensuremath{\wedge}}
\newunicodechar{⇔}{\ensuremath{\Leftrightarrow}}
\newunicodechar{⟺}{\ensuremath{\Longleftrightarrow}}
\newunicodechar{⇒}{\ensuremath{\Rightarrow}}
\newunicodechar{⟹}{\ensuremath{\Longrightarrow}}
\newunicodechar{∘}{\ensuremath{\circ}}
\newunicodechar{∪}{\ensuremath{\cup}}
\newunicodechar{∩}{\ensuremath{\cap}}
\newunicodechar{≡}{\ensuremath{\equiv}}
\newunicodechar{⊂}{\ensuremath{\subset}}
\newunicodechar{⊥}{\ensuremath{\perp}}
\newunicodechar{∃}{\ensuremath{\exists}}
\newunicodechar{∀}{\ensuremath{\forall}}
\newunicodechar{∛}{\ensuremath{\sqrt[3]{\,}}}
\newunicodechar{∜}{\ensuremath{\sqrt[4]{\,}}}
\newunicodechar{⃗}{\ensuremath{{}^{\to}}}
\newunicodechar{ⁿ}{\ifmmode^{n}\else\textsuperscript{\ensuremath{n}}\fi}
\newunicodechar{ˣ}{\ifmmode^{x}\else\textsuperscript{\ensuremath{x}}\fi}
\newunicodechar{ᵇ}{\ifmmode^{b}\else\textsuperscript{\ensuremath{b}}\fi}
\newunicodechar{ʳ}{\ifmmode^{r}\else\textsuperscript{\ensuremath{r}}\fi}
\newunicodechar{ᵘ}{\ifmmode^{u}\else\textsuperscript{\ensuremath{u}}\fi}
\newunicodechar{ʸ}{\ifmmode^{y}\else\textsuperscript{\ensuremath{y}}\fi}
\newunicodechar{ᵃ}{\ifmmode^{a}\else\textsuperscript{\ensuremath{a}}\fi}
\newunicodechar{ᵉ}{\ifmmode^{e}\else\textsuperscript{\ensuremath{e}}\fi}
\newunicodechar{ᵏ}{\ifmmode^{k}\else\textsuperscript{\ensuremath{k}}\fi}
\newunicodechar{ᵖ}{\ifmmode^{p}\else\textsuperscript{\ensuremath{p}}\fi}
\newunicodechar{ᴺ}{\ifmmode^{N}\else\textsuperscript{\ensuremath{N}}\fi}
\newunicodechar{ᶜ}{\ifmmode^{c}\else\textsuperscript{\ensuremath{c}}\fi}
\newunicodechar{ʰ}{\ifmmode^{h}\else\textsuperscript{\ensuremath{h}}\fi}
\newunicodechar{ᵠ}{\ifmmode^{q}\else\textsuperscript{\ensuremath{q}}\fi}
\newunicodechar{ₙ}{\ifmmode_{n}\else\textsubscript{\ensuremath{n}}\fi}
\newunicodechar{ₐ}{\ifmmode_{a}\else\textsubscript{\ensuremath{a}}\fi}
\newunicodechar{ᵢ}{\ifmmode_{i}\else\textsubscript{\ensuremath{i}}\fi}
\newunicodechar{ᵣ}{\ifmmode_{r}\else\textsubscript{\ensuremath{r}}\fi}
\newunicodechar{ₖ}{\ifmmode_{k}\else\textsubscript{\ensuremath{k}}\fi}
\newunicodechar{ᵦ}{\ifmmode_{\beta}\else\textsubscript{\ensuremath{\beta}}\fi}
\newunicodechar{ₓ}{\ifmmode_{x}\else\textsubscript{\ensuremath{x}}\fi}
\newfontfamily\popdisplay{ArchivoBlack-Regular.ttf}[Path=../../../fonts/]
\newfontfamily\poplogo{SpaceGrotesk-Bold.ttf}[Path=../../../fonts/]
\newfontfamily\popsemi{PlusJakartaSans-SemiBold.ttf}[Path=../../../fonts/]
\newfontfamily\popxbold{PlusJakartaSans-ExtraBold.ttf}[Path=../../../fonts/]
\newfontfamily\popxboldls{PlusJakartaSans-ExtraBold.ttf}[Path=../../../fonts/,LetterSpace=3.0]

\setlist[enumerate]{leftmargin=1.7em,itemsep=5pt,topsep=4pt}
\setlist[itemize]{leftmargin=1.7em,itemsep=4pt,topsep=4pt,label={\color{poporange}\textbullet}}
\setlength{\parindent}{0pt}
\setlength{\parskip}{5pt}
\renewcommand{\arraystretch}{1.25}

% pgfplots : style Pop par défaut
\pgfplotsset{
  every axis/.append style={axis line style={popink,line width=1pt},tick style={popink},
    grid style={popline,line width=0.5pt},label style={font=\small},tick label style={font=\footnotesize}},
  popcurve/.style={poporange,line width=1.8pt,no markers,smooth},
}
% Même style disponible dans un \draw TikZ simple (hors pgfplots)
\tikzset{popcurve/.style={poporange,line width=1.8pt}}
\tikzset{popgrid/.style={popline,very thin}}
\colorlet{popcurve}{poporange} % le nom du style est parfois utilisé comme couleur

% ============================================================
% Pill / squiggle
% ============================================================
\newcommand{\poppill}[3]{%
  \tikz[baseline=-0.55ex]\node[draw=popink,line width=1.2pt,rounded corners=2.6pt,%
    fill=#2,inner xsep=6.5pt,inner ysep=3pt]{%
    \popxboldls\fontsize{7}{7}\selectfont\color{#3}\MakeUppercase{#1}};%
}
\newcommand{\popsquiggle}[1]{%
  \tikz[baseline=-0.4ex]{
    \draw[#1,line width=2.6pt,line cap=round]
      plot [smooth] coordinates {(0,0.09) (0.34,0.01) (0.68,0.21) (1.02,0.01) (1.36,0.10)};
  }%
}
% Mot-clé surligné (marqueur orange sous le texte)
\newcommand{\cle}[1]{{\popxbold\color{popdark}#1}}

% ============================================================
% En-tête / pied
% ============================================================
\pagestyle{fancy}
\fancyhf{}
\renewcommand{\headrulewidth}{0pt}
\renewcommand{\footrulewidth}{0pt}
\lhead{\raisebox{-0.15ex}{\poplogo\fontsize{13}{13}\selectfont\color{popink}OnBuch\color{poporange}+}}
\rhead{\poppill{\DOCMATIERE\ \textbullet\ \DOCNIVEAU\ \textbullet\ Leçon \DOCLECON}{white}{popink}}
\lfoot{\popxbold\fontsize{7.6}{7.6}\selectfont\color{popmuted}\MakeUppercase{Cours signé OnBuch+}\ \textbullet\ {\popsemi\DOCTITRE}}
\rfoot{\tikz[baseline=-0.6ex]\node[rounded corners=8.5pt,fill=popink,minimum height=17pt,minimum width=17pt,inner xsep=5pt,inner sep=0]{\popxbold\fontsize{8}{8}\selectfont\color{popcream}\thepage\,/\,\pageref*{LastPage}};}

% ============================================================
% Titres
% ============================================================
\newcommand{\chaptitre}[2]{%
  \begin{tcolorbox}[enhanced,colback=poporange,colframe=popink,boxrule=2.6pt,arc=11pt,
      drop shadow=popink,left=15pt,right=15pt,top=12pt,bottom=11pt,before skip=3pt,after skip=18pt]
    \poppill{#2}{popink}{popcream}\par\vspace{9pt}
    {\raggedright\hyphenpenalty=10000\exhyphenpenalty=10000\popdisplay\fontsize{22}{24}\selectfont\color{popcream}\MakeUppercase{#1}\par}
  \end{tcolorbox}
}
% Section numérotée : pastille orange avec le numéro + titre Archivo
\newcounter{popsec}
\newcommand{\coursec}[1]{%
  \stepcounter{popsec}\setcounter{popsub}{0}%
  \par\vspace{16pt}\noindent
  \begin{minipage}{\linewidth}
  \tikz[baseline=-0.7ex]\node[draw=popink,line width=1.6pt,fill=poporange,rounded corners=4pt,minimum size=22pt,inner sep=0]{\popdisplay\fontsize{12}{12}\selectfont\color{popcream}\thepopsec};%
  \hspace{8pt}{\popdisplay\fontsize{15}{17}\selectfont\color{popink}\MakeUppercase{#1}}\par
  \vspace{4pt}\noindent{\color{poporange}\rule{36pt}{3.6pt}}
  \end{minipage}\par\nopagebreak\vspace{8pt}
}
\newcounter{popsub}
\newcommand{\courssub}[1]{%
  \stepcounter{popsub}%
  \par\vspace{9pt}\noindent{\popxbold\fontsize{11.5}{13}\selectfont\color{popink}{\color{poporange}\thepopsec.\thepopsub}\hspace{6pt}#1}\par\nopagebreak\vspace{4pt}
}

% ============================================================
% Boîtes pédagogiques — même recette : fond blanc, bordure épaisse colorée,
% ombre dure encre, badge plein. Titre optionnel [#1].
% ============================================================
\tcbset{popbase/.style={breakable,enhanced,colback=white,boxrule=2.3pt,arc=9pt,
  shadow={3pt}{-3pt}{0pt}{popink},left=14pt,right=14pt,top=11pt,bottom=11pt,before skip=11pt,after skip=15pt,
  fontupper=\popsemi\fontsize{10.6}{14.5}\selectfont\color{popink},
  fontlower=\popsemi\fontsize{10.6}{14.5}\selectfont\color{popink}}}
% Coche dessinée (le glyphe ✓ n'existe pas dans les polices du document)
\newcommand{\popcheck}[1]{\tikz[baseline=-0.5ex]\draw[#1,line width=1.6pt,line cap=round,line join=round](-0.13,0.0)--(-0.04,-0.09)--(0.13,0.1);}
\providecommand{\coche}{\popcheck{popgreen}}
\providecommand{\resultat}[1]{\par\smallskip\noindent\fcolorbox{popgreen}{popgreenL}{\parbox{\dimexpr\linewidth-2\fboxsep-2\fboxrule\relax}{\textbf{Résultat.} #1}}\par\smallskip}
\newcommand{\popboxhead}[3]{\poppill{#1}{#2}{white}\ifx\relax#3\relax\else\hspace{6pt}{\popxbold\fontsize{10.6}{12}\selectfont\color{popink}#3}\fi\par\vspace{7pt}}

\newtcolorbox{aretenir}[1][]{popbase,colframe=popgreen,before upper={\popboxhead{À retenir}{popgreen}{#1}}}
\newtcolorbox{exemplebox}[1][]{popbase,colframe=poporange,before upper={\popboxhead{Exemple}{poporange}{#1}}}
\newtcolorbox{definition}[1][]{popbase,colframe=popblue,before upper={\popboxhead{Définition}{popblue}{#1}}}
\newtcolorbox{propriete}[1][]{popbase,colframe=poppurple,before upper={\popboxhead{Propriété}{poppurple}{#1}}}
\newtcolorbox{methode}[1][]{popbase,colframe=popgold,before upper={\popboxhead{Méthode}{popgold}{#1}}}
\newtcolorbox{attention}[1][]{popbase,colframe=poppink,before upper={\popboxhead{Attention}{poppink}{#1}}}
\newtcolorbox{experience}[1][]{popbase,colframe=popblue,colback=popblueL,before upper={\popboxhead{Expérience}{popblue}{#1}}}
% « Le savais-tu ? » : carte violette pleine (contexte, culture, Cameroun)
\newtcolorbox{savaistu}[1][]{popbase,colback=poppurple,colframe=popink,
  fontupper=\popsemi\fontsize{10.4}{14.2}\selectfont\color{white},
  before upper={\poppill{Le savais-tu\,?}{white}{poppurple}\ifx\relax#1\relax\else\hspace{6pt}{\popxbold\color{white}#1}\fi\par\vspace{7pt}}}
% Exercice résolu : énoncé en haut, \tcblower puis la solution sur fond crème
\newcounter{popexo}\newcounter{popexr}
\newtcolorbox{exoresolu}[1][]{popbase,colframe=poporange,colbacklower=popcream,
  segmentation style={popink,line width=1.2pt,dashed},
  before upper={\stepcounter{popexr}\popboxhead{Exercice résolu \thepopexr}{poporange}{#1}},
  before lower={\poppill{Solution}{popink}{popcream}\par\vspace{6pt}}}
% Exercice (non résolu en place) — la correction est dans la partie Corrigés
\newtcolorbox{exercice}[1][]{popbase,colframe=popink,
  before upper={\stepcounter{popexo}\popboxhead{Exercice \thepopexo}{popink}{#1}}}
% Corrigé
\newtcolorbox{corrige}[1][]{popbase,colframe=popgreen,colback=popgreenL,
  before upper={\popboxhead{Corrigé}{popgreen}{#1}}}
% Objectifs / prérequis / bilan
\newtcolorbox{objectifs}{popbase,colframe=popink,colback=poporangeL,
  before upper={\popboxhead{Objectifs de la leçon}{popink}{}}}
\newtcolorbox{prerequis}{popbase,colframe=popink,colback=popblueL,
  before upper={\popboxhead{Prérequis}{popblue}{}}}
\newtcolorbox{bilan}{popbase,colframe=popink,colback=white,boxrule=2.8pt,
  before upper={\popboxhead{Fiche bilan}{poporange}{L'essentiel à connaître pour l'examen}}}
% Figure encadrée avec légende
\newtcolorbox{popfigure}[1][]{enhanced,breakable=false,colback=white,colframe=popline,boxrule=1.4pt,arc=8pt,
  left=10pt,right=10pt,top=10pt,bottom=8pt,before skip=10pt,after skip=12pt,halign=center,
  lower separated=false,halign lower=center,
  fontlower=\popsemi\fontsize{9}{11.5}\selectfont\color{popmuted},
  #1}
\newcommand{\legende}[1]{\par\vspace{6pt}{\popsemi\fontsize{9}{11.5}\selectfont\color{popmuted}#1}}

% ============================================================
% Signature OnBuch+
% ============================================================
\newcommand{\poplogoinline}[1]{{\poplogo\fontsize{#1}{#1}\selectfont\color{popink}OnBuch\color{poporange}+}}
\newcommand{\signatureonbuch}{%
  \begin{tcolorbox}[enhanced,colback=popink,colframe=popink,boxrule=2.2pt,arc=10pt,
      drop shadow=poporange,left=14pt,right=14pt,top=10pt,bottom=10pt,before skip=0pt,after skip=0pt]
    \tikz[baseline=-0.7ex]\node[circle,fill=popgreen,draw=popcream,line width=1.3pt,minimum size=22pt,inner sep=0]{\popcheck{white}};%
    \hspace{9pt}\begin{minipage}[c]{0.62\linewidth}
      {\popxbold\fontsize{12}{13}\selectfont\color{popcream}Signé\ \textbullet\ {\poplogo OnBuch\color{poporange}+}}\par\vspace{2pt}
      {\raggedright\popsemi\fontsize{8}{10.5}\selectfont\color{popline}\MakeUppercase{Équipe pédagogique OnBuch+}\\\MakeUppercase{Conforme au programme officiel MINESEC}\par}
    \end{minipage}\hfill
    {\poplogo\fontsize{22}{22}\selectfont\color{popcream}OnBuch\color{poporange}+}
  \end{tcolorbox}%
}

% ============================================================
% Page de garde
% #1 titre · #2 matière · #3 niveau · #4 module · #5 n° leçon · #6 durée ·
% #7 description / objectifs (texte ou itemize)
% ============================================================
\newcommand{\pagegarde}[7]{%
  \thispagestyle{empty}
  \newgeometry{a4paper,margin=1.7cm,top=1.6cm,bottom=1.4cm}%
  \begin{tikzpicture}[remember picture,overlay]
    \fill[poporange!18] ([xshift=-3.2cm,yshift=-3.6cm]current page.north east) circle (4.6cm);
    \fill[poppurple!14] ([xshift=2.4cm,yshift=7.6cm]current page.south west) circle (3.8cm);
  \end{tikzpicture}%
  {\poplogo\fontsize{30}{30}\selectfont\color{popink}OnBuch\color{poporange}+}\hfill
  \raisebox{0.6ex}{\poppill{Cours signé \textbullet\ Premium}{popink}{popcream}}\par
  \vspace{1pt}\popsquiggle{poppurple}\par
  \vspace{8pt}
  \begin{tcolorbox}[enhanced,colback=poporange,colframe=popink,boxrule=2.8pt,arc=13pt,
      drop shadow=popink,left=18pt,right=18pt,top=12pt,bottom=13pt,after skip=10pt]
    \poppill{#2 \textbullet\ #3}{popink}{popcream}\hspace{5pt}\poppill{#4}{white}{popink}\par\vspace{8pt}
    {\popdisplay\fontsize{14}{14}\selectfont\color{popink}LEÇON #5}\par\vspace{4pt}
    {\raggedright\hyphenpenalty=10000\exhyphenpenalty=10000\popdisplay\fontsize{25}{27}\selectfont\color{popcream}\MakeUppercase{#1}\par}
    \vspace{8pt}
    {\popxbold\fontsize{9.5}{9.5}\selectfont\color{popink}\MakeUppercase{Durée conseillée : #6}}
  \end{tcolorbox}
  \begin{tcolorbox}[enhanced,colback=white,colframe=popink,boxrule=2.2pt,arc=10pt,
      drop shadow=popink,left=16pt,right=16pt,top=10pt,bottom=10pt,after skip=8pt]
    \poppill{Dans cette leçon}{poppurple}{white}\par\vspace{5pt}
    \popsemi\fontsize{9.3}{12.6}\selectfont\color{popink}#7
  \end{tcolorbox}
  \noindent\begin{minipage}[t]{0.487\linewidth}
    \begin{tcolorbox}[enhanced,colback=popgreen,colframe=popink,boxrule=2.2pt,arc=10pt,
        drop shadow=popink,left=11pt,right=11pt,top=10pt,bottom=10pt,height=3.1cm,valign=center]
      \tcbox[colback=white,colframe=popink,boxrule=1.3pt,arc=5pt,boxsep=0pt,nobeforeafter,
        left=3pt,right=3pt,top=3pt,bottom=3pt,valign=center]{\qrcode[height=1.9cm]{https://play.google.com/store/apps/details?id=com.luvvix.onbuch&hl=fr}}%
      \hspace{8pt}\begin{minipage}[c]{0.5\linewidth}
        {\popdisplay\fontsize{11}{12.5}\selectfont\color{white}\MakeUppercase{Télécharge OnBuch}}\par\vspace{4pt}
        {\raggedright\popsemi\fontsize{8.4}{11}\selectfont\color{white}Gratuit \textbullet\ Google Play\\Tuteur IA, annales, concours, résultats\par}
      \end{minipage}
    \end{tcolorbox}
  \end{minipage}\hfill
  \begin{minipage}[t]{0.487\linewidth}
    \begin{tcolorbox}[enhanced,colback=poppurple,colframe=popink,boxrule=2.2pt,arc=10pt,
        drop shadow=popink,left=11pt,right=11pt,top=10pt,bottom=10pt,height=3.1cm,valign=center]
      \tikz[baseline=-0.7ex]\node[circle,fill=white,draw=popink,line width=1.3pt,minimum size=1.6cm,inner sep=0]{\popdisplay\fontsize{24}{24}\selectfont\color{poppurple}+};%
      \hspace{8pt}\begin{minipage}[c]{0.6\linewidth}
        {\popdisplay\fontsize{11}{12.5}\selectfont\color{white}\MakeUppercase{Passe à OnBuch+}}\par\vspace{4pt}
        {\raggedright\popsemi\fontsize{8.4}{11}\selectfont\color{white}Tous les cours signés, Tuteur IA illimité, fiches PDF — onglet OnBuch+ de l'app\par}
      \end{minipage}
    \end{tcolorbox}
  \end{minipage}
  \vspace{8pt}
  \popcontenu
  \vfill
  \signatureonbuch
  \restoregeometry
  \newpage
}

% Rangée « Ce cours contient » (page de garde)
\newcommand{\poptile}[3]{% #1 couleur #2 gros texte #3 légende
  \begin{tcolorbox}[enhanced,colback=#1,colframe=popink,boxrule=2pt,arc=9pt,shadow={2.5pt}{-2.5pt}{0pt}{popink},
      left=6pt,right=6pt,top=9pt,bottom=9pt,halign=center,valign=center,height=1.75cm,width=\linewidth,nobeforeafter]
    {\popdisplay\fontsize{12.5}{13}\selectfont\color{white}\MakeUppercase{#2}}\par\vspace{3pt}
    {\centering\popsemi\fontsize{7.8}{9.5}\selectfont\color{white}#3\par}
  \end{tcolorbox}}
\newcommand{\popcontenu}{%
  \par\noindent{\popxboldls\fontsize{7.5}{7.5}\selectfont\color{popmuted}CE COURS SIGNÉ CONTIENT}\par\vspace{7pt}
  \noindent\begin{minipage}[t]{0.235\linewidth}\poptile{popblue}{Cours}{complet, illustré, pas à pas}\end{minipage}\hfill
  \begin{minipage}[t]{0.235\linewidth}\poptile{poporange}{Méthodes}{et exercices résolus}\end{minipage}\hfill
  \begin{minipage}[t]{0.235\linewidth}\poptile{poppink}{Exercices}{type examen + corrigés}\end{minipage}\hfill
  \begin{minipage}[t]{0.235\linewidth}\poptile{popgreen}{Fiche bilan}{l'essentiel à retenir}\end{minipage}\par}

% Sommaire
\newtcolorbox{sommaire}{enhanced,colback=white,colframe=popink,boxrule=2.2pt,arc=10pt,
  drop shadow=popink,left=18pt,right=18pt,top=13pt,bottom=13pt,before skip=6pt,after skip=20pt,
  before upper={\poppill{Sommaire}{poppurple}{white}\par\vspace{9pt}\popsemi\fontsize{10.6}{14.5}\selectfont\color{popink}}}

% Signature de fin de cours — poussée tout en bas de la dernière page
\newcommand{\findecours}{%
  \par\vfill\nopagebreak
  \begin{center}\popsquiggle{poporange}\end{center}\vspace{4pt}
  \signatureonbuch
}
\AtBeginDocument{\def\llbracket{\mathopen{[\mkern-3.5mu[}}\def\rrbracket{\mathclose{]\mkern-3.5mu]}}} % intervalles d'entiers (glyphe ⟦ absent de la police)
% Glyphes absents de Fira Math : redéfinis à partir de glyphes disponibles
\AtBeginDocument{%
  \def\setminus{\mathbin{\backslash}}\let\smallsetminus\setminus
  \def\vdots{\mathord{\vbox{\baselineskip3.2pt\lineskiplimit0pt\kern4pt\hbox{.}\hbox{.}\hbox{.}}}}%
  \def\ddots{\mathinner{\mkern1mu\raise6.5pt\hbox{.}\mkern2mu\raise3.6pt\hbox{.}\mkern2mu\raise0.7pt\hbox{.}\mkern1mu}}%
  \def\triangle{\mathord{\scalebox{0.9}{$\symup{\Delta}$}}}%
  \def\nabla{\mathord{\raisebox{\depth}{\scalebox{1}[-1]{$\symup{\Delta}$}}}}%
  \def\checkmark{\popcheck{popdark}}%
  \def\star{\mathbin{\ast}}\def\bigstar{\mathord{\ast}}%
  \def\diamond{\mathbin{\scalebox{0.6}{$\Diamond$}}}%
  \def\bigcup{\mathop{\mathchoice{\raisebox{-0.3em}{\scalebox{1.7}{$\cup$}}}{\scalebox{1.25}{$\cup$}}{\cup}{\cup}}\displaylimits}%
  \def\bigcap{\mathop{\mathchoice{\raisebox{-0.3em}{\scalebox{1.7}{$\cap$}}}{\scalebox{1.25}{$\cap$}}{\cap}{\cap}}\displaylimits}%
  \def\lesseqqgtr{\mathrel{\lessgtr}}\def\bigoplus{\mathop{\oplus}\displaylimits}%
}
\providecommand{\ssi}{si et seulement si} % abréviation fréquente
\AtBeginDocument{\providecommand{\wideparen}{\overparen}} % arc de cercle

%% FIGFIX-BEGIN v4 — correctifs globaux des figures (docs/onbuch-generation/tools/figfix)
\usepackage{adjustbox}\usepackage{etoolbox}
% toute figure TikZ/pgfplots plus large que la ligne est réduite à la largeur disponible
\BeforeBeginEnvironment{tikzpicture}{\begin{adjustbox}{max width=\linewidth}}
\AfterEndEnvironment{tikzpicture}{\end{adjustbox}}
% légendes pgfplots lisibles par-dessus les courbes
\pgfplotsset{every axis/.append style={legend style={fill=white,fill opacity=0.92,text opacity=1,draw=black!15,inner sep=3pt}}}
% étiquettes d'arêtes (syntaxe edge["..."]) avec fond blanc
\tikzset{every edge quotes/.append style={fill=white,inner sep=1.2pt}}
%% FIGFIX-END
\begin{document}

\pagegarde{Consolidation des acquis de la dissertation philosophique}{Philosophie}{Tle Tech}{Philosophie – classe de Terminale technique}{N02}{3 séances (fiche de progression harmonisée)}{Cette leçon consolide les acquis méthodologiques de la dissertation philosophique en Terminale technique. Elle outille l'élève pour maîtriser chaque étape — de l'introduction à la conclusion en passant par l'articulation thèse/antithèse/synthèse et l'usage critique des citations — afin de répondre aux exigences du Baccalauréat.
\vspace{4pt}
\begin{multicols}{2}\small
\begin{itemize}
\item Analyser un sujet de dissertation pour en dégager les termes clés, la tension conceptuelle et la problématique pertinente.
\item Construire une introduction complète (accroche, définitions, problématique, annonce de plan) en respectant le canon académique.
\item Rédiger des paragraphes argumentatifs structurés (Idée directrice, Argument, Exemple) pour la thèse et l'antithèse.
\item Élaborer une véritable synthèse dialectique qui dépasse l'opposition thèse/antithèse par une distinction ou une réconciliation conceptuelle.
\item Rédiger une conclusion en deux temps (bilan nuancé et ouverture pertinente) sans ajouter d'arguments nouveaux.
\item Sélectionner, intégrer et analyser des citations philosophiques exactes et pertinentes (auteurs au programme et penseurs africains).
\end{itemize}\end{multicols}}

\begin{sommaire}
\begin{multicols}{2}
\begin{enumerate}
\item 1. Les fondements et l'architecture globale de la dissertation
\item 2. L'introduction : la "porte d'entrée" (4 étapes canoniques)
\item 3. Le développement (1) : Thèse et Antithèse — L'art du paragraphe argumentatif (I-A-E)
\item 4. Le développement (2) : La Synthèse — Le dépassement dialectique
\item 5. La conclusion : Bilan définitif et Ouverture maîtrisée
\item 6. Les citations : Règles de choix, d'intégration et de référencement
\item 7. TP Consolidation : Simulation Bac \& Méthode d'auto-évaluation
\item Activité d'intégration
\item Méthodes et exercices résolus
\item Exercices
\item Corrigés des exercices
\item Fiche bilan
\end{enumerate}
\end{multicols}
\end{sommaire}

\begin{objectifs}
\begin{itemize}
\item Analyser un sujet de dissertation pour en dégager les termes clés, la tension conceptuelle et la problématique pertinente.
\item Construire une introduction complète (accroche, définitions, problématique, annonce de plan) en respectant le canon académique.
\item Rédiger des paragraphes argumentatifs structurés (Idée directrice, Argument, Exemple) pour la thèse et l'antithèse.
\item Élaborer une véritable synthèse dialectique qui dépasse l'opposition thèse/antithèse par une distinction ou une réconciliation conceptuelle.
\item Rédiger une conclusion en deux temps (bilan nuancé et ouverture pertinente) sans ajouter d'arguments nouveaux.
\item Sélectionner, intégrer et analyser des citations philosophiques exactes et pertinentes (auteurs au programme et penseurs africains).
\item Gérer le temps de l'épreuve (4h) : 30 min de brouillon, 3h15 de rédaction, 15 min de relecture.
\item S'auto-évaluer sur grille officielle (compréhension, problématisation, argumentation, structuration, expression).
\end{itemize}
\end{objectifs}

\begin{prerequis}
\begin{itemize}
\item Distinction entre dissertation et explication de texte (acquis de 1ère).
\item Connaissance des grands repères historiques de la philosophie (Antique, Moderne, Contemporain, Africaine).
\item Maîtrise des connecteurs logiques (cause, conséquence, concession, opposition, illustration).
\item Culture générale sur les enjeux techniques, sociaux et environnementaux du Cameroun (urbanisation, agriculture, NTIC, décentralisation).
\item Notion de concept philosophique vs sens commun (travail en début d'année de Tle).
\end{itemize}
\end{prerequis}

\newpage

\coursec{Les fondements et l'architecture globale de la dissertation}

Dans une salle d'examen du Baccalauréat technique à Douala, un candidat fixe le sujet : \og La technique libère-t-elle l'homme ?\fg. Il connaît ses auteurs, mais sa copie reste blanche : par où commencer ? Comment structurer sa pensée sans se perdre ? Cette leçon transforme l'angoisse de la page blanche en une méthode rigoureuse, étape par étape, pour construire une dissertation philosophique solide, notée et réussie.

\courssub{Qu'est-ce que la dissertation philosophique ?}

\begin{definition}[La dissertation philosophique]
La dissertation philosophique est une \cle{réponse argumentée, structurée et problématisée} à une question philosophique. Elle n'est ni une récitation de cours, ni un exposé d'opinions personnelles, ni une compilation de citations. C'est un \cle{exercice de problématisation} : il s'agit de faire émerger la tension conceptuelle cachée dans le sujet, puis d'y répondre par un raisonnement dialectique (thèse / antithèse / synthèse) étayé d'exemples précis et de références philosophiques pertinentes.
\end{definition}

\begin{attention}[Différence fondamentale avec l'explication de texte]
Ne confondez pas les deux épreuves du Bac :
\begin{itemize}
    \item \textbf{Explication de texte} : Le support est un \emph{texte d'auteur} imposé. La tâche : comprendre, analyser, commenter la pensée d'un autre.
    \item \textbf{Dissertation} : Le support est un \emph{sujet} (question ou citation). La tâche : construire \emph{sa propre pensée} en réponse à ce sujet, en mobilisant sa culture philosophique comme outillage, non comme fin en soi.
\end{itemize}
\emph{Piège :} Traiter la dissertation comme un « cours sur le thème » (ex : « La technique chez Marx, chez Heidegger, chez Ellul… ») sans jamais répondre à la question précise posée.
\end{attention}

\begin{exemplebox}[Sujet : « La technique libère-t-elle l'homme ? »]
\emph{Mauvaise approche (récitation) :} « Marx dit que la technique aliène le prolétaire, Heidegger qu'elle réduit l'être à du fonds de réserve, Ellul qu'elle devient autonome… »
\emph{Bonne approche (problématisation) :} Le terme « libérer » suppose un état d'esclavage initial. La technique libère-t-elle de la peine (thèse) ou asservit-elle par ses propres lois (antithèse) ? La libération authentique ne réside-t-elle pas dans une \cle{maîtrise critique} de la technique (synthèse) ?
\end{exemplebox}

\courssub{Les 4 critères d'évaluation au Baccalauréat technique}

Le correcteur note sur \textbf{20 points} répartis en quatre compétences invariantes. Les connaître, c'est savoir où placer ses efforts.

\begin{propriete}[Grille d'évaluation officielle (Bac Tech — Philosophie)]
\begin{center}
\begin{tabularx}{\linewidth}{|l|X|c|}\hline
\textbf{Compétence} & \textbf{Descriptif} & \textbf{Points} \\ \hline
\rowcolor{popblueL}
\cle{Compréhension / Problématisation} & Analyse des termes, détection de la tension, formulation d'une problématique pertinente, annonce d'un plan dialectique. & 5 \\ \hline
\cle{Argumentation / Exemples} & Qualité du raisonnement (I-A-E), pertinence des exemples (concrets, littéraires, scientifiques, historiques), justesse des références philosophiques. & 6 \\ \hline
\rowcolor{popgreenL}
\cle{Structuration / Articulation} & Respect de l'architecture (Intro / 3 parties / Conclusion), transitions logiques, paragraphes unifiés (une idée directrice), connecteurs logiques. & 4 \\ \hline
\cle{Expression / Orthographe} & Correction syntaxique, orthographique, richesse lexicale, style philosophique (précision, nuances, absence de familier). & 5 \\ \hline
\end{tabularx}
\end{center}
\end{propriete}

\begin{aretenir}[Priorité stratégique]
L'\cle{argumentation} (6 pts) et la \cle{problématisation} (5 pts) valent à elles seules 11 points. Une copie bien structurée mais vide de contenu réel ne dépasse pas 8/20. Une copie brillante mais mal structurée plafonne à 11/20. L'équilibre est la clé.
\end{aretenir}

\begin{exoresolu}[Analyse d'une copie type « 8/20 »]
\textbf{Énoncé :} Au Bac Tech 2023, un candidat traite « Le travail est-il une contrainte ? ». Il écrit 4 pages. Son introduction définit le travail (Larousse), cite Marx et Hannah Arendt, mais ne pose \emph{aucune problématique} : « Nous verrons si le travail est une contrainte ou une libération. » Son développement aligne 3 paragraphes : 1. Le travail contraint (Marx), 2. Le travail libère (Arendt), 3. Le travail aujourd'hui (statistiques OIT). Pas de transitions, pas de synthèse. Conclusion : « En conclusion, le travail est les deux. »
\textbf{Diagnostic :} Problématisation absente (0/5), Argumentation superficielle (2/6), Structure invisible (1/4), Expression correcte (4/5). \textbf{Total : 7/20}.
\textbf{Leçon :} Sans problématique explicite et sans articulation dialectique, le savoir reste muet.
\end{exoresolu}

\courssub{L'architecture invariante : Introduction — Développement dialectique — Conclusion}

Toute dissertation de Bac suit une \cle{architecture unique}, garante de la rigueur philosophique. Ce n'est pas une convention administrative, c'est la forme même de la pensée dialectique : \textbf{poser le problème, le déployer dans sa contradiction, le dépasser}.

\begin{propriete}[Architecture canonique de la dissertation]
\begin{center}
\begin{tabular}{|c|>{\centering\arraybackslash}p{5cm}|>{\centering\arraybackslash}p{5cm}|}\hline
\textbf{Partie} & \textbf{Rôle} & \textbf{Étapes obligatoires} \\ \hline
\rowcolor{popblueL}
\textbf{Introduction} & \cle{La « porte d'entrée »} : capter, situer, problématiser, annoncer. & 1. Accroche / 2. Analyse du sujet / 3. Problématique / 4. Annonce du plan \\ \hline
\textbf{Développement} & \cle{Le « corps »} : déployer la pensée en 3 temps dialectiques. & \textbf{I. Thèse} (réponse spontanée, arguments, exemples) \\
& & \textbf{II. Antithèse} (objection radicale, arguments, exemples) \\
& & \textbf{III. Synthèse} (dépassement, nouvelle thèse, ouverture) \\ \hline
\rowcolor{popgreenL}
\textbf{Conclusion} & \cle{Le « bilan »} : clore sans répéter, ouvrir sans divaguer. & 1. Réponse synthétique / 2. Portée / 3. Ouverture maîtrisée \\ \hline
\end{tabular}
\end{center}
\end{propriete}

\begin{popfigure}
\begin{tikzpicture}[
    node distance=1.2cm and 1.5cm,
    box/.style={draw=popdark, thick, rounded corners, fill=popcream, minimum width=3.2cm, minimum height=1.4cm, align=center, font=\small},
    arrow/.style={-Stealth, thick, popblue},
    bigbox/.style={draw=popblue, thick, rounded corners, fill=popblueL, minimum width=7cm, minimum height=5.5cm, align=center},
    medbox/.style={draw=poporange, thick, rounded corners, fill=poporangeL, minimum width=5.5cm, minimum height=1.6cm, align=center},
]
% Global container
\node[bigbox] (global) {};
% Title inside global
\node[above=0.2cm of global.north, font=\bfseries\large, popdark] {Architecture globale de la dissertation};
% Input
\node[box, anchor=south, align=center] (sujet) at (global.north){SUJET\\Question ou Citation};
% Analysis/Problematique
\node[box, below=1.5cm of sujet, align=center] (prob){ANALYSE \& PROBLÉMATIQUE\\Termes, Tension, Question directrice};
% Plan
\node[box, below=1.5cm of prob, align=center] (plan){PLAN DÉTAILLÉ\\Thèse / Antithèse / Synthèse\\+ Exemples \& Références};
% Three parts inside a grouped box
\node[medbox, below=1.8cm of plan] (dev) {DÉVELOPPEMENT (3 parties × 2–3 §)};
\node[box, below left=0.8cm and 1.5cm of dev.south, align=center] (th){I. THÈSE\\Réponse immédiate\\Arguments + Exemples};
\node[box, below=0.8cm of dev.south, align=center] (ant){II. ANTITHÈSE\\Objection critique\\Arguments + Exemples};
\node[box, below right=0.8cm and 1.5cm of dev.south, align=center] (syn){III. SYNTHÈSE\\Dépassement dialectique\\Nouvelle perspective};
% Conclusion
\node[box, below=1.8cm of ant, align=center] (ccl){CONCLUSION\\Bilan + Ouverture};
% Arrows
\draw[arrow] (sujet.south) -- (prob.north);
\draw[arrow] (prob.south) -- (plan.north);
\draw[arrow] (plan.south) -- (dev.north);
\draw[arrow] (dev.south west) -- (th.north);
\draw[arrow] (dev.south) -- (ant.north);
\draw[arrow] (dev.south east) -- (syn.north);
\draw[arrow] (th.south) -- ++(0,-0.3) -| (ccl.north west);
\draw[arrow] (ant.south) -- (ccl.north);
\draw[arrow] (syn.south) -- ++(0,-0.3) -| (ccl.north east);
% Annotations
\node[left=0.3cm of prob, font=\tiny\itshape, popmuted, align=right] {Brouillon :\\Arbre conceptuel};
\node[right=0.3cm of plan, font=\tiny\itshape, popmuted, align=left] {Brouillon :\\Plan numéroté};
\node[left=0.3cm of th, font=\tiny\itshape, popmuted, align=right] {Structure § :\\I-A-E};
\node[right=0.3cm of syn, font=\tiny\itshape, popmuted, align=left] {Clé :\\« Ni l'un ni l'autre,\\mais… »};
\end{tikzpicture}
\legende{Architecture globale : du sujet à la conclusion, le cheminement logique et dialectique. Chaque flèche représente une étape de validation (brouillon → rédaction).}
\end{popfigure}

\courssub{La gestion du temps (4h) et le rôle stratégique du brouillon}

Quatre heures (240 minutes) paraissent longues ; elles passent vite si l'on rédige sans avoir pensé. La règle d'or : \textbf{30 minutes de brouillon intense valent 3 heures de rédaction fluide}.

\begin{methode}[Les 4 temps du brouillon (30 min chrono)]
\begin{enumerate}
    \item \textbf{Analyser les termes (5 min) :} Définir philosophiquement chaque mot-clé (pas le dictionnaire courant), repérer les présupposés, les oppositions internes. Ex : « Technique » $\neq$ « technologie » ; « Libérer » $\rightarrow$ de quoi ? pour qui ?
    \item \textbf{Trouver la tension / le problème (7 min) :} Pourquoi ce sujet n'est-il pas évident ? Quelle opposition (liberté/nécessité, nature/culture, moyen/fin) le traverse ? Formuler la \cle{problématique} en une question précise, ni trop large, ni trop étroite.
    \item \textbf{Construire le plan dialectique détaillé (13 min) :} Numéroter I, II, III. Pour chaque partie : 2 à 3 arguments (A, B, C), chaque argument appuyé par un \cle{exemple concret} (historique, scientifique, littéraire, vécu camerounais) et/ou une \cle{référence philosophique} (auteur + concept précis). Rédiger les \cle{phrases d'introduction} et de \cle{conclusion} de chaque grande partie. Pré-rédiger l'\cle{accroche}, la \cle{problématique}, l'\cle{annonce de plan}, la \cle{réponse finale} et l'\cle{ouverture}.
    \item \textbf{Sélectionner et vérifier les citations (5 min) :} Choisir 3 à 5 citations exactes, courtes, pertinentes. Vérifier l'auteur, l'œuvre, le sens. Les écrire en marge du brouillon pour les copier sans erreur.
\end{enumerate}
\end{methode}

\begin{popfigure}
\begin{tikzpicture}
\begin{axis}[
    width=\linewidth,
    height=7cm,
    xbar stacked,
    xmin=0, xmax=240,
    ytick={1,2,3},
    yticklabels={Relecture (15 min), Rédaction soignée (195 min), Brouillon structuré (30 min)},
    xlabel={Minutes},
    bar width=1.2cm,
    legend style={at={(0.5,-0.2)}, anchor=north, legend columns=-1, font=\small},
    axis x line*=bottom, axis y line=none,
    tickwidth=0pt,
    enlarge y limits=0.4,
    xmajorgrids=true,
    grid style={popline},
]
\addplot[popblueL, fill=popblue] coordinates {(30,1)};
\addplot[poporangeL, fill=poporange] coordinates {(195,2)};
\addplot[popgreenL, fill=popgreen] coordinates {(15,3)};
\legend{Brouillon, Rédaction, Relecture}
\end{axis}
% Annotations
\node[anchor=west, font=\small\bfseries, popblue] at (0.15\linewidth, 7.2cm) {Analyse, Problématique, Plan, Citations};
\node[anchor=west, font=\small\bfseries, poporange] at (0.15\linewidth, 5.8cm) {Intro (20') — I (45') — II (45') — III (45') — Ccl (20') — Transitions (20')};
\node[anchor=west, font=\small\bfseries, popgreen] at (0.15\linewidth, 4.4cm) {Correction orthographe, syntaxe, cohérence, citations};
\end{tikzpicture}
\legende{Chronogramme type des 240 minutes. Le brouillon (12,5\% du temps) structure 81\% de la rédaction. La relecture (6\%) sécurise les 5 points d'expression.}
\end{popfigure}

\begin{attention}[Pièges de la gestion du temps]
\begin{itemize}
    \item \textbf{Zéro minute de brouillon :} La copie « part en vrille » au milieu de la thèse ; le plan change en cours de route ; les transitions sont inexistantes.
    \item \textbf{Brouillon « brouillon » :} Mots-clés seulement, pas de phrases-clés. On perd 20 min à reformuler en propre pendant la rédaction.
    \item \textbf{Rédaction au stylo directement :} Ratures, surcharges, lisibilité dégradée $\rightarrow$ pénalité expression.
    \item \textbf{Oublier la relecture :} 5 points faciles perdus (accords, majuscules, noms d'auteurs, numérotation des parties).
\end{itemize}
\end{attention}

\begin{savaistu}[L'effet de halo du correcteur]
Les correcteurs du Bac lisent souvent l'\textbf{introduction} et la \textbf{conclusion} en premier (parfois uniquement si la copie est faible). Une introduction claire (problématique nette + plan annoncé) et une conclusion ferme (réponse + ouverture) créent un \emph{effet de halo} positif : le correcteur entre dans le développement avec une note mentale déjà haute (ex : 14/20). Une introduction floue ou une conclusion absente ancrent une note basse (ex : 8/20) avant même de lire les arguments. \textbf{Soignez ces deux moments comme s'ils valaient 10 points chacun.}
\end{savaistu}

\begin{exemplebox}[Exemple concret camerounais pour « La technique libère-t-elle l'homme ? »]
\textbf{Thèse (Libération) :} La machette, la pirogue motorisée, le téléphone mobile libèrent le paysan de la fatigue, de l'enclavement, de l'ignorance des cours. \emph{Exemple :} Le paysan de Dschang vend ses plantains via WhatsApp, évite les intermédiaires.
\textbf{Antithèse (Aliénation) :} La même technique crée dépendance (carburant, pièces, réseau), standardise les gestes, soumet au rythme machine. \emph{Exemple :} L'orpaillage artisanat au Cameroun-Est : la pompe à eau et le détecteur de métaux augmentent le rendement mais détruisent la rivière, endettent le creuseur auprès du fournisseur chinois.
\textbf{Synthèse (Maîtrise critique) :} La technique n'est ni bonne ni mauvaise ; elle est \cle{pharmakon} (remède/poison). La libération réelle = \emph{appropriation technique} : formation locale à la maintenance, choix souverains d'innovation (ex : fours à bois améliorés conçus à l'ENSTP Yaoundé), gouvernance des données numériques.
\end{exemplebox}

\begin{exoresolu}[Exercice : Analyser le sujet « L'art est-il un luxe ? »]
\textbf{Étape 1 — Termes :} « Art » (création esthétique, pas seulement beaux-arts), « Luxe » (superflu, coûteux, réservé à une élite, non nécessaire à la survie).
\textbf{Étape 2 — Tension :} L'art semble superflu face à la faim (luxe), mais aucune société humaine n'existe sans art (nécessité anthropologique).
\textbf{Étape 3 — Problématique :} « Si l'art ne nourrit pas le corps, pourquoi aucune culture ne peut-elle s'en passer ? L'art est-il un luxe \emph{économique} ou une nécessité \emph{existentielle} ? »
\textbf{Étape 4 — Plan :}
I. Thèse : L'art comme luxe social (distinction bourgeoise, marché de l'art, Bourdieu \emph{La Distinction}).
II. Antithèse : L'art comme besoin vital (rituels, cohésion sociale, catharsis, Aristote \emph{Poétique}, rites Bamiléké).
III. Synthèse : L'art comme \cle{fabrique du sens} — le « luxe » qui rend la vie supportable et la société pensable (Merleau-Ponty, \emph{L'Œil et l'Esprit}).
\end{exoresolu}

\begin{aretenir}[Check-list avant de rendre la copie]
\begin{itemize}
    \item [$\square$] Introduction : 4 étapes respectées ? Problématique en \textbf{question} ? Plan \textbf{annoncé} (I, II, III) ?
    \item [$\square$] Développement : 3 parties équilibrées ? Paragraphes I-A-E ? Transitions \textbf{logiques} (pas « Ensuite », mais « Cependant », « Dès lors », « C'est pourquoi ») ?
    \item [$\square$] Synthèse : Vrai dépassement (pas simple résumé) ? Nouvelle thèse formulée ?
    \item [$\square$] Conclusion : Réponse directe à la problématique ? Ouverture pertinente (autre domaine, actualité, auteur) ?
    \item [$\square$] Citations : Exactes ? Auteur/Œuvre cités ? Intégrées dans la phrase (pas « comme dit Kant : … ») ?
    \item [$\square$] Relecture : Orthographe, accords, majuscules (Kant, Marx, Heidegger), numérotation romaine (I, II, III) ?
\end{itemize}
\end{aretenir}

\coursec{L'introduction : la « porte d'entrée » (4 étapes canoniques)}

Dans la section précédente, nous avons posé les fondements de la dissertation philosophique : une architecture globale en trois moments (introduction, développement, conclusion) répondant à une exigence de rigueur argumentative. Mais une architecture, aussi solide soit-elle, ne vaut que si l'on franchit correctement sa porte d'entrée. Or, en philosophie, cette porte d'entrée a un nom : l'\cle{introduction}. C'est elle que le correcteur lit en premier, c'est elle qui donne le ton de toute la copie, et c'est souvent elle — disons-le sans détour — qui décide de la note. Une introduction ratée condamne même un bon développement ; une introduction réussie dispose favorablement le lecteur. Voyons donc, étape par étape, comment construire une introduction irréprochable.

\courssub{Pourquoi l'introduction est-elle décisive ?}

Imaginez un visiteur qui se présente chez vous. Avant même d'entrer dans le salon, il sonne, se présente, annonce l'objet de sa visite et la manière dont il compte la conduire. Vous ne laisseriez jamais entrer quelqu'un qui déboule sans se nommer ni dire ce qu'il veut. L'introduction philosophique obéit exactement à cette logique de politesse intellectuelle : elle \textbf{présente le sujet}, elle \textbf{précise le sens des mots}, elle \textbf{formule le problème} et elle \textbf{annonce l'itinéraire} de la réflexion.

\begin{definition}[L'introduction philosophique]
L'\cle{introduction} est la première partie de la dissertation. Elle conduit le lecteur du sujet posé vers le problème philosophique qu'il recèle, puis annonce le plan du développement. Elle comprend obligatoirement \textbf{quatre étapes canoniques} : l'accroche, l'analyse des termes du sujet, la problématique et l'annonce du plan. Sa longueur idéale est de \textbf{15 à 20 lignes manuscrites} (environ un quart de page).
\end{definition}

\begin{methode}[La règle des 4 étapes — l'acronyme ADPA]
Pour ne jamais rien oublier, retiens l'acronyme \cle{ADPA} :
\begin{enumerate}
\item \textbf{A}ccrocher : capter l'attention du lecteur par une amorce pertinente ;
\item \textbf{D}éfinir : analyser et définir les termes du sujet (sens philosophique) ;
\item \textbf{P}roblématiser : transformer le sujet en une question philosophique porteuse d'une tension ;
\item \textbf{A}nnoncer : présenter le plan du développement (thèse, antithèse, synthèse).
\end{enumerate}
Ces quatre étapes sont \textbf{séquentielles} : on ne peut pas problématiser avant d'avoir défini, ni annoncer un plan avant d'avoir dégagé le problème.
\end{methode}

\begin{popfigure}
\begin{center}
\begin{tikzpicture}[
node distance=0.55cm,
etape/.style={rectangle, rounded corners=4pt, draw=popblue, thick, fill=popblueL, text width=6.6cm, align=center, font=\small, inner sep=6pt},
fleche/.style={-{Stealth[length=3mm]}, thick, poporange},
num/.style={circle, draw=poporange, fill=poporange, text=white, font=\bfseries\small, inner sep=2pt, minimum size=0.55cm}
]
\node[etape, align=center] (e1){\textbf{Étape 1 — L'accroche}\\ Citation, fait d'actualité, chiffre clé, référence culturelle};
\node[etape, below=of e1, align=center] (e2){\textbf{Étape 2 — Définition des termes}\\ Sens commun / sens technique / sens philosophique};
\node[etape, below=of e2, align=center] (e3){\textbf{Étape 3 — La problématique}\\ Une seule question, porteuse d'une tension};
\node[etape, below=of e3, align=center] (e4){\textbf{Étape 4 — L'annonce du plan}\\ Thèse $\rightarrow$ Antithèse $\rightarrow$ Synthèse};
\node[num, left=0.25cm of e1] {1};
\node[num, left=0.25cm of e2] {2};
\node[num, left=0.25cm of e3] {3};
\node[num, left=0.25cm of e4] {4};
\draw[fleche] (e1) -- (e2);
\draw[fleche] (e2) -- (e3);
\draw[fleche] (e3) -- (e4);
\end{tikzpicture}
\end{center}
\legende{Organigramme des quatre étapes séquentielles de l'introduction (méthode ADPA) : chaque étape prépare la suivante, et aucune ne peut être sautée.}
\end{popfigure}

\courssub{Étape 1 — L'accroche (ou amorce) : capter l'attention}

\courssub{Qu'est-ce qu'une accroche ?}

L'\cle{accroche} (ou amorce) est la première phrase de la copie. Sa fonction est double : \textbf{capter l'attention} du correcteur et \textbf{introduire naturellement le thème} du sujet. Elle doit être courte (une ou deux phrases), pertinente (directement reliée au sujet) et exacte (aucune erreur de fait ni d'attribution).

L'intuition est simple : un correcteur qui lit des centaines de copies est comme un auditeur d'un concours de chant. Les dix premières secondes suffisent à le convaincre — ou à le lasser. Une accroche vive et précise signale immédiatement un candidat cultivé et maîtrisant son sujet.

\begin{propriete}[Les 5 types d'accroches efficaces — boîte à outils]
\begin{enumerate}
\item \textbf{La citation courte} d'un philosophe ou d'un penseur (vérifiée et attribuée) : « La technique n'est plus un simple moyen, elle est devenue un milieu » (idée développée par J. Ellul).
\item \textbf{Le chiffre clé ou la statistique} : « Selon l'INS, la grande majorité des jeunes Camerounais possède aujourd'hui un téléphone portable... »
\item \textbf{Le fait d'actualité} : la généralisation du paiement mobile au Cameroun, la tenue des premières élections régionales en décembre 2020, la construction de grands barrages hydroélectriques.
\item \textbf{La référence culturelle, littéraire ou artistique} : un proverbe, une œuvre, un fait historique (« Un proverbe africain affirme que... »).
\item \textbf{Le paradoxe apparent} : « Jamais l'homme n'a disposé d'autant d'outils pour gagner du temps, et jamais il n'a semblé en manquer autant. »
\end{enumerate}
\end{propriete}

\begin{attention}[Les formules interdites]
Sont \textbf{bannies} car vides, fausses et usées jusqu'à la corde :
\begin{itemize}
\item « Depuis la nuit des temps... »
\item « De tout temps, les hommes... »
\item « L'homme a toujours... »
\end{itemize}
Ces formules prétendent à une connaissance de toute l'histoire humaine que personne ne possède. Elles signalent au correcteur une copie scolaire et paresseuse. Préfère toujours une accroche \textbf{datée, située, vérifiable}.
\end{attention}

\begin{exemplebox}[Accroche camerounaise sur un sujet politique]
Sujet : \emph{« L'État a-t-il pour seul rôle d'assurer la sécurité ? »}

Accroche possible : « Avec la tenue des premières élections régionales en décembre 2020 et la mise en place progressive des conseils régionaux, l'État camerounais s'est vu confier, et a délégué, des missions qui débordent largement la seule sécurité : électrification rurale, éducation, santé. Dès lors, se pose la question du périmètre exact de son rôle. »

L'accroche est ici un \textbf{fait d'actualité nationale}, précis et daté, qui mène naturellement au sujet.
\end{exemplebox}

\courssub{Étape 2 — L'analyse et la définition des termes du sujet}

\courssub{Pourquoi définir ?}

Un sujet de dissertation est une phrase composée de mots. Or les mots sont \textbf{polysémiques} : ils ont plusieurs sens selon les contextes. Si l'on ne fixe pas le sens des termes, on risque de répondre à côté du sujet. Définir, c'est donc \textbf{délimiter le terrain} de la réflexion : c'est tracer les lignes du terrain avant de jouer le match.

\begin{propriete}[Les trois niveaux de sens d'un terme]
Pour chaque terme important du sujet, distingue :
\begin{enumerate}
\item le \textbf{sens commun} : celui du langage quotidien (ce que tout le monde entend) ;
\item le \textbf{sens technique} : celui d'une discipline particulière (droit, économie, informatique) ;
\item le \textbf{sens philosophique} : le sens conceptuel, le plus large et le plus rigoureux, qui sera retenu pour la dissertation.
\end{enumerate}
La définition retenue dans l'introduction est toujours le \textbf{sens philosophique}, éventuellement éclairé par le sens technique.
\end{propriete}

\begin{exemplebox}[Distinguer les sens : deux cas classiques]
\textbf{Cas 1 — « Technique » et « Technologie ».} Au sens commun, on les confond. Or la \emph{technique} désigne l'ensemble des procédés et des moyens que l'homme emploie pour transformer la nature et atteindre ses fins (du grec \emph{tekhnè}, art, savoir-faire) ; la \emph{technologie} est le discours scientifique (\emph{logos}) sur ces techniques, c'est-à-dire leur étude rationnelle. Confondre les deux, c'est confondre l'outil et la science de l'outil.

\textbf{Cas 2 — « Libérer » et « Faciliter ».} \emph{Faciliter} signifie rendre plus aisé (la machine facilite le travail) ; \emph{libérer} signifie rendre libre, affranchir d'une contrainte ou d'une servitude. Dire que la technique « libère » l'homme est une thèse bien plus forte que dire qu'elle « facilite » sa vie : toute la dissertation se joue dans cet écart.
\end{exemplebox}

\begin{attention}[Piège majeur : la définition-dictionnaire]
Ne recopie \textbf{jamais} une définition du Larousse ou du Petit Robert dans une dissertation de philosophie. Le dictionnaire donne le sens \emph{usuel} ; la philosophie exige le sens \emph{conceptuel}. Exemple : le dictionnaire définit le \emph{travail} comme une « activité pénible » ou une « occupation ». Philosophiquement, le travail est une \textbf{activité transformatrice} par laquelle l'homme agit sur la nature, se transforme lui-même et produit des valeurs (Hegel, Marx). Une définition-dictionnaire dans l'introduction fait perdre des points au critère « Analyse du sujet ».
\end{attention}

\courssub{Étape 3 — La problématique : le cœur de l'introduction}

\courssub{De l'observation au problème}

Voici une démarche d'investigation que tu peux reproduire sur n'importe quel sujet. Prenons une observation banale : au marché Mokolo à Yaoundé, une vendeuse utilise le paiement mobile (Orange Money) au lieu de compter les billets. La technique lui \emph{facilite} la vie. Mais observons plus loin : si le réseau tombe en panne, elle ne peut plus vendre ; si son téléphone est volé, elle perd l'accès à son compte. La technique qui la servait devient une \emph{dépendance}. L'observation révèle une \textbf{tension} : le même outil semble à la fois libérer et enchaîner. C'est précisément cette tension qui, formulée en question, devient une problématique philosophique.

\begin{definition}[La problématique]
La \cle{problématique} est la \textbf{question philosophique précise} qui révèle le \textbf{conflit d'interprétations} inhérent au sujet. Elle ne demande pas une réponse de fait (oui/non immédiat), mais met en tension deux réponses également défendables, que le développement devra départager. Elle se formule en \textbf{une seule phrase interrogative}.
\end{definition}

\begin{propriete}[Comment construire une problématique]
\begin{enumerate}
\item Repérer les deux termes (ou concepts) du sujet qui peuvent entrer en tension ;
\item Formuler d'abord la réponse spontanée, évidente (sens commun) ;
\item Montrer que cette réponse évidente peut être contestée (doute philosophique) ;
\item Condenser cette tension en \textbf{une seule question} contenant les termes du sujet.
\end{enumerate}
Exemple réussi : « La technique, en permettant à l'homme de maîtriser la nature, ne risque-t-elle pas de l'asservir à ses propres créations ? » — On y sent les deux pôles : maîtrise (thèse possible) et asservissement (antithèse possible).
\end{propriete}

\begin{attention}[Erreurs fréquentes sur la problématique]
\begin{itemize}
\item \textbf{Multiplier les questions} : trois ou quatre interrogations dispersées ne font pas une problématique, mais un désordre. Une seule question, une seule.
\item \textbf{Reformuler le sujet tel quel} : si le sujet est « La technique libère-t-elle l'homme ? », recopier cette phrase n'apporte rien. Il faut \emph{approfondir} : « La technique, en libérant l'homme des contraintes naturelles, ne crée-t-elle pas de nouvelles dépendances ? »
\item \textbf{Poser une question de fait} : « La technique existe-t-elle ? » n'est pas philosophique, car la réponse est triviale.
\end{itemize}
\end{attention}

\begin{savaistu}[Le critère « Problématisation » au Bac technique]
Dans les grilles de correction harmonisées de la dissertation philosophique, la problématisation est un critère noté à part entière. Une introduction \textbf{sans problématique explicite} plafonne généralement la copie autour de \textbf{6/20}, même si le reste est correct : sans problème posé, le développement ne peut pas être un raisonnement, il n'est qu'une suite d'opinions. La problématique est donc littéralement la clé qui débloque la note.
\end{savaistu}

\courssub{Étape 4 — L'annonce du plan : l'itinéraire de la réflexion}

Dernière étape : indiquer au lecteur la route que suivra le développement. Au Baccalauréat technique, le plan attendu est le \textbf{plan dialectique} en trois parties : thèse, antithèse, synthèse. L'annonce du plan reprend cette structure en deux ou trois phrases, au futur ou avec le pronom « nous ».

\begin{propriete}[Formule canonique de l'annonce du plan]
« Nous examinerons d'abord... \emph{(thèse)}, puis nous verrons que... \emph{(antithèse)}, pour enfin montrer que... \emph{(synthèse)}. »

Exemple appliqué : « Nous examinerons d'abord en quoi la technique apparaît comme un instrument de libération de l'homme ; nous verrons ensuite qu'elle engendre de nouvelles formes de servitude ; nous montrerons enfin que, comprise comme moyen et non comme fin, elle peut concilier maîtrise et liberté. »
\end{propriete}

\begin{attention}[Ce qu'il ne faut pas écrire]
\begin{itemize}
\item Jamais « \textbf{Je vais dire que...} » : la dissertation n'est pas un bavardage personnel ; on écrit « nous examinerons », « il convient d'étudier ».
\item Jamais l'annonce d'un plan « catalogue » : « Nous parlerons des avantages, puis des inconvénients » n'est pas un plan dialectique.
\item Ne pas \textbf{dévoiler la conclusion} : l'annonce du plan donne l'itinéraire, pas la destination finale argumentée.
\end{itemize}
\end{attention}

\courssub{Exemple complet et annoté : une introduction réussie}

\begin{exoresolu}[Introduction type — Sujet : « La technique libère-t-elle l'homme de la nature ? »]
Rédige l'introduction complète de ce sujet en respectant les quatre étapes ADPA.
\tcblower
\textbf{Introduction rédigée :}

« En 2022, le Cameroun poursuivait ses programmes de modernisation agricole fondés sur l'irrigation et la mécanisation, promettant de libérer le paysan de l'aléa des saisons. \emph{(Accroche : fait d'actualité national.)} La \emph{technique}, au sens philosophique, désigne l'ensemble des procédés par lesquels l'homme transforme la nature pour atteindre ses fins ; la \emph{nature} renvoie à l'ordre du monde indépendant de la volonté humaine ; enfin, \emph{libérer} ne signifie pas simplement faciliter, mais affranchir d'une contrainte. \emph{(Définitions des trois termes, au sens philosophique.)} Or, si la technique semble soustraire l'homme aux nécessités naturelles, elle le rend parallèlement dépendant de ses propres artefacts : la technique, en maîtrisant la nature, ne risque-t-elle pas d'asservir l'homme à ses propres créations ? \emph{(Problématique : une seule question, porteuse de la tension maîtrise/asservissement.)} Nous examinerons d'abord la technique comme instrument d'émancipation vis-à-vis de la nature ; nous verrons ensuite qu'elle engendre de nouvelles dépendances ; nous montrerons enfin que l'usage raisonné de la technique permet de concilier maîtrise de la nature et liberté humaine. \emph{(Annonce du plan dialectique.)} »

\textbf{Bilan} : 4 étapes présentes, 13 à 18 lignes manuscrites, une seule question de problématique, aucune formule interdite. Cette introduction remplit tous les critères de la grille.
\end{exoresolu}

\begin{popfigure}
\begin{center}
\begin{tikzpicture}[font=\small]
\node[rectangle, draw=popdark, thick, fill=popcream, text width=10.2cm, align=justify, inner sep=8pt] (intro) {%
« En 2022, le Cameroun poursuivait ses programmes de modernisation agricole fondés sur l'irrigation et la mécanisation. La \emph{technique} désigne l'ensemble des procédés par lesquels l'homme transforme la nature ; \emph{libérer} signifie affranchir d'une contrainte, et non simplement faciliter. Or, si la technique soustrait l'homme aux nécessités naturelles, elle le rend dépendant de ses artefacts : la technique, en maîtrisant la nature, ne risque-t-elle pas d'asservir l'homme à ses propres créations ? Nous examinerons d'abord la technique comme émancipation, puis les nouvelles dépendances qu'elle crée, pour montrer enfin qu'un usage raisonné concilie maîtrise et liberté. »};
\node[rectangle, rounded corners, draw=popblue, fill=popblueL, font=\footnotesize\bfseries, right=0.5cm of intro.north east, anchor=north west] (a1) {1. Accroche};
\node[rectangle, rounded corners, draw=popgreen, fill=popgreenL, font=\footnotesize\bfseries, below=0.35cm of a1] (a2) {2. Définitions};
\node[rectangle, rounded corners, draw=poporange, fill=poporangeL, font=\footnotesize\bfseries, below=0.35cm of a2] (a3) {3. Problématique};
\node[rectangle, rounded corners, draw=poppurple, fill=poppurpleL, font=\footnotesize\bfseries, below=0.35cm of a3] (a4) {4. Plan};
\draw[-{Stealth}, popblue, thick] (a1.west) -- ++(-0.3,0) |- ([yshift=-0.15cm]intro.north east);
\draw[-{Stealth}, popgreen, thick] (a2.west) -- (intro.east);
\draw[-{Stealth}, poporange, thick] (a3.west) -- ([yshift=-0.4cm]intro.east);
\draw[-{Stealth}, poppurple, thick] (a4.west) -- ([yshift=-1.1cm]intro.south east);
\end{tikzpicture}
\end{center}
\legende{Introduction type annotée sur le sujet « Technique et Nature » : les flèches repèrent les quatre étapes canoniques (1 : accroche factuelle ; 2 : définitions philosophiques ; 3 : problématique en une question ; 4 : annonce du plan dialectique).}
\end{popfigure}

\begin{aretenir}[L'essentiel de la section]
\begin{itemize}
\item L'introduction compte \textbf{4 étapes obligatoires} (ADPA) : Accroche, Définitions, Problématique, Annonce du plan.
\item L'accroche est \textbf{datée, située, vérifiable} ; jamais « depuis la nuit des temps ».
\item On définit les termes au \textbf{sens philosophique}, jamais avec le dictionnaire courant.
\item La problématique est \textbf{une seule question} révélant une tension ; sans elle, la copie plafonne.
\item L'annonce du plan suit la dialectique \textbf{thèse / antithèse / synthèse}, sans « je vais dire que ».
\item Longueur idéale : \textbf{15 à 20 lignes manuscrites}.
\end{itemize}
\end{aretenir}

\begin{exercice}[À toi de jouer]
Sujet : \emph{« L'État a-t-il pour seul rôle d'assurer la sécurité ? »} Rédige une introduction complète (15 à 20 lignes) en respectant les quatre étapes ADPA. Tu utiliseras comme accroche un fait d'actualité camerounaise (décentralisation, électrification rurale, construction d'infrastructures), tu définiras « État », « rôle » et « sécurité » au sens philosophique, tu formuleras une problématique en une seule question, puis tu annonceras un plan dialectique.
\end{exercice}

\begin{corrige}[Exercice — éléments attendus]
\textbf{Accroche} : fait national précis, ex. « Avec la tenue des premières élections régionales en décembre 2020 et l'installation des conseils régionaux, l'État camerounais a vu ses missions s'étendre au développement local, bien au-delà du maintien de l'ordre. »
\textbf{Définitions} : l'\emph{État} est l'organisation politique souveraine exerçant un pouvoir légitime sur un territoire et une population ; le \emph{rôle} désigne la fonction ou la fin assignée à une institution ; la \emph{sécurité} est la protection des personnes et des biens contre les menaces (sens politique : ordre public).
\textbf{Problématique} (une seule question) : « Si l'État se réduisait à garantir la sécurité, ne trahirait-il pas sa vocation de justice et de bien commun ? » ou « La sécurité est-elle la fin unique de l'État, ou seulement la condition d'autres missions ? »
\textbf{Annonce du plan} : « Nous examinerons d'abord la thèse de l'État-gendarme, pour qui la sécurité est la fonction première ; nous verrons ensuite que cette conception minimaliste est insuffisante ; nous montrerons enfin que la sécurité est la condition nécessaire, mais non la fin unique, de l'action étatique. »
\end{corrige}

\coursec{Le développement (1) : Thèse et Antithèse — L'art du paragraphe argumentatif (I-A-E)}

Dans la section précédente, nous avons construit l'introduction, cette « porte d'entrée » qui se referme sur une problématique claire et un plan annoncé. Mais une porte, si belle soit-elle, ne fait pas une maison. C'est dans le \cle{développement} que se joue l'essentiel de la note : c'est là que le correcteur vérifie si vous savez réellement \emph{penser}, c'est-à-dire défendre une position, puis la contester, avant de la dépasser. Le développement d'une dissertation dialectique au Baccalauréat technique comprend deux grands moments : la \cle{Thèse} et l'\cle{Antithèse} (la Synthèse fera l'objet de la section suivante). Chacun de ces moments est constitué de paragraphes, et chaque paragraphe obéit à une architecture unique et non négociable : la règle \cle{I-A-E}.

\courssub{La structure unique du paragraphe : la règle I-A-E}

Commençons par une intuition simple. Imaginez que vous vouliez convaincre un ami que « la technique libère l'homme ». Vous ne pouvez pas vous contenter de l'affirmer : il vous répondra « pourquoi ? ». Vous donnerez alors une raison (« parce qu'elle nous débarrasse des contraintes naturelles »). Mais il pourra encore douter : « ah bon ? montre-moi ! ». Vous devrez alors montrer un cas précis (« regarde le barrage de Lom Pangar : l'eau qui menaçait les récoltes produit désormais l'électricité de toute une région »). Vous venez, sans le savoir, de construire un paragraphe philosophique : une \textbf{Idée}, un \textbf{Argument}, un \textbf{Exemple}.

\begin{definition}[Le paragraphe argumentatif]
Un paragraphe de dissertation philosophique est une unité de pensée qui développe \textbf{une seule idée} selon trois temps indissociables :
\begin{enumerate}
\item \textbf{I — Idée directrice} : la phrase-thèse du paragraphe, qui annonce clairement ce que l'on va démontrer ;
\item \textbf{A — Argument(s)} : le ou les raisonnements philosophiques qui justifient l'idée (appuyés, si possible, sur un auteur du programme) ;
\item \textbf{E — Exemple(s) analysé(s)} : un cas concret, précis et \emph{explicitement relié} à l'argument.
\end{enumerate}
\end{definition}

\begin{popfigure}
\begin{center}
\begin{tikzpicture}[
  node distance=0.55cm,
  box/.style={draw=popblue, thick, rounded corners, fill=popblueL, text width=0.27\linewidth, align=center, minimum height=1.7cm, font=\small},
  arr/.style={-{Stealth[length=3mm]}, very thick, poporange},
  lab/.style={font=\small\bfseries, text=popdark}
]
\node[box, align=center] (idee){\textbf{I — Idée directrice}\\ Phrase-thèse : ce que je vais prouver};
\node[box, right=of idee, align=center] (arg){\textbf{A — Argument}\\ Raisonnement philosophique (auteur, concept)};
\node[box, right=of arg, align=center] (ex){\textbf{E — Exemple analysé}\\ Cas concret relié à l'argument};
\draw[arr] (idee) -- (arg);
\draw[arr] (arg) -- (ex);
\node[lab, above=0.15cm of idee] {1};
\node[lab, above=0.15cm of arg] {2};
\node[lab, above=0.15cm of ex] {3};
\end{tikzpicture}
\end{center}
\legende{Schéma 1 — La chaîne I-A-E : un paragraphe = une idée, justifiée par un argument, incarnée dans un exemple analysé.}
\end{popfigure}

\begin{methode}[Appliquer la règle I-A-E à chaque paragraphe]
\begin{enumerate}
\item Rédiger d'abord la phrase-thèse du paragraphe, au brouillon, en une ligne. Si vous n'y arrivez pas, c'est que le paragraphe n'a pas d'idée : supprimez-le.
\item Développer l'argument en 3 à 6 lignes : pourquoi cette idée est-elle vraie ? Quel concept, quel auteur la soutient ?
\item Choisir un exemple précis (lieu, date, mécanisme) et l'\emph{analyser} : montrer en quoi il illustre exactement l'argument.
\item Vérification finale : \textbf{mon exemple prouve-t-il mon argument ?} Si l'exemple illustre une autre idée, ou s'il n'illustre rien, le remplacer.
\end{enumerate}
\end{methode}

\begin{attention}[Le point de contrôle décisif]
L'erreur la plus fréquente est le « collage » d'exemple : poser un exemple à la fin du paragraphe sans le relier à l'argument. Un exemple non analysé ne vaut presque rien. Il faut toujours ajouter une phrase d'analyse du type : « Ce cas montre bien que... », « Cet exemple illustre le fait que... ».
\end{attention}

\courssub{La Thèse : la réponse spontanée et évidente}

\begin{definition}[La Thèse]
La \cle{Thèse} est la première réponse à la problématique : la réponse la plus spontanée, la plus évidente, celle du sens commun ou d'une tradition philosophique bien établie. Elle constitue la première partie du développement (2 à 3 paragraphes).
\end{definition}

Pourquoi commencer par l'évidence ? Parce que la démarche dialectique reproduit le mouvement naturel de la pensée : on épouse d'abord l'opinion commune, on l'éprouve, puis on la dépasse. Sur le sujet « La technique libère-t-elle l'homme ? », la thèse sera : \emph{« Oui, la technique libère l'homme par la maîtrise des contraintes naturelles. »} C'est l'opinion de tout un chacun, mais aussi celle de philosophes comme Descartes, qui voulait nous rendre « comme maîtres et possesseurs de la nature ».

\begin{exemplebox}[Un paragraphe de Thèse construit selon I-A-E]
\textbf{(I)} En premier lieu, la technique libère l'homme en le soustrayant aux contraintes de la nature. \textbf{(A)} \emph{En effet}, l'homme est un être naturellement démuni : sans outils, il subit la faim, le froid, les maladies et les crues. La technique, comprise comme l'ensemble des procédés par lesquels l'homme transforme son milieu, compense cette faiblesse originelle et le rend maître de son destin. \textbf{(E)} \emph{Par exemple}, le barrage de Lom Pangar, dans la région de l'Est du Cameroun, retient les eaux de la Sanaga : ce qui était une menace d'inondation et d'irrégularité devient une source d'électricité et de sécurité pour des millions de Camerounais. Ce cas montre bien que la technique, en maîtrisant une contrainte naturelle, élargit la liberté humaine.
\end{exemplebox}

Observez les \cle{connecteurs logiques} : « En premier lieu » introduit l'idée, « En effet » introduit l'argument, « Par exemple » introduit l'exemple, et la dernière phrase effectue l'analyse. Ces charnières sont obligatoires : elles rendent la démonstration lisible.

\courssub{L'Antithèse : l'objection radicale et le renversement de perspective}

Une pensée qui ne rencontre aucune résistance n'est pas une pensée philosophique : c'est une opinion endormie. L'\cle{Antithèse} est le moment où l'on conteste ce que l'on vient d'affirmer.

\begin{definition}[L'Antithèse]
L'\cle{Antithèse} est l'objection radicale opposée à la thèse : une critique interne (la thèse se contredit elle-même) ou externe (un autre point de vue la réfute), qui opère un \emph{renversement de perspective}. Elle constitue la deuxième partie du développement (2 à 3 paragraphes).
\end{definition}

Sur notre sujet, l'antithèse sera : \emph{« Non, la technique aliène l'homme : elle crée la dépendance, renforce la bureaucratie et produit l'aliénation dénoncée par Marx, pour qui l'ouvrier est dépossédé du produit de son travail. »} Le renversement est total : ce qui apparaissait comme instrument de liberté devient instrument de servitude.

\begin{exemplebox}[Un paragraphe d'Antithèse construit selon I-A-E]
\textbf{(I)} \emph{Cependant}, la technique, loin de libérer, peut enfermer l'homme dans de nouvelles dépendances. \textbf{(A)} \emph{En effet}, chaque outil exige celui qui le possède, l'entretient et le répare : l'usager devient tributaire d'objets qu'il ne maîtrise plus. Marx a montré que le travailleur peut être aliéné par la machine au lieu d'être libéré par elle. \emph{En outre}, la technique produit ses propres nuisances. \textbf{(E)} \emph{Par exemple}, les embouteillages quotidiens de Yaoundé et de Douala transforment l'automobile — symbole de liberté de mouvement — en prison mobile où des milliers de citoyens perdent chaque jour des heures de vie. De même, les déchets électroniques déversés à Agbogbloshie, au Ghana, montrent que le progrès technique des uns devient l'empoisonnement des autres. Ces exemples illustrent que la maîtrise technique peut se retourner contre l'homme.
\end{exemplebox}

\begin{propriete}[La règle du même plan conceptuel]
La Thèse et l'Antithèse doivent s'opposer sur le \textbf{même plan conceptuel}. Si la thèse affirme que la technique \emph{libère}, l'antithèse doit affirmer qu'elle \emph{aliène} : l'opposition porte sur la liberté. Il est fautif d'opposer « la technique libère » (plan de la liberté) à « la technique rend heureux » (plan du bonheur) : ce ne serait plus une contradiction, mais deux sujets différents. Le correcteur sanctionne immédiatement ce glissement.
\end{propriete}

\begin{popfigure}
\begin{center}
\begin{tikzpicture}[
  grow=right,
  level distance=3.6cm,
  level 1/.style={sibling distance=3.4cm},
  level 2/.style={sibling distance=1.6cm, level distance=3.9cm},
  root/.style={draw=popdark, thick, rounded corners, fill=popcream, align=center, font=\small\bfseries, text width=2.6cm},
  branche/.style={draw=popblue, thick, rounded corners, fill=popblueL, align=center, font=\small, text width=3.1cm},
  feuille/.style={draw=poporange, rounded corners, fill=poporangeL, align=center, font=\scriptsize, text width=3.4cm},
  edge from parent/.style={draw=popmuted, thick}
]
\node[root] {La technique libère-t-elle l'homme ?}
  child { node[branche, align=center]{ANTITHÈSE\\ La technique aliène}
    child { node[feuille, align=center]{Obj. 2 : nuisances\\ \emph{Ex.} : déchets électroniques d'Agbogbloshie} }
    child { node[feuille, align=center]{Obj. 1 : dépendance\\ \emph{Ex.} : embouteillages Yaoundé/Douala} }
  }
  child { node[branche, align=center]{THÈSE\\ La technique libère}
    child { node[feuille, align=center]{Arg. 2 : lien social\\ \emph{Ex.} : mobile money, inclusion bancaire} }
    child { node[feuille, align=center]{Arg. 1 : maîtrise de la nature\\ \emph{Ex.} : barrage de Lom Pangar} }
  };
\end{tikzpicture}
\end{center}
\legende{Schéma 2 — Arbre argumentatif : la thèse et l'antithèse se déploient chacune en arguments étayés d'exemples camerounais et africains, sur le même plan conceptuel (liberté / aliénation).}
\end{popfigure}

\courssub{Les règles d'or du développement : quantité, connecteurs et exemples}

\begin{aretenir}[Les règles d'or]
\begin{itemize}
\item \textbf{1 idée = 1 paragraphe.} Jamais deux idées dans le même paragraphe, jamais une idée sur deux paragraphes.
\item \textbf{2 à 3 paragraphes pour la Thèse, 2 à 3 pour l'Antithèse}, soit 4 à 6 paragraphes de développement.
\item \textbf{Connecteurs logiques obligatoires} : \emph{En effet}, \emph{Par exemple}, \emph{Certes}, \emph{Cependant}, \emph{En outre}, \emph{Par conséquent}. Ils signalent au correcteur la progression de votre raisonnement.
\item \textbf{Chaque exemple est analysé} : il doit être explicitement relié à l'argument qu'il illustre.
\end{itemize}
\end{aretenir}

\begin{attention}[L'erreur fatale : l'exemple « M. Tout-le-monde »]
Sont interdits : l'anecdote personnelle non conceptualisée (« Mon oncle a un tracteur, donc la technique libère »), l'exemple vague (« les gens, de nos jours... ») et l'exemple « M. Tout-le-monde » sans lieu ni précision. Un bon exemple est \emph{daté, situé et analysé}. Votre vie privée n'est pas une preuve philosophique ; elle ne peut servir que si elle est élevée au rang de cas général.
\end{attention}

\begin{exemplebox}[Le poids chiffré de l'exemple]
Au Baccalauréat technique, l'argumentation est notée sur 6 points. Un paragraphe sans exemple analysé fait perdre \textbf{1,5 à 2 points} sur ce critère. Sur 4 paragraphes sans exemples, c'est l'ensemble du critère Argumentation qui s'effondre — et avec lui, souvent, la moyenne.
\end{exemplebox}

L'argument peut aussi s'appuyer sur une citation, à condition de l'\emph{analyser} et non de la « poser ». Ainsi, André Leroi-Gourhan définit la technique comme « l'ensemble des procédés » par lesquels l'homme agit sur son milieu : cette définition, une fois expliquée, fonde rigoureusement l'argument de la maîtrise de la nature. Une citation jetée sans commentaire ne rapporte rien ; une citation analysée vaut un argument.

\begin{savaistu}[Les exemples qui font la différence]
Les correcteurs valorisent particulièrement les exemples issus du programme d'Enseignement Moral, Civique et de Réarmement (EMCRE) et de l'Histoire du Cameroun : la construction du barrage de Lom Pangar, l'inclusion bancaire par le téléphone mobile (mobile money) dans les quartiers de Douala ou les villages du Nord, la lutte contre les déchets électroniques en Afrique de l'Ouest. Ces exemples montrent que vous savez relier la philosophie à votre environnement national.
\end{savaistu}

\begin{exoresolu}[Construire un paragraphe d'antithèse selon I-A-E]
Énoncé. Sujet : « La technique libère-t-elle l'homme ? » Rédiger un paragraphe d'antithèse dont l'idée directrice est : « La technique crée une dépendance qui aliène l'usager. »
\tcblower
\textbf{Solution détaillée.}
\textbf{(I)} \emph{En outre}, la technique enchaîne l'homme à ses propres créations en le rendant dépendant de ses outils. \textbf{(A)} \emph{En effet}, celui qui adopte un outil en adopte aussi les contraintes : il doit l'entretenir, le recharger, le réparer, et il se trouve démuni dès que l'outil fait défaut. La liberté apparente se paie ainsi d'une servitude invisible : plus l'outil est puissant, plus son absence est paralysante. \textbf{(E)} \emph{Par exemple}, le téléphone mobile a certes permis l'inclusion bancaire de millions de Camerounais grâce au mobile money ; mais une simple panne de réseau ou une coupure d'électricité suffit à bloquer les paiements, les transferts et le commerce de tout un quartier de Douala. Ce cas montre bien que l'instrument de libération devient, par la dépendance qu'il instaure, une nouvelle forme d'aliénation. \emph{Par conséquent}, la technique ne libère pas sans condition : elle exige en retour une soumission à ses propres exigences.
\end{exoresolu}

\begin{exercice}[À vous de construire]
Sujet : « L'État doit-il tout contrôler ? »
\begin{enumerate}
\item Rédiger la phrase-thèse (idée directrice) d'un paragraphe de thèse et d'un paragraphe d'antithèse, en veillant à rester sur le même plan conceptuel.
\item Rédiger entièrement le paragraphe de thèse selon la règle I-A-E, avec un exemple camerounais analysé.
\item Souligner dans votre texte les connecteurs logiques utilisés.
\end{enumerate}
\end{exercice}

\begin{corrige}[Exercice 1]
\textbf{1.} Thèse : « L'État doit garantir l'ordre et la sécurité, conditions de la liberté de tous » (plan de la liberté). Antithèse : « Mais un État qui contrôle tout étouffe les libertés individuelles qu'il prétend protéger » (même plan : la liberté). L'opposition est correcte car les deux portent sur la liberté.
\textbf{2.} Exemple attendu (modèle) : \textbf{(I)} L'État est d'abord nécessaire comme garant de l'ordre, sans lequel aucune liberté n'est possible. \textbf{(A)} \emph{En effet}, selon Hobbes, sans pouvoir commun qui fasse respecter les lois, la vie serait « solitaire, pauvre, dangereuse, brutale et brève » : la liberté de chacun suppose donc l'autorité de tous. \textbf{(E)} \emph{Par exemple}, l'organisation par l'État camerounais des examens officiels comme le Baccalauréat technique garantit à chaque candidat, de Garoua à Buea, les mêmes règles et les mêmes chances ; sans ce cadre, ce serait la loi du plus fort ou du plus riche. Ce cas illustre que le contrôle étatique, loin d'opprimer, rend la liberté effective.
\textbf{3.} Connecteurs à souligner : \emph{En effet}, \emph{Par exemple}, ainsi que la phrase d'analyse finale.
\end{corrige}

Ainsi, thèse et antithèse, construites en paragraphes rigoureusement structurés selon la règle I-A-E, dressent le conflit des idées dans toute sa force. Reste à dépasser ce conflit : ce sera l'objet de la synthèse, que nous étudierons dans la section suivante.

\coursec{Le développement (2) : La Synthèse — Le dépassement dialectique}

\courssub{De l'affrontement à la résolution : pourquoi la synthèse est le moment de vérité}

Après avoir construit une \textbf{Thèse} solide (section 3) et une \textbf{Antithèse} rigoureuse qui en a mis en lumière les limites, l'esprit du candidat se trouve face à une contradiction apparente. La dissertation philosophique ne saurait s'arrêter sur un match nul (\og{} les deux ont raison \fg{}) ni sur un éclectisme mou (\og{} un peu l'un, un peu l'autre \fg{}). Elle exige un \textbf{troisième moment} : la \cle{Synthèse}.

\emph{Transition depuis la section 3 :} Vous avez maîtrisé la structure I-A-E (Idée - Argument - Exemple) pour vos paragraphes de Thèse et d'Antithèse. La Synthèse ne suit \textbf{pas} ce schéma. Elle ne rajoute pas un argument ; elle \textbf{change de niveau d'analyse}. C'est le passage de la \emph{multiplicité} des arguments à l'\emph{unité} de la compréhension.

\begin{definition}[La Synthèse (Aufheben hégélien)]
Dans la dialectique hégélienne, \emph{Aufheben} porte trois sens indissociables que la synthèse philosophique doit réaliser simultanément :
\begin{enumerate}
    \item \textbf{Supprimer (Annuler) :} La contradiction brute entre Thèse et Antithèse est levée ; elles ne s'affrontent plus comme des opinions rivales.
    \item \textbf{Conserver (Maintenir) :} La part de vérité contenue dans chacune n'est pas jetée, mais intégrée comme un \emph{moment} nécessaire de la vérité totale.
    \item \textbf{Élever (Sublimer) :} La pensée accède à un concept plus riche, plus concret, qui explique \emph{pourquoi} les deux moments partiels étaient vrais mais insuffisants pris séparément.
\end{enumerate}
La Synthèse n'est pas la \og{} conclusion \fg{} au sens scolaire (résumé), c'est la \textbf{réponse philosophique définitive} au sujet.
\end{definition}

\courssub{Les trois formes canoniques de la Synthèse}

Toute synthèse de qualité au Baccalauréat technique s'appuie sur l'une de ces trois figures logiques (souvent combinées). Les identifier vous permet de \textbf{construire} votre synthèse au lieu de la \og{} chercher \fg{} au hasard.

\begin{center}
\begin{popfigure}
\begin{tikzpicture}[
    node distance=1.2cm and 1.5cm,
    box/.style={draw=popdark, thick, rounded corners=6pt, fill=popcream, text width=4.2cm, align=center, font=\small, inner sep=6pt},
    title/.style={font=\bfseries\small, text=popblue},
    arrow/.style={-Stealth, thick, popblue},
    brace/.style={decorate, decoration={brace, amplitude=6pt}, thick, poporange}
]
% Nodes
\node[box, title, align=center] (dist){\textbf{1. DISTINCTION}\\[4pt] (Séparation conceptuelle)};
\node[box, title, right=of dist, align=center] (recon){\textbf{2. RÉCONCILIATION}\\[4pt] (Condition d'unité)};
\node[box, title, right=of recon, align=center] (depass){\textbf{3. DÉPASSEMENT}\\[4pt] (Niveau ontologique supérieur)};

% Content Distinction
\node[box, below=0.2cm of dist, fill=popblueL!30, text width=4.2cm, font=\small] (distc) {
\textbf{Connecteurs :} \emph{Pourvu que, à condition que, dès lors que, si et seulement si.}\\
\textbf{Mécanisme :} Les deux thèses parlent de \textbf{choses différentes} masquées par un même mot.\\
\textbf{Exemple :} \og{} Technique \fg{} = Outil (moyen neutre) \textbf{vs} Système technique (Ellul : milieu autonome).
};

% Content Réconciliation
\node[box, below=0.2cm of recon, fill=popgreenL!30, text width=4.2cm, font=\small] (reconc) {
\textbf{Connecteurs :} \emph{Dans la mesure où, relativement à, sous réserve de, en tant que.}\\
\textbf{Mécanisme :} Les deux thèses sont vraies \textbf{sous des conditions distinctes} qu'il faut articuler.\\
\textbf{Exemple :} La technique libère \textbf{si} l'homme en reste le maître politique (éthique/droit).
};

% Content Dépassement
\node[box, below=0.2cm of depass, fill=poporangeL!30, text width=4.2cm, font=\small] (depassc) {
\textbf{Connecteurs :} \emph{En définitive, radicalement, structurellement, c'est pourquoi.}\\
\textbf{Mécanisme :} La contradiction révèle la \textbf{condition d'être} du sujet (nature humaine, essence de la technique).\\
\textbf{Exemple :} \og{} L'homme est un animal technique \fg{} (Bergson/Platon/Simondon) : la technique \emph{constitue} l'humain.
};

% Arrows
\draw[arrow] (dist.south) -- ++(0,-0.3) -| (distc.north);
\draw[arrow] (recon.south) -- ++(0,-0.3) -| (reconc.north);
\draw[arrow] (depass.south) -- ++(0,-0.3) -| (depassc.north);

% Braces
\draw[brace] ($(distc.south west)!0.5!(reconc.south east)$) -- node[below=4pt, font=\bfseries\small, text=poporange] {Choix selon la nature du sujet} ($(distc.south east)!0.5!(reconc.south west)$);
\draw[brace] ($(reconc.south west)!0.5!(depassc.south east)$) -- node[below=4pt, font=\bfseries\small, text=poporange] {Profondeur croissante} ($(reconc.south east)!0.5!(depassc.south west)$);

\end{tikzpicture}
\legende{Tableau comparatif des trois formes canoniques de la Synthèse : connecteurs logiques, mécanisme intellectuel et exemples types pour le Bac Tech.}
\end{popfigure}
\end{center}

\paragraph{1. La Synthèse par Distinction (La plus fréquente au Bac).}
La contradiction naît d'une \textbf{équivocité} : un même terme cache deux concepts. La thèse a raison \emph{au sens A}, l'antithèse a raison \emph{au sens B}.
\begin{itemize}
    \item \textbf{Structure mentale :} \og{} La Thèse traite de X\emph{1}, l'Antithèse de X\emph{2}. X\emph{1} $\neq$ X\emph{2}. \fg{}
    \item \textbf{Sujets types :} \og{} La technique est-elle neutre ? \fg{} (Outil vs Système), \og{} Le travail libère-t-il ? \fg{} (Activité anthropologique vs Aliénation salariée), \og{} L'État est-il au service de la nation ? \fg{} (Idéal républicain vs Appareil bureaucratique).
\end{itemize}

\paragraph{2. La Synthèse par Réconciliation (Conditionnelle).}
Les deux moments sont vrais mais \textbf{dépendent d'une condition externe} (éthique, politique, juridique, anthropologique) qui n'était pas explicitée.
\begin{itemize}
    \item \textbf{Structure mentale :} \og{} La Thèse est vraie \emph{si condition C1}, l'Antithèse est vraie \emph{si condition C2}. La vérité totale exige l'articulation de C1 et C2. \fg{}
    \item \textbf{Sujets types :} \og{} La science progresse-t-elle ? \fg{} (Vérité épistémologique vs Responsabilité éthique), \og{} La démocratie garantit-elle la liberté ? \fg{} (Procédure formelle vs Effectivité sociale).
\end{itemize}

\paragraph{3. La Synthèse par Dépassement (Ontologique / Anthropologique).}
La contradiction n'est pas une erreur de définition, elle \textbf{structure l'être même} du sujet. On ne \og{} choisit \fg{} pas, on \og{} comprend \fg{} la nécessité de l'opposition.
\begin{itemize}
    \item \textbf{Structure mentale :} \og{} Thèse et Antithèse sont les deux faces d'une même structure fondamentale (ex: la condition humaine technique). \fg{}
    \item \textbf{Auteurs clés :} \textbf{Gilbert Simondon} (l'individuation technique : l'objet technique n'est pas aliénant \emph{par nature}, il faut passer de l'objet \og{} aliéné \fg{} à l'objet \og{} concret \fg{} par la culture technique), \textbf{Jacques Ellul} (le \og{} système technicien \fg{} comme milieu total qui détermine les fins), \textbf{Paulin Hountondji} (la technique comme levier de \og{} développement endogène \fg{} vs importation aliénante).
    \item \textbf{Sujets types :} \og{} L'homme est-il un animal technique ? \fg{}, \og{} La technique détermine-t-elle l'histoire ? \fg{}.
\end{itemize}

\begin{methode}[3 questions pour trouver VOTRE synthèse (Brouillon obligatoire)]
Ne commencez pas à rédiger la synthèse avant d'avoir répondu par écrit à ces trois questions sur votre brouillon :
\begin{enumerate}
    \item \textbf{Sur quel point précis (concept, présupposé, définition) la Thèse et l'Antithèse s'opposent-elles vraiment ?} (Formulez l'antinomie : \og{} Pour la Thèse, X = A ; pour l'Antithèse, X = non-A \fg{}).
    \item \textbf{Quel concept opératoire (distinction, condition, structure) permet de dissoudre cette opposition sans nier l'une ni l'autre ?} (Cherchez le \emph{tertium datur} : le tiers exclu qui les unit).
    \item \textbf{Quelle thèse unique, plus forte, plus précise, émerge de cette résolution ?} (Écrivez-la en une phrase affirmative, sans \og{} peut-être \fg{}, sans \og{} il faut \fg{}).
\end{enumerate}
\emph{Astuce :} Si vous ne trouvez pas de distinction conceptuelle nette (Q1), tentez la réconciliation conditionnelle (Q2). Si le sujet porte sur l'essence de l'homme/histoire/technique, visez le dépassement (Q3).
\end{methode}

\courssub{Architecture rédactionnelle : 1 ou 2 paragraphes, zéro argument nouveau}

\begin{propriete}[Structure standard de la Synthèse (1 à 2 paragraphes max)]
La Synthèse ne contient \textbf{aucun nouvel exemple}, \textbf{aucune nouvelle citation}, \textbf{aucun nouvel auteur}. Elle ne fait que \textbf{relire} les arguments déjà produits sous un jour nouveau.
\begin{enumerate}
    \item \textbf{Phrase d'amorce (Obligatoire) :} Marque le basculement dialectique.
        \begin{itemize}
            \item \og{} Au terme de cette analyse, il apparaît que... \fg{} (Classique, solide).
            \item \og{} La contradiction se résout si l'on distingue... \fg{} (Distinction).
            \item \og{} L'opposition trouve son unité dans la condition suivante : ... \fg{} (Réconciliation).
            \item \og{} C'est que, radicalement, ... \fg{} (Dépassement).
        \end{itemize}
    \item \textbf{Explicitation du concept résolveur :} Définissez la distinction / la condition / la structure découverte.
    \item \textbf{Relecture unificatrice :} Montrez comment la Thèse et l'Antithèse s'inscrivent comme \emph{moments} de cette vérité supérieure (\og{} La thèse avait raison en ce sens que... ; l'antithèse complétait en montrant que... \fg{}).
    \item \textbf{Formulation finale de la réponse :} Une phrase affirmative, nuancée, qui \textbf{est} la réponse au sujet.
\end{enumerate}
\end{propriete}

\begin{attention}[Pièges mortels (Note maximale 11/20 si présence)]
\begin{itemize}
    \item \textbf{\og{} Synthèse = Résumé \fg{} :} \og{} En résumé, la thèse dit A, l'antithèse dit B. \fg{} $\rightarrow$ \textbf{Zéro point} argumentation. La synthèse \textbf{produit} du sens, elle ne le \textbf{répète} pas.
    \item \textbf{\og{} Synthèse = Nouvelle Thèse \fg{} :} Introduire un argument d'autorité ou un exemple inédit au dernier moment. $\rightarrow$ Hors structure, perte de cohérence.
    \item \textbf{\og{} Synthèse = Compromis mou \fg{} :} \og{} Il y a du vrai dans les deux, tout dépend des cas. \fg{} $\rightarrow$ Relativisme, absence de pensée conceptuelle. Remplacez \og{} tout dépend \fg{} par \og{} cela dépend \emph{précisément de la distinction entre X et Y} \fg{}.
    \item \textbf{Marqueurs scolaires interdits :} \og{} En synthèse, je dirais que... \fg{}, \og{} Pour synthétiser... \fg{}, \og{} Thèse / Antithèse / Synthèse : \fg{} (Listes). $\rightarrow$ Style non philosophique, sanctionné.
    \item \textbf{Synthèse trop longue :} Plus de 2 paragraphes = vous êtes en train de faire un 3ème développement. La synthèse est \textbf{dense et brève} (environ 15-20 lignes max).
\end{itemize}
\end{attention}

\courssub{Exemple résolu : Sujet type Bac Tech — \og{} La technique libère-t-elle ? \fg{}}

\begin{exoresolu}[Simulation Bac : Construction de la Synthèse]
\textbf{Énoncé :} \og{} La technique libère-t-elle ? \fg{} (Sujet classique, programme : Technique / Liberté).

\textbf{Rappel rapide Thèse / Antithèse (Section 3) :}
\begin{itemize}
    \item \textbf{Thèse :} Oui. Technique = puissance d'agir (maîtrise nature), gain de temps (loisirs), prothèses (handicap), communication. Auteurs : Aristote (esclaves/outils), Bergson (homo faber), \textbf{Simondon} (objet technique libérateur si \og{} concret \fg{}).
    \item \textbf{Antithèse :} Non. Technique = aliénation (Marx : machine capitaliste), asservissement (Ellul : \og{} système technicien \fg{}, efficacité contrainte), dépendance (panne, énergie), destruction milieu. Auteur : \textbf{Ellul}, Jonas (responsabilité).
\end{itemize}

\textbf{Application de la Méthode (3 questions) :}
\begin{enumerate}
    \item \textbf{Point d'opposition :} Le mot \og{} Technique \fg{}. Thèse : \og{} Technique = ensemble d'outils/moyens neutres \fg{}. Antithèse : \og{} Technique = Système autonome / Milieu totalitaire \fg{}.
    \item \textbf{Concept résolveur :} \textbf{Distinction conceptuelle} (Simondon / Ellul) entre \emph{Technique comme moyen (objet technique)} et \emph{Technique comme système (système technicien)}. + \textbf{Condition politique} : La maîtrise collective (droit, éthique, choix démocratiques) du système.
    \item \textbf{Thèse unifiée :} \og{} La technique \emph{outil} libère la puissance d'agir ; le \emph{système technique} tend à asservir par sa logique propre ; la liberté réelle réside dans la \textbf{maîtrise politique et éthique} du système technique par l'homme. \fg{}
\end{enumerate}

\textbf{Proposition de rédaction (Synthèse - 1 paragraphe dense) :}
\tcblower
\textbf{Solution détaillée :}

Au terme de cette analyse, il apparaît que la contradiction ne porte pas sur la technique \emph{en général}, mais sur l'équivocité du concept lui-même. La thèse avait raison de souligner que la technique, entendue comme \textbf{ensemble de moyens techniques} (outils, machines, procédés), augmente objectivement la puissance d'agir de l'homme sur la nature et libère du temps vital : c'est la définition bergsonienne de l'\emph{homo faber}, confirmée par Simondon quand il montre que l'objet technique \og{} concret \fg{} (abouti, transparent à son usage) est un facteur d'autonomie. L'antithèse, quant à elle, visait juste en dévoilant la logique du \textbf{système technicien} (Ellul) : lorsque les moyens deviennent autonomes, s'enchaînent par la seule loi de l'efficacité et s'imposent comme un milieu total, la technique cesse d'être un moyen pour devenir une fin en soi, asservissant l'homme à sa propre création. \textbf{La synthèse opère donc une distinction conceptuelle décisive :} la technique libère \emph{en tant qu'outil maîtrisé} ; elle asservit \emph{en tant que système autonome}. Dès lors, la liberté ne réside ni dans le refus luddite de la technique, ni dans l'adoration béate du progrès, mais dans la \textbf{maîtrise politique, juridique et éthique du système technique} (régulation, choix démocratiques des fins, éducation technique critique \`a la Simondon). \textbf{La technique libère donc \emph{sous condition} : celle d'une société qui refuse de confondre l'efficacité technique et la finalité humaine.}
\end{exoresolu}

\courssub{Lien organique avec la Conclusion : La Synthèse \emph{est} la réponse}

\begin{aretenir}[Synthèse vs Conclusion : Distinction fonctionnelle]
\begin{center}
\begin{tabularx}{\linewidth}{|l|X|X|}\hline
\rowcolor{popblueL}
\textbf{Critère} & \textbf{Synthèse (Dernier § Développement)} & \textbf{Conclusion (Dernière partie copie)} \\ \hline
\textbf{Fonction} & \textbf{Résout} le problème philosophique (Aufheben). & \textbf{Scelle} la réponse et \textbf{ouvre}. \\ \hline
\textbf{Contenu} & Argumentation conceptuelle dense (Distinction/Condition/Structure). & Reformulation synthétique de la thèse finale + Portée. \\ \hline
\textbf{Ton} & \og{} Il apparaît que... / C'est pourquoi... \fg{} (Démonstratif). & \og{} En définitive... / Au total... \fg{} (Assertif / Réflexif). \\ \hline
\textbf{Ouverture} & Interne (résolution de la dialectique interne au sujet). & Externe (autre sujet, autre auteur, actualité, limite). \\ \hline
\end{tabularx}
\end{center}
\textbf{Règle d'or :} Si votre conclusion apporte un \emph{nouvel} élément de réponse absent de la synthèse, votre développement est \textbf{inachevé}. La conclusion ne fait que \textbf{réécrire la phrase finale de la synthèse} en style \og{} clôture \fg{} et ajouter l'ouverture.
\end{aretenir}

\begin{savaistu}[Le poids de la Synthèse au Barème (Indicatif Bac Tech)]
\begin{itemize}
    \item La dissertation est notée sur \textbf{20 points} : \og{} Argumentation \fg{} (8-10 pts) + \og{} Structuration \fg{} (6-8 pts) + \og{} Connaissances \fg{} (4-6 pts).
    \item La \textbf{Synthèse} est le \textbf{point de cristallisation} de l'Argumentation et de la Structuration.
    \item \textbf{Absence de Synthèse identifiable} (pas de paragraphe distinct, pas de résolution) $\rightarrow$ Plafond \textbf{11/20} (même avec de bonnes parties Thèse/Antithèse). Le correcteur ne peut pas mettre \og{} Très bien \fg{} en structuration si le mouvement dialectique n'est pas bouclé.
    \item \textbf{Synthèse de qualité (Distinction nette / Réconciliation rigoureuse / Dépassement pertinent)} $\rightarrow$ Accessit direct aux notes \textbf{14-16+/20} (si connaissances correctes).
    \item \emph{Conseil stratégique :} Sacrifiez 2-3 lignes d'exemples en Thèse/Antithèse pour gagner 5 lignes d'une Synthèse \textbf{parfaite}. C'est le meilleur retour sur investissement temporel.
\end{itemize}
\end{savaistu}

\begin{savaistu}[Auteurs clés pour nourrir la Synthèse en Terminale Tech]
\begin{itemize}
    \item \textbf{Gilbert Simondon (\og{} Du mode d'existence des objets techniques \fg{}, 1958) :} \cle{Individuation technique}. Distinction \og{} Objet technique abstrait/aliéné \fg{} (pièces détachées, opacité) vs \og{} Objet technique concret \fg{} (surdéterminé, transparent, fiable). La culture technique (éducation, maintenance, invention) fait passer du premier au second. \emph{Arme absolue pour les synthèses par distinction sur la technique.}
    \item \textbf{Jacques Ellul (\og{} La Technique ou l'enjeu du siècle \fg{}, 1954) :} \cle{Système technicien}. La technique n'est plus un ensemble de moyens, c'est un \og{} milieu \fg{} autonome, totalitaire, qui détermine les fins (l'efficacité). L'homme n'est plus le sujet mais l'objet de la technique. \emph{Base théorique de l'Antithèse forte et de la Synthèse par condition politique (résistance/démocratie).}
    \item \textbf{Paulin Hountondji (\og{} La technique et le développement en Afrique \fg{}, 1980s) :} Critique de l'importation technologique aliénante (\og{} technique greffée \fg{}). Plaidoyer pour un \cle{développement endogène} : appropriation cognitive, adaptation au milieu local, innovation propre. \emph{Exemple concret camerounais/africain idéal pour illustrer la Synthèse : \og{} Maîtrise locale vs Dépendance externe \fg{}.}
    \item \textbf{Hans Jonas (\og{} Le principe responsabilité \fg{}, 1979) :} \cle{Éthique de la responsabilité} face à la puissance technique (nucléaire, génétique, IA). \og{} Agis de façon que les effets de ton action soient compatibles avec la permanence d'une vie authentiquement humaine. \fg{} \emph{Fondement éthique de la condition de validité (Réconciliation).}
\end{itemize}
\end{savaistu}

\courssub{Entraînement ciblé : De l'Antithèse à la Synthèse}

\begin{exercice}[Atelier \og{} Trouver le concept résolveur \fg{}]
Pour chacun des couples Thèse/Antithèse ci-dessous (issus de sujets probables Bac Tech), identifiez : (a) Le type de synthèse visé (Distinction / Réconciliation / Dépassement), (b) Le concept clé (le \og{} tiers \fg{}) qui résout l'opposition, (c) La phrase d'amorce type (\og{} Au terme de cette analyse... \fg{}).

\begin{enumerate}
    \item \textbf{Sujet : \og{} Le travail est-il une contrainte ? \fg{}}
        \begin{itemize}
            \item Thèse : Oui, contrainte naturelle (besoins), contrainte sociale (salariat, subordination), pénibilité.
            \item Antithèse : Non, réalisation de soi (Hegel/Marx jeune), activité anthropologique fondatrice, lien social, créativité.
        \end{itemize}
    \item \textbf{Sujet : \og{} L'État garantit-il la justice ? \fg{}}
        \begin{itemize}
            \item Thèse : Oui, monopole violence légitime, lois, tribunaux, protection faibles, redistribution.
            \item Antithèse : Non, justice de classe (Marx), inégalités structurelles, justice formelle vs réelle, corruption, État failli.
        \end{itemize}
    \item \textbf{Sujet : \og{} L'histoire a-t-elle un sens ? \fg{}}
        \item Thèse : Oui, progrès (Condorcet), esprit hégélien, téléologie, émancipation.
        \item Antithèse : Non, contingence, cycles, absurdité (Camus), catastrophes, pas de finalité immanente.
    \item \textbf{Sujet : \og{} La technique détruit-elle la nature ? \fg{}}
        \item Thèse : Oui, extraction, pollution, artificialisation, Anthropocène, technique = anti-nature (Heidegger : \og{} Gestell \fg{}).
        \item Antithèse : Non, technique permet protection (énergies renouvelables, dépollution), nature n'existe plus \og{} vierge \fg{}, technique = prolongement nature (Simondon, Stiegler).
\end{enumerate}
\end{exercice}

\begin{corrige}[Exercice n°1 - Éléments de réponse]
\begin{enumerate}
    \item \textbf{Travail/Contrainte :} \textbf{Distinction} (Travail \emph{activité anthropologique} vs Travail \emph{emploi salarié/aliéné}). Concept : \og{} Travail libre \fg{} (artisan, artiste, care) vs \og{} Travail contraint \fg{} (salariat taylorien). \og{} Au terme de cette analyse, il apparaît que la contrainte ne réside pas dans le travail \emph{comme tel} (activité vitale), mais dans sa \emph{forme sociale aliénée} (division, subordination). \fg{}
    \item \textbf{État/Justice :} \textbf{Réconciliation conditionnelle}. Concept : \og{} État de droit \fg{} (séparation pouvoirs, indépendance justice, démocratie effective) vs \og{} État légal \fg{} (lois formelles seulement). Condition : \textbf{Democraticité effective}. \og{} La contradiction se résout si l'on distingue la justice \emph{formelle} (lois) de la justice \emph{réelle} (effectivité des droits) : l'État garantit la première, mais la seconde exige une démocratie vigilante. \fg{}
    \item \textbf{Histoire/Sens :} \textbf{Dépassement} (ou Réconciliation critique). Concept : \og{} Sens \emph{immanent} \fg{} (donné d'avance) vs \og{} Sens \emph{construit} \fg{} (par les hommes). Auteurs : Ricoeur, Arendt, Merleau-Ponty. \og{} C'est que, radicalement, l'histoire n'a pas de sens \emph{donné} (théologique/métaphysique), mais elle \emph{devient} sens par la liberté humaine qui l'interprète et l'oriente. \fg{}
    \item \textbf{Technique/Nature :} \textbf{Distinction + Dépassement (Simondon/Stiegler)}. Concept : \og{} Technique \emph{contre} nature \fg{} (Heidegger/Ellul) vs \og{} Technique \emph{comme} nature \fg{} (prothèse, organe artificiel, milieu technique). Distinction : \og{} Technologie prédatrice \fg{} vs \og{} Technologie symbiotique \fg{}. \og{} Au terme de cette analyse, il apparaît que la technique ne détruit la nature que si elle est conçue comme \emph{maîtrise/extraction} (Gestell) ; elle peut au contraire \emph{régénérer} la nature si elle s'inscrit dans une logique d'individuation concrète et de soin du milieu (pharmacologie positive, Stiegler). \fg{}
\end{enumerate}
\end{corrige}

\begin{methode}[Auto-évaluation Synthèse (Check-list 30 secondes avant de rendre la copie)]
Cochez mentalement chaque case. Si une case est vide $\rightarrow$ Réécrivez le paragraphe.
\begin{enumerate}
    \item [ ] \textbf{Présence visuelle :} Un paragraphe distinct (ou deux courts), bien séparé de l'Antithèse (saut de ligne + alinéa).
    \item [ ] \textbf{Amorce dialectique :} Commence par \og{} Au terme de cette analyse... \fg{} / \og{} La contradiction se résout... \fg{} / \og{} C'est que... \fg{} (PAS \og{} En synthèse... \fg{}).
    \item [ ] \textbf{Zéro nouveauté :} Aucun nouvel auteur, aucune nouvelle citation, aucun nouvel exemple concret.
    \item [ ] \textbf{Concept résolveur explicite :} Une distinction nette (\og{} Il faut distinguer X de Y \fg{}) OU une condition claire (\og{} Pourvu que... \fg{}) OU une structure ontologique (\og{} C'est que l'homme est... \fg{}).
    \item [ ] \textbf{Relecture unificatrice :} La phrase \og{} La thèse avait raison en ce que... ; l'antithèse complétait en montrant que... \fg{} (ou équivalent) est présente.
    \item [ ] \textbf{Réponse finale affirmative :} La dernière phrase répond \textbf{précisément} au sujet (mot pour mot si possible) par une thèse nuancée.
    \item [ ] \textbf{Lien Conclusion :} La première phrase de ma Conclusion future reprendra \textbf{exactement} la thèse finale de cette synthèse.
\end{enumerate}
\end{methode}

\begin{exemplebox}[Modèle de phrase de transition Antithèse $\rightarrow$ Synthèse (à adapter)]
\og{} Cependant, opposer frontalement la thèse et l'antithèse revient à maintenir une alternative stérile. \textbf{Au terme de cette analyse, il apparaît que} l'opposition ne porte pas sur [Sujet] en général, mais sur la confusion entre [Concept A] et [Concept B]. \fg{}
\emph{(Variante Réconciliation) :} \og{} Cependant, cette opposition trouve son unité dès lors que l'on pose la condition suivante : [Condition C]. \fg{}
\emph{(Variante Dépassement) :} \og{} C'est que, radicalement, [Sujet] ne saurait se comprendre hors de la structure fondamentale qu'est [Structure S]. \fg{}
\end{exemplebox}

\courssub{Résumé visuel : Le Triangle Dialectique (Mémorisation)}

\begin{center}
\begin{popfigure}
\begin{tikzpicture}[
    scale=1.1,
    every node/.style={font=\small},
    triangle node/.style={draw=popdark, thick, fill=popcream, rounded corners=4pt, inner sep=6pt, text width=3.2cm, align=center, minimum height=2.2cm},
    arrow main/.style={-Stealth, very thick, popblue},
    arrow aufheben/.style={-Stealth, thick, densely dashed, poporange},
    label box/.style={fill=popblue!10, draw=popblue, rounded corners=3pt, inner sep=3pt, font=\tiny\bfseries, text=popdark}
]
% Sommets du triangle
\node[triangle node, align=center] (thesis) at (-3, 2){\textbf{THÈSE}\\\emph{(Moment affirmatif)}\\\og{} Oui, car... \fg{}\\\textit{Arguments I-A-E}};
\node[triangle node, align=center] (antithesis) at (3, 2){\textbf{ANTITHÈSE}\\\emph{(Moment négatif)}\\\og{} Non, car... \fg{}\\\textit{Arguments I-A-E}};
\node[triangle node, align=center] (synthesis) at (0, -2.5){\textbf{SYNTHÈSE}\\\emph{(Moment résolutif)}\\\og{} Au terme... \fg{}\\\textit{Distinction / Condition / Structure}\\\textbf{Réponse finale}};

% Flèches principales
\draw[arrow main] (thesis) -- node[above, label box] {S'oppose} (antithesis);
\draw[arrow main] (thesis) -- node[left, label box, rotate=30] {Dépasse \& Conserve} (synthesis);
\draw[arrow main] (antithesis) -- node[right, label box, rotate=-30] {Dépasse \& Conserve} (synthesis);

% Flèches Aufheben (retour conceptuel)
\draw[arrow aufheben] (synthesis) to[bend left=20] node[left=2pt, label box] {Relit \& Intègre} (thesis);
\draw[arrow aufheben] (synthesis) to[bend right=20] node[right=2pt, label box] {Relit \& Intègre} (antithesis);

% Annotation centrale
\node[align=center, font=\footnotesize\itshape, text=popmuted] at (0, 0.2) {\textbf{Aufheben}\\(Supprimer / Conserver / Élever)};

% Légende bas
\node[anchor=north, font=\footnotesize, text width=12cm, align=center] at (0, -4.8) {
\textbf{Thèse} = Vérité partielle (Unilatéralité) \quad $\rightarrow$ \quad \textbf{Antithèse} = Vérité partielle inverse (Négation) \quad $\rightarrow$ \quad \textbf{Synthèse} = Vérité concrète (Négation de la négation)
};

\end{tikzpicture}
\legende{Le triangle dialectique : la Synthèse ne se place pas \og{} après \fg{} chronologiquement, mais \og{} au-dessus \fg{} logiquement. Elle relie les deux sommets par une double flèche \og{} Dépasse \& Conserve \fg{} (Aufheben).}
\end{popfigure}
\end{center}

\begin{attention}[Dernier rappel avant la Section 5 (Conclusion)]
La Synthèse est le \textbf{cœur nucléaire} de votre dissertation. C'est le seul endroit où vous \textbf{prouvez} que vous savez \emph{philosopher} (conceptualiser, distinguer, unifier) et pas seulement \emph{réciter} des arguments. Soignez sa rédaction : \textbf{clarté conceptuelle, densité argumentative, absence de scories}. La Section 5 vous montrera comment en faire le socle inébranlable de votre Conclusion.
\end{attention}

\coursec{La conclusion : bilan définitif et ouverture maîtrisée}

La synthèse dialectique, que nous venons d'étudier, a dépassé l'opposition entre la thèse et l'antithèse pour proposer une réponse élaborée au problème posé. Mais cette réponse ne peut pas rester suspendue dans le dernier paragraphe du développement : il faut encore la \emph{sceller}, la formuler une dernière fois de manière définitive, puis ouvrir l'esprit du lecteur vers un nouvel horizon. C'est précisément le rôle de la \cle{conclusion}. Trop souvent négligée, écrite à la hâte dans les cinq dernières minutes de l'épreuve, elle est pourtant la dernière impression laissée au correcteur — et l'on sait qu'en philosophie comme dans la vie, \emph{la dernière impression compte autant que la première}.

\courssub{Qu'est-ce qu'une conclusion philosophique ?}

\textbf{Intuition.} Imaginez un avocat qui, après trois heures de plaidoirie brillante, quitterait la salle d'audience sans mot dire, sans résumer sa demande au tribunal. Son travail, aussi remarquable soit-il, resterait inachevé : les juges n'auraient pas entendu le verdict qu'il sollicite. La conclusion est ce moment où le philosophe en herbe « demande le verdict » : il rappelle au lecteur la réponse que tout son raisonnement a construite, puis il montre que cette réponse n'épuise pas la pensée, qu'elle en appelle d'autres.

\begin{definition}[La conclusion philosophique]
La \cle{conclusion} est la dernière partie de la dissertation philosophique. Elle comprend \textbf{deux temps obligatoires et distincts} :
\begin{enumerate}
\item le \cle{bilan} : une réponse claire, nuancée et définitive à la problématique annoncée dans l'introduction, qui reprend le fil de la synthèse sans la copier mot à mot ;
\item l'\cle{ouverture} : une question nouvelle, brève et pertinente, qui prolonge la réflexion vers un autre domaine, une autre discipline ou une actualité, sans prétendre la traiter.
\end{enumerate}
Sa longueur idéale est de \textbf{10 à 15 lignes}. Elle se termine par une phrase percutante qui marque durablement le correcteur.
\end{definition}

\begin{attention}[Deux temps, pas un seul]
L'erreur la plus fréquente est de réduire la conclusion au seul bilan (« En conclusion, la technique libère et asservit l'homme »). Une conclusion sans ouverture est une conclusion incomplète : elle ferme la pensée au lieu de la maintenir vivante. Inversement, une ouverture sans bilan laisse le lecteur sans réponse au problème posé. \textbf{Les deux temps sont exigibles au Baccalauréat technique.}
\end{attention}

\courssub{Premier temps : le bilan, ou la réponse définitive}

\textbf{Ce qu'il faut comprendre.} Le bilan n'est pas un résumé de la dissertation. Le correcteur sait ce que vous avez écrit : il vient de le lire. Le bilan est la \emph{réponse} que tout le parcours a rendue possible. Si l'introduction posait la question « La technique libère-t-elle l'homme ? », le bilan doit répondre : oui, non, ou — le plus souvent en philosophie — \emph{oui, mais à certaines conditions}.

\begin{propriete}[Règle d'or du bilan]
La conclusion ne doit rien contenir qui ne soit \textbf{déjà présent dans la synthèse}. Le bilan condense la synthèse en une ou deux phrases ; il n'y ajoute ni argument, ni exemple, ni auteur nouveau. Toute idée apparaissant pour la première fois en conclusion est une faute méthodologique : soit elle aurait dû être développée (et sa présence révèle un défaut de plan), soit elle est superflue (et elle affaiblit la fin du devoir).
\end{propriete}

\textbf{La forme du bilan.} Le bilan idéal tient en \emph{une phrase unique}, longue mais construite, qui articule les deux pôles du problème et leur dépassement. Cette phrase doit être :

\begin{itemize}
\item \textbf{nuancée} : elle évite les réponses brutales (« la technique libère l'homme » tout court) au profit d'une réponse conditionnelle ou dialectique ;
\item \textbf{impersonnelle} : on ne dit jamais « \emph{Je pense que...} », « \emph{À mon avis...} », « \emph{Personnellement...} ». La philosophie ne s'intéresse pas à votre opinion, mais à la vérité que le raisonnement établit. Écrivez : « Il apparaît que... », « La technique libère... », jamais « Je crois que la technique libère... » ;
\item \textbf{reliée à la problématique} : elle reprend les termes mêmes de la question posée dans l'introduction, pour montrer que la boucle est bouclée.
\end{itemize}

\begin{exemplebox}[Un bilan réussi]
Sujet : \emph{La technique libère-t-elle l'homme ?}

« Ainsi, la technique libère l'homme en tant qu'outil d'émancipation placé sous le contrôle de sa volonté et de sa raison, mais elle l'asservit dès qu'elle devient un système autonome échappant à toute gouvernance politique et morale. »

On remarque : le mot « Ainsi » qui annonce le bilan ; la nuance (libération \emph{conditionnelle}) ; l'absence totale de « je » ; la reprise des termes du sujet (« libère », « l'homme ») ; l'articulation des moments de la thèse (outil d'émancipation) et de l'antithèse (système autonome) dépassés dans la synthèse (la question du contrôle politique).
\end{exemplebox}

\begin{methode}[Rédiger le bilan en trois étapes]
\begin{enumerate}
\item \textbf{Relire la problématique} de l'introduction et la recopier mentalement : à quelle question dois-je répondre ?
\item \textbf{Relire la phrase de synthèse} (dernier paragraphe du développement) : quelle position dépassée ai-je défendue ?
\item \textbf{Formuler la réponse en une phrase} commençant par « Ainsi » ou « Il résulte de cette analyse que », articulant la nuance (« en tant que... mais dès que... », « si... alors... sinon... »), sans « je », sans nouvel argument.
\end{enumerate}
\end{methode}

\courssub{Deuxième temps : l'ouverture, ou la pensée qui continue}

\textbf{Intuition.} Une bonne conversation ne s'achève jamais sur un point final sec : elle se termine souvent par « cela me fait penser à... ». L'ouverture est ce geste intellectuel : après avoir répondu à une question, le philosophe montre que cette réponse en soulève une autre. Il prouve ainsi qu'il n'a pas seulement traité un sujet, mais qu'il \emph{pense}.

\begin{definition}[L'ouverture]
L'\cle{ouverture} est une question nouvelle, formulée en \textbf{2 à 3 lignes maximum}, qui prolonge la réflexion sans la traiter. Elle peut s'orienter vers :
\begin{itemize}
\item un \textbf{prolongement du problème} (une dimension non étudiée du sujet) ;
\item une \textbf{autre discipline} (sociologie, anthropologie, droit, économie, histoire) ;
\item une \textbf{actualité} ou une situation concrète (enjeux contemporains, contexte camerounais) ;
\item un \textbf{auteur non traité} dans le devoir, dont la pensée permettrait d'aller plus loin.
\end{itemize}
\end{definition}

\begin{propriete}[Conditions de validité de l'ouverture]
Une ouverture est recevable si et seulement si elle est :
\begin{enumerate}
\item \textbf{pertinente} : elle naît réellement du sujet traité, elle n'est pas plaquée artificiellement ;
\item \textbf{brève} : 2 à 3 lignes, jamais un nouveau paragraphe de développement ;
\item \textbf{non traitée} : on la pose, on ne la résout pas — sinon ce serait un nouvel argument, interdit en conclusion ;
\item \textbf{hors du sujet immédiat mais dans son horizon} : elle ne doit pas être le sujet voisin qu'on aurait pu traiter à la place.
\end{enumerate}
\end{propriete}

\begin{exemplebox}[Une ouverture réussie (contexte camerounais)]
« Dès lors, cette analyse invite à repenser l'éducation technique au Cameroun : ne s'agit-il pas de former des ingénieurs-citoyens, capables non seulement de maîtriser la technique, mais de la gouverner démocratiquement ? »

On remarque : le connecteur « Dès lors » qui lie l'ouverture au bilan ; la brièveté (deux lignes) ; la forme interrogative (on ne traite pas la question) ; la pertinence (l'éducation prolonge naturellement le problème du contrôle de la technique) ; l'ancrage concret dans la réalité nationale.
\end{exemplebox}

\begin{exemplebox}[Une ouverture vers une autre discipline (exemple camerounais)]
Toujours sur le sujet de la technique : « Il reste alors à interroger la place des savoirs endogènes — pharmacopée traditionnelle, architecture vernaculaire en banco, techniques agricoles locales — face à la technique importée : l'anthropologie et la sociologie pourraient nous apprendre si ces savoirs constituent eux aussi des techniques au sens plein, et ce que leur marginalisation révèle de notre rapport à la modernité. »

Cette ouverture prolonge le sujet vers l'anthropologie, évoque une réalité camerounaise concrète (la pharmacopée, le banco) et pose une vraie question sans y répondre.
\end{exemplebox}

\begin{attention}[Le piège classique : l'ouverture « clé sur porte »]
L'ouverture la plus fréquente — et la plus sanctionnée — est celle qui consiste à poser \emph{le sujet voisin} : « On pourrait aussi se demander si la science libère l'homme... » alors que le sujet portait sur la technique. C'est ce qu'on appelle une ouverture « clé sur porte » : on referme la porte du sujet en montrant simplement la porte d'à côté. Une telle ouverture est \textbf{non problématisée} : elle ne naît pas du raisonnement mené, elle flotte à côté. Une bonne ouverture doit pouvoir justifier son lien avec la conclusion : « \emph{Puisque} nous avons montré que la technique exige une gouvernance politique, \emph{il reste à} se demander... »
\end{attention}

\begin{attention}[Les interdits absolus de la conclusion]
\begin{itemize}
\item \textbf{Aucun nouvel argument} : tout argument nouveau prouve que le développement était incomplet.
\item \textbf{Aucune citation nouvelle non analysée} : plaquer une citation en conclusion sans l'expliquer est une faute grave (voir la section sur les citations).
\item \textbf{Aucun lieu commun} : « Ce sujet est vaste... », « Cette question reste ouverte depuis la nuit des temps... », « Chacun aura son opinion... » — ces formules vides agacent le correcteur et ne disent rien.
\item \textbf{Aucun retour au sujet suivant} : ne pas confondre ouverture et annonce d'un autre devoir.
\item \textbf{Aucune formule personnelle} : « Je pense que », « Pour moi », « En tant que jeune camerounais, je... ».
\end{itemize}
\end{attention}

\courssub{Schéma de la structure de la conclusion}

\begin{popfigure}
\begin{center}
\begin{tikzpicture}[node distance=0.9cm and 2.6cm,
  bloc/.style={rectangle, rounded corners=4pt, draw=popdark, thick, fill=popblueL, text width=4.6cm, align=center, minimum height=1.4cm, font=\small},
  bilan/.style={rectangle, rounded corners=4pt, draw=popgold, thick, fill=popgoldL, text width=4.6cm, align=center, minimum height=1.4cm, font=\small},
  ouvert/.style={rectangle, rounded corners=4pt, draw=popgreen, thick, fill=popgreenL, text width=4.6cm, align=center, minimum height=1.4cm, font=\small},
  horizon/.style={rectangle, rounded corners=4pt, draw=poppurple, thick, fill=poppurpleL, text width=4.2cm, align=center, minimum height=1.4cm, font=\small},
  fleche/.style={-{Stealth[length=3mm]}, very thick, popdark}]
\node[bloc, align=center] (synth){\textbf{Synthèse}\\ (fin du développement)\\ position dépassée};
\node[bilan, right=of synth, align=center] (bil){\textbf{Bilan}\\ réponse nuancée à la problématique (1 phrase)};
\node[ouvert, below=of bil, align=center] (ouv){\textbf{Ouverture}\\ question nouvelle, brève, non traitée};
\node[horizon, left=of ouv, align=center] (hor){\textbf{Nouvel horizon}\\ autre discipline, actualité, auteur non traité};
\draw[fleche] (synth) -- node[above, font=\small\itshape]{résout} (bil);
\draw[fleche] (bil) -- node[right, font=\small\itshape]{puis} (ouv);
\draw[fleche] (ouv) -- node[above, font=\small\itshape]{prolonge} (hor);
\end{tikzpicture}
\end{center}
\legende{Structure de la conclusion : le bilan \emph{résout} la problématique en condensant la synthèse ; l'ouverture \emph{prolonge} la réflexion vers un nouvel horizon sans le traiter.}
\end{popfigure}

\courssub{Un exemple complet annoté}

Voici une conclusion type, rédigée pour le sujet « La technique libère-t-elle l'homme ? », puis annotée visuellement.

\begin{exemplebox}[Conclusion type (sujet : La technique libère-t-elle l'homme ?)]
« Ainsi, la technique libère l'homme en tant qu'outil d'émancipation placé sous le contrôle de sa volonté et de sa raison, mais elle l'asservit dès qu'elle devient un système autonome échappant à toute gouvernance politique et morale. La question n'est donc plus de savoir s'il faut accepter ou refuser la technique — ce dilemme est dépassé — mais de déterminer \emph{qui} la gouverne et \emph{au nom de quelles valeurs}. Dès lors, cette analyse invite à repenser l'éducation technique au Cameroun : ne s'agit-il pas de former des ingénieurs-citoyens, capables non seulement de produire, mais de décider collectivement du sens du progrès ? La technique ne libère que celui qui refuse d'être son instrument. »
\end{exemplebox}

\begin{popfigure}
\begin{center}
\begin{tikzpicture}
\node[rectangle, draw=popgold, thick, fill=popgoldL, rounded corners=3pt, text width=11.5cm, align=justify, font=\small] (b) at (0,0) {\textbf{[Bilan]} « Ainsi, la technique libère l'homme en tant qu'outil d'émancipation placé sous le contrôle de sa volonté et de sa raison, mais elle l'asservit dès qu'elle devient un système autonome échappant à toute gouvernance politique et morale. »};
\node[rectangle, draw=popgreen, thick, fill=popgreenL, rounded corners=3pt, text width=11.5cm, align=justify, font=\small] (o) at (0,-2.6) {\textbf{[Ouverture]} « Dès lors, cette analyse invite à repenser l'éducation technique au Cameroun : ne s'agit-il pas de former des ingénieurs-citoyens, capables non seulement de produire, mais de décider collectivement du sens du progrès ? »};
\node[rectangle, draw=popink, thick, fill=popcream, rounded corners=3pt, text width=11.5cm, align=justify, font=\small] (p) at (0,-4.9) {\textbf{[Phrase finale percutante]} « La technique ne libère que celui qui refuse d'être son instrument. »};
\end{tikzpicture}
\end{center}
\legende{Anatomie d'une conclusion réussie : en jaune, le bilan nuancé qui condense la synthèse ; en vert, l'ouverture brève et pertinente ; la phrase finale, courte et mémorable, frappe l'esprit du correcteur.}
\end{popfigure}

\begin{methode}[L'astuce des connecteurs de conclusion]
Pour baliser clairement les deux temps, utilisez des connecteurs conventionnés que le correcteur repérera immédiatement :
\begin{itemize}
\item pour lancer le \textbf{bilan} : « \cle{Ainsi}, ... », « Il résulte de cette analyse que... », « Au terme de ce parcours, ... » ;
\item pour lancer l'\textbf{ouverture} : « \cle{Dès lors}, ... », « \cle{Il reste à} se demander... », « Cette réponse en appelle une autre : ... ».
\end{itemize}
Ces mots sont des signaux méthodologiques : ils montrent au correcteur que vous maîtrisez la structure canonique.
\end{methode}

\begin{methode}[Rédiger la conclusion sur le brouillon AVANT la rédaction finale]
Voici une stratégie de candidat efficace, en quatre étapes :
\begin{enumerate}
\item \textbf{Au brouillon}, juste après avoir construit votre plan détaillé, rédigez la problématique de l'introduction \emph{et} la phrase de bilan de la conclusion, l'une en face de l'autre.
\item \textbf{Vérifiez la cohérence} : la phrase de bilan répond-elle exactement à la question posée ? Si non, ajustez l'un ou l'autre avant de rédiger.
\item \textbf{Rédigez le devoir} au propre, en gardant ce fil : chaque partie doit rapprocher de la réponse annoncée.
\item \textbf{En fin d'épreuve}, reprenez la phrase de bilan du brouillon, reformulez-la pour tenir compte du chemin réellement parcouru, puis ajoutez l'ouverture.
\end{enumerate}
Cette méthode garantit que l'introduction et la conclusion se répondent : c'est le critère de \emph{cohérence} du barème.
\end{methode}

\begin{exoresolu}[Corriger une conclusion défaillante]
\textbf{Énoncé.} Voici la conclusion d'un élève sur le sujet « La technique libère-t-elle l'homme ? ». Relevez toutes les fautes méthodologiques, puis proposez une version corrigée.

« En conclusion, je pense que la technique est à la fois bonne et mauvaise. Ce sujet est très vaste et les philosophes ne sont pas d'accord depuis la nuit des temps. Comme le dit un proverbe, "toute médaille a son revers". On pourrait aussi se demander si la science libère l'homme, ce qui est un autre sujet intéressant. D'ailleurs, l'intelligence artificielle pose aujourd'hui des problèmes nouveaux qu'il faudrait analyser en détail, car les robots pourraient un jour remplacer les ouvriers dans les usines de Douala, ce qui créerait du chômage et poserait la question du revenu universel. »
\tcblower
\textbf{Solution détaillée.} Relevé des fautes :
\begin{enumerate}
\item « \emph{Je pense que} » : formule personnelle interdite. Le bilan doit être impersonnel.
\item « \emph{bonne et mauvaise} » : bilan non nuancé philosophiquement, simple juxtaposition sans dépassement ; la synthèse n'apparaît pas.
\item « \emph{Ce sujet est très vaste... nuit des temps} » : lieu commun vide, interdit.
\item « \emph{toute médaille a son revers} » : citation (proverbe) plaquée, non analysée, et sans valeur philosophique.
\item « \emph{si la science libère l'homme} » : ouverture « clé sur porte » — sujet voisin, non problématisé.
\item Le dernier développement sur l'intelligence artificielle et le chômage à Douala introduit un \textbf{nouvel argument} et un \textbf{nouvel exemple} : c'est un paragraphe de développement déguisé, alors que l'ouverture doit tenir en 2-3 lignes et ne rien traiter.
\end{enumerate}
\textbf{Version corrigée} : « Ainsi, la technique libère l'homme lorsqu'elle demeure un instrument au service de ses fins, mais elle l'asservit dès qu'elle s'érige en système autonome imposant ses propres finalités. Tout l'enjeu est donc politique : il s'agit moins de juger la technique que de la gouverner. Dès lors, il reste à se demander si l'école technique camerounaise forme aujourd'hui de simples exécutants ou de véritables citoyens capables de décider du sens du progrès. »
\end{exoresolu}

\begin{savaistu}[Le coût réel d'une conclusion bâclée]
Au Baccalauréat technique, le barème de la dissertation réserve des points précis à la \emph{structuration} du devoir et à la \emph{qualité de l'expression}. Une conclusion absente, réduite à une ligne (« Donc la technique libère et asservit l'homme ») ou truffée de lieux communs fait perdre de \textbf{2 à 3 points} sur la note finale — soit souvent la différence entre une mention et un simple passable. À l'inverse, une conclusion nette, en deux temps, terminée par une phrase percutante, laisse au correcteur une impression de maîtrise qui rejaillit sur l'appréciation globale du devoir. Les cinq dernières minutes de l'épreuve doivent être \emph{sacrées} : elles appartiennent à la conclusion.
\end{savaistu}

\begin{aretenir}[L'essentiel de la conclusion]
\begin{itemize}
\item La conclusion comprend \textbf{deux temps obligatoires} : le \cle{bilan} (réponse nuancée à la problématique, condensé de la synthèse) et l'\cle{ouverture} (question nouvelle, brève, pertinente, non traitée).
\item Le bilan tient idéalement en \textbf{une phrase}, impersonnelle (jamais « je pense »), introduite par « Ainsi ».
\item L'ouverture tient en \textbf{2 à 3 lignes}, introduite par « Dès lors » ou « Il reste à », et ne doit jamais être le sujet voisin (« clé sur porte »).
\item Interdits : nouvel argument, citation non analysée, lieu commun, formule personnelle.
\item Longueur totale : \textbf{10 à 15 lignes}, avec une dernière phrase courte et percutante.
\item Stratégie : rédiger le bilan \emph{au brouillon} avant la rédaction, pour garantir la cohérence avec l'introduction et la synthèse.
\end{itemize}
\end{aretenir}

\begin{exercice}[Entraînement au bilan et à l'ouverture]
Pour chacun des deux sujets suivants, rédigez une conclusion complète (10 à 15 lignes) respectant les deux temps obligatoires :
\begin{enumerate}
\item « Le travail n'est-il qu'une contrainte ? » (ouverture possible : vers la question du loisir ou de la formation professionnelle au Cameroun) ;
\item « L'État doit-il tout contrôler ? » (ouverture possible : vers la décentralisation et les collectivités locales camerounaises, ou vers la philosophie politique de John Locke, non traité dans le devoir).
\end{enumerate}
Vous veillerez à : formuler le bilan en une phrase nuancée commençant par « Ainsi » ; éviter tout « je » ; construire une ouverture problématisée en 2-3 lignes ; terminer par une phrase percutante.
\end{exercice}

\begin{corrige}[Exercice — Entraînement au bilan et à l'ouverture]
\textbf{Éléments de correction pour le sujet 1} (« Le travail n'est-il qu'une contrainte ? »).

\emph{Bilan attendu} : « Ainsi, le travail n'est une pure contrainte que lorsqu'il est subi comme une nécessité étrangère ; dès lors qu'il devient activité par laquelle l'homme se réalise, transforme la nature et se reconnaît dans son œuvre, il s'élève au rang de condition de la liberté. » (On vérifie : une phrase, « Ainsi », pas de « je », reprise des termes du sujet, dépassement de l'alternative contrainte/libération.)

\emph{Ouverture attendue} : « Dès lors, si le travail peut libérer, il reste à se demander quelle place accorder au loisir dans une vie accomplie : la formation professionnelle des jeunes Camerounais devrait-elle viser le seul rendement, ou aussi l'épanouissement de l'homme tout entier ? »

\emph{Critères d'évaluation} : bilan nuancé (2 points) ; impersonnalité (1 point) ; ouverture brève, problématisée et liée au bilan (2 points) ; phrase finale percutante (1 point). Toute conclusion contenant un lieu commun (« le travail, notion vaste et complexe... ») ou un nouvel argument perd les points correspondants.
\end{corrige}

\coursec{Les citations : Règles de choix, d'intégration et de référencement}

Après avoir maîtrisé l'art de conclure — ce bilan définitif qui scelle la démonstration et ouvre l'horizon —, il nous reste à perfectionner un outil indispensable du philosophe : l'usage des auteurs. Une dissertation sans références précises est comme un édifice sans fondations : elle peut tenir un temps, mais elle manque de solidité intellectuelle. Au Probatoire comme au Baccalauréat technique, l'examinateur attend que le candidat démontre sa culture philosophique non par du \og name-dropping \fg{} stérile, mais par une \textbf{intégration rigoureuse} des pensées majeures au cœur de son argumentation. Cette section transforme la citation d'un ornement en un \textbf{levier argumentatif}.

\courssub{Les quatre fonctions légitimes de la citation}

Une citation n'est jamais « décorative ». Elle remplit une fonction précise dans l'économie de la démonstration. En identifier la fonction \emph{avant} de l'insérer garantit sa pertinence.

\begin{definition}[Les quatre fonctions de la citation]
Dans une dissertation philosophique, toute citation doit assumer l'une de ces quatre fonctions :
\begin{enumerate}
    \item \textbf{L'autorité (appui) :} Elle invoque un penseur reconnu pour valider une thèse ou un concept clé. Exemple : citer Kant pour définir l'impératif catégorique.
    \item \textbf{L'illustration (concrétisation) :} Elle éclaire une idée abstraite par une formule imagée ou une analyse fine. Exemple : citer Bergson sur l'élan vital pour illustrer la créativité évolutive.
    \item \textbf{Le contre-point (objection) :} Elle expose une thèse adverse que l'on va ensuite réfuter ou nuancer (antithèse). Exemple : citer Hobbes sur l'état de nature pour le contester par Rousseau.
    \item \textbf{L'amorce (introduction) ou la chute (conclusion) :} Elle ouvre le sujet (accroche) ou le referme sur une résonance forte (ouverture), à condition qu'elle soit \emph{reprise} et \emph{analysée} immédiatement.
\end{enumerate}
\end{definition}

\begin{attention}[Citation décorative = Zéro point]
Placer « Socrate : Connais-toi toi-même » en exergue d'un sujet sur \og La technique libère-t-elle l'homme ? \fg{} sans aucun lien avec la problématique ni analyse ultérieure est une \textbf{citation décorative}. Elle signale au correcteur une culture de surface. \textbf{Sanction :} elle ne rapporte aucun point et alourdit la copie.
\end{attention}

\courssub{Critères de choix : Exactitude, Pertinence, Identification}

Choisir une citation, c'est faire un travail de philologue miniature. Trois critères sont non négociables.

\begin{propriete}[Critères de sélection d'une citation]
\begin{itemize}
    \item \textbf{Exactitude textuelle :} Citer le texte original (ou la traduction de référence : Bréhier pour Platon/Aristote, Vrin pour Descartes, PUF pour Kant/Hegel, etc.). Toute modification (modernisation de l'orthographe, suppression de mots) doit être signalée par des crochets [...] ou des points de suspension [...]. \textbf{Interdit :} la paraphrase présentée comme citation.
    \item \textbf{Pertinence argumentative :} La citation doit \og coller \fg{} à l'argument du paragraphe. Elle ne doit pas être un \og corps étranger \fg{}. Question test : \og Si je retire cette citation, mon argument tient-il encore ? \fg{} Si oui, la citation est superflue. Si non, elle est structurelle.
    \item \textbf{Identification rigoureuse :} Nom de l'auteur + Titre de l'œuvre (en italique). Exemple : \og Kant, \emph{Critique de la raison pure} \fg{}. Pas de note de bas de page en dissertation (gain de temps, lisibilité).
\end{itemize}
\end{propriete}

\begin{aretenir}[Corpus minimal « Bac Tech » — Auteurs incontournables]
Le programme technique exige une culture philosophique solide, ouverte sur la pensée africaine. Mémorisez au moins \textbf{une citation précise par auteur} pour les thèmes majeurs :
\begin{center}
\begin{tabularx}{\linewidth}{|l|X|X|}\hline
\textbf{Thème} & \textbf{Classiques Occidentaux} & \textbf{Auteurs Africains (Impératif)} \\ \hline
Technique & Descartes (\emph{Discours}, VI), Marx (\emph{Le Capital}, I), Ellul (\emph{Le Système technicien}), Heidegger (\emph{La Question de la technique}) & \textbf{Fabien Eboussi Boulaga} (\emph{La Crise du Muntu}), \textbf{Paulin Hountondji} (\emph{La Connaissance endogène}), \textbf{Valentin Mudimbe} (\emph{L'Invention de l'Afrique}) \\ \hline
Politique / Droit & Platon (\emph{République}), Aristote (\emph{Politique}), Hobbes (\emph{Léviathan}), Rousseau (\emph{Contrat social}), Kant (\emph{Projet de paix perpétuelle}) & \textbf{Marcien Towa} (\emph{Essai sur la problématique philosophique en Afrique}), \textbf{Kwame Nkrumah} (\emph{Consciencisme}), \textbf{Stanislas Adotevi} (\emph{Négritude et Négrologues}), \textbf{Achille Mbembe} (\emph{Critique de la raison nègre}) \\ \hline
Morale / Éthique & Kant (\emph{Fondements}), Bentham/Mill (\emph{Utilitarisme}), Sartre (\emph{L'Existentialisme est un humanisme}), Jonas (\emph{Le Principe responsabilité}) & \textbf{Eboussi Boulaga} (éthique de la libération), \textbf{Hountondji} (éthique de la responsabilité), \textbf{Towa} (éthique communautaire) \\ \hline
Connaissance / Vérité & Descartes (\emph{Méditations}), Kant (\emph{Critique}), Hegel (\emph{Phénoménologie}), Foucault (\emph{Les Mots et les Choses}) & \textbf{Hountondji} (critique de l'ethnophilosophie), \textbf{Mudimbe} (gnosis africaine), \textbf{Mbembe} (politique de la connaissance) \\ \hline
\end{tabularx}
\end{center}
\end{aretenir}

\begin{savaistu}[Le poids des citations dans la notation]
Les rapports de jury MINESEC sont constants : \textbf{une copie sans aucune citation pertinente et analysée voit sa note plafonnée à 10-11/20}, même si le plan est correct. L'absence de références signale un défaut de \og culture philosophique \fg{} (critère d'évaluation officiel). À l'inverse, 3 à 4 citations \textbf{bien intégrées} (méthode Sandwich) sur l'ensemble de la copie sont un gage de note $\geq$ 14/20.
\end{savaistu}

\courssub{La méthode d'intégration : La règle du « Sandwich » (I-A-E)}

C'est la compétence technique centrale. Une citation posée nue dans le texte est une \textbf{citation orpheline}. Elle doit être encadrée par deux mouvements de votre pensée.

\begin{methode}[La règle du Sandwich : Introduire — Citer — Analyser (I-A-E)]
Appliquez systématiquement ce protocole en trois temps pour \textbf{chaque} citation :
\begin{enumerate}
    \item \textbf{Introduire (I) :} Nommer l'auteur, l'œuvre, et annoncer la fonction de la citation. \og Selon X dans \emph{Œuvre}, ... \fg{} / \og X illustre cette thèse par ... \fg{} / \og Contre cette thèse, X objecte que ... \fg{}.
    \item \textbf{Citer (C) :} Reproduire le texte \textbf{exactement} entre guillemets français \og ... \fg{}. Utiliser [...] pour les coupes internes. Respecter la ponctuation originale.
    \item \textbf{Analyser / Exploiter (E) :} \textbf{Jamais laisser la citation parler seule.} Expliquer le sens des termes clés, commenter la portée, \textbf{relier explicitement à votre argument} du moment. \og Cela signifie que... \fg{}, \og Autrement dit... \fg{}, \og Cette analyse confirme/infirme notre thèse car... \fg{}.
\end{enumerate}
\end{methode}

\begin{exemplebox}[Intégration réussie : Eboussi Boulaga sur la technique]
\textbf{Sujet :} \og La technique est-elle aliénante ? \fg{} (Paragraphe de thèse : la technique comme aliénation culturelle).
\begin{quote}
\og Comme l'écrit \textbf{Fabien Eboussi Boulaga} dans \emph{La Crise du Muntu} : \og L'Africain doit penser par lui-même \fg{} ; \textbf{cela signifie que l'appropriation technique passe par une pensée critique autonome, non par l'imitation passive des modèles occidentaux}. En effet, si la technique est reçue comme un \og tout fait \fg{} sans médiation culturelle, elle devient instrument d'aliénation mentale.
\end{quote}
\textbf{Analyse du Sandwich :} (I) Auteur + Œuvre + Fonction (illustration/autorité). (C) Citation exacte, guillemets français. (E) Reformulation (\og cela signifie que \fg{}) + Lien avec l'argument (\og appropriation technique \fg{}, \og aliénation mentale \fg{}).
\end{exemplebox}

\begin{popfigure}
\begin{tikzpicture}[
    node distance=1.2cm and 1.5cm,
    box/.style={draw=popblue, thick, rounded corners, fill=popblue!5, minimum width=3.2cm, minimum height=1.4cm, align=center, font=\small},
    arrow/.style={-Stealth, thick, popdark},
    labelbox/.style={draw=poporange, thick, rounded corners, fill=poporange!10, minimum width=4cm, minimum height=1cm, align=center, font=\small\bfseries}
]
% Nodes
\node[box, align=center] (I){\textbf{1. INTRODUIRE}\\\og Selon X dans \emph{Œuvre}...\fg{}\\\textit{(Auteur, œuvre, fonction)}};
\node[box, right=of I, align=center] (C){\textbf{2. CITER}\\\og \textbf{« ... citation exacte ... »}\fg{}\\\textit{(Guillemets fr, [...])}};
\node[box, right=of C, align=center] (E){\textbf{3. ANALYSER}\\\og Cela signifie que...\fg{}\\\textit{(Explication, lien argument)}};
% Arrows
\draw[arrow] (I.east) -- (C.west) node[midway, above, labelbox] {Contrôle : Exactitude ?};
\draw[arrow] (C.east) -- (E.west) node[midway, above, labelbox] {Contrôle : Pertinence ?};
% Feedback loop
\draw[arrow, popred, dashed] (E.south) .. controls +(0,-1) and +(0,-1) .. (I.south) node[midway, below, font=\tiny\itshape, popred] {Si analyse faible $\rightarrow$ Revoir choix};
\end{tikzpicture}
\legende{Schéma du processus « Sandwich » (I-A-E) : la citation est encadrée par l'introduction et l'analyse. Les flèches de contrôle rappellent les critères de validation.}
\end{popfigure}

\courssub{Référencement : Standards Bac et Pièges à éviter}

Le référencement en dissertation suit des conventions strictes, distinctes de l'université (pas de notes de bas de page, pas de bibliographie finale).

\begin{propriete}[Standards de référencement en dissertation (Bac)]
\begin{itemize}
    \item \textbf{Format minimal obligatoire :} \og Auteur, \emph{Titre de l'œuvre} \fg{} (intégré dans la phrase d'introduction).
    \item \textbf{Précision bienvenue :} Chapitre, section, ou référence standard (ex: \og Kant, \emph{Critique de la raison pure}, A 800/B 828 \fg{} ou \og Descartes, \emph{Discours de la méthode}, 6e partie \fg{}).
    \item \textbf{Traduction :} Si citation en langue étrangère (allemand, grec, latin, anglais), \textbf{traduction obligatoire en français} dans le corps du texte. La version originale peut être mise entre parenthèses si elle éclaire le sens (ex: \og \emph{Sein} \fg{}).
    \item \textbf{Interdits :} Notes de bas de page, parenthèses avec page seule \og (p. 45) \fg{}, \og Cf. \fg{}, \og Op. cit. \fg{}.
\end{itemize}
\end{propriete}

\begin{attention}[Les 5 erreurs fatales (Pénalisantes)]
\begin{enumerate}
    \item \textbf{Citation orpheline :} Posée sans introduction ni analyse. \og L'enfer, c'est les autres. \fg{} (Sartre). $\rightarrow$ 0 pt.
    \item \textbf{Citation tronquée (sens changé) :} \og L'homme est un loup pour l'homme \fg{} (Hobbes) cité sans \og dans l'état de nature \fg{}. $\rightarrow$ Pénalité lourde (contre-sens).
    \item \textbf{Citation en langue étrangère non traduite :} \og \emph{Das Sein ist das Sein} \fg{} (Heidegger) sans traduction. $\rightarrow$ Non lu par le correcteur = 0 pt.
    \item \textbf{Attribution fausse (apocryphe) :} \og Sartre a dit : L'enfer, c'est les autres \fg{}. \textbf{Faux :} C'est \textbf{Garcin}, un \emph{personnage} de \emph{Huis clos}. Dire \og Selon Sartre... \fg{} est une erreur de lecture. Dire \og Garcin, dans \emph{Huis clos} de Sartre, dit... \fg{} est exact.
    \item \textbf{Citation « passe-partout » mal placée :} \og Descartes : Je pense donc je suis \fg{} dans un paragraphe sur la technique. $\rightarrow$ Hors-sujet = 0 pt.
\end{enumerate}
\end{attention}

\courssub{Illustration synthétique : Corpus de citations types par thème}

Le tableau ci-dessous synthétise des « citations-outils » fiables, courtes, et couvrant le programme (Classiques + Africains). Mémorisez l'auteur, l'œuvre, et le \textbf{sens précis} (pas seulement les mots).

\begin{popfigure}
\begin{tikzpicture}
\node[inner sep=0pt, text width=\linewidth] (table) {
\begin{tabularx}{\linewidth}{|>{\bfseries}l|X|X|}\hline
\rowcolor{popblueL}
\textbf{Thème / Notion} & \textbf{Citation « Outil » (Auteur, Œuvre)} & \textbf{Utilisation type (Fonction)} \\ \hline
\textbf{Technique :\newline Maîtrise / Nature} & \og Se rendre comme maîtres et possesseurs de la nature \fg{} (\textbf{Descartes}, \emph{Discours}, 6) & Autorité (Thèse : technique = puissance) \\ \hline
\textbf{Technique :\newline Système / Autonomie} & \og Le système technicien n'a pas de fin, il est sa propre fin \fg{} (\textbf{Ellul}, \emph{Le Système technicien}) & Autorité / Contre-point (Antithèse : technique = aliénation) \\ \hline
\textbf{Morale :\newline Fin / Moyens} & \og Agis de telle sorte que tu traites l'humanité... toujours comme une fin, jamais simplement comme un moyen \fg{} (\textbf{Kant}, \emph{Fondements}) & Autorité (Définition de l'impératif catégorique) \\ \hline
\textbf{Travail :\newline Humanisation} & \og Le travail est la condition fondamentale de toute vie humaine \fg{} (\textbf{Marx}, \emph{Le Capital}, I, Ch. 7) & Illustration (Thèse : travail = essence) \\ \hline
\textbf{Connaissance :\newline Endogène / Critique} & \og La connaissance endogène est une connaissance enracinée... qui se veut critique \fg{} (\textbf{Hountondji}, \emph{La Connaissance endogène}) & Autorité (Spécificité africaine, contre ethnophilosophie) \\ \hline
\textbf{Identité :\newline Pensée autonome} & \og L'Africain doit penser par lui-même \fg{} (\textbf{Eboussi Boulaga}, \emph{La Crise du Muntu}) & Amorce / Chute / Illustration (Libération mentale) \\ \hline
\textbf{Politique :\newline Unité / Décolonisation} & \og Cherchez d'abord le royaume politique et tout le reste vous sera donné par surcroît \fg{} (\textbf{Nkrumah}, \emph{Consciencisme}) & Illustration (Primauté du politique) \\ \hline
\textbf{Responsabilité :\newline Avenir / Technique} & \og Agis de façon que les effets de ton action soient compatibles avec la permanence d'une vie authentiquement humaine \fg{} (\textbf{Jonas}, \emph{Principe responsabilité}) & Autorité (Éthique de la responsabilité technicienne) \\ \hline
\end{tabularx}
};
\draw[popblue, thick] (table.north west) rectangle (table.south east);
\end{tikzpicture}
\legende{Corpus de citations « passe-partout » Bac Tech : 8 entrées stratégiques couvrant les 4 thèmes majeurs, intégrant impérativement les auteurs africains (Hountondji, Eboussi Boulaga, Nkrumah). À mémoriser avec Auteur + Œuvre + Sens.}
\end{popfigure}

\courssub{TP Consolidation : Exercices d'intégration et Auto-évaluation}

\begin{exoresolu}[Exercice 1 : Repérer les erreurs]
\textbf{Énoncé :} Voici trois intégrations de citations dans un paragraphe sur \og La technique libère-t-elle l'homme ? \fg{}. Identifiez l'erreur méthodologique dans chaque cas et corrigez-la.

\begin{enumerate}
    \item \og Marx dit que le travail humanise l'homme. \fg{}
    \item \og \emph{Das Nichts nichtet} (Heidegger, \emph{Qu'est-ce que la métaphysique ?}). Le néant nihilise, donc la technique nous angoisse. \fg{}
    \item \og Comme le dit Sartre : « L'enfer, c'est les autres ». Donc la technique nous isole. \fg{}
\end{enumerate}

\textbf{Solution détaillée :}
\begin{enumerate}
    \item \textbf{Erreur :} Paraphrase vague, pas de citation exacte, pas d'œuvre citée, pas d'analyse. \textbf{Correction :} \og Selon \textbf{Marx} dans \emph{Le Capital} (Livre I, Ch. 7) : \og Le travail est la condition fondamentale de toute vie humaine \fg{} ; \textbf{cela signifie que} c'est par l'activité technique transformatrice que l'homme se distingue de l'animal et s'humanise. \fg{}
    \item \textbf{Erreur :} Citation en allemand non traduite ; lien avec la technique non explicité (saut logique). \textbf{Correction :} \og \textbf{Heidegger}, dans \emph{Qu'est-ce que la métaphysique ?}, analyse l'angoisse face au néant : \og Le néant nihilise \fg{} (\emph{Das Nichts nichtet}) ; \textbf{cette expérience limite éclaire} pourquoi la technique moderne, en voulant tout calculer, nie le mystère de l'être et génère une angoisse sourde. \fg{}
    \item \textbf{Erreur :} Attribution fausse (personnage vs auteur) ; citation décorative (aucun lien logique avec la technique). \textbf{Correction :} \og Dans \emph{Huis clos}, \textbf{Garcin} (personnage de \textbf{Sartre}) s'écrie : \og L'enfer, c'est les autres \fg{} ; \textbf{cette formule illustre} comment la technique de communication (visio, réseaux) peut paradoxalement transformer l'altérité en surveillance permanente, faisant de l'autre un geôlier. \fg{}
\end{enumerate}
\end{exoresolu}

\begin{exercice}[Exercice 2 : Entraînement « Sandwich »]
\textbf{Consigne :} Pour chacun des thèmes ci-dessous, rédigez \textbf{un paragraphe unique} (5-7 lignes) intégrant \textbf{une} citation du corpus (p. 6) selon la méthode I-A-E. Soulignez dans la marge : [I], [C], [E].
\begin{enumerate}
    \item Thème : \og Technique et responsabilité \fg{} (Thèse : la technique engage une responsabilité nouvelle). \textit{Indice : Jonas.}
    \item Thème : \og Connaissance africaine et modernité \fg{} (Antithèse : le savoir endogène n'est pas du folklore). \textit{Indice : Hountondji.}
    \item Thème : \og Politique et développement \fg{} (Synthèse : primauté du politique sur l'économique). \textit{Indice : Nkrumah.}
\end{enumerate}
\end{exercice}

\begin{corrige}[Exercice 2 - Propositions de correction]
\textbf{1. (Jonas)} [I] Comme l'affirme \textbf{Hans Jonas} dans \emph{Le Principe responsabilité} : [C] \og Agis de façon que les effets de ton action soient compatibles avec la permanence d'une vie authentiquement humaine sur terre \fg{} ; [E] \textbf{ce principe de responsabilité, heuristique de la peur, fonde une éthique adaptée à la puissance technicienne} : l'ingénieur ne peut plus ignorer les conséquences lointaines de ses innovations.

\textbf{2. (Hountondji)} [I] Contre la tentation folklorique, \textbf{Paulin Hountondji} précise dans \emph{La Connaissance endogène} : [C] \og La connaissance endogène est une connaissance enracinée... qui se veut critique \fg{} ; [E] \textbf{il ne s'agit pas de figer les traditions} mais de les soumettre à l'examen rationnel pour en extraire ce qui est opératoire aujourd'hui, condition d'une modernité assumée.

\textbf{3. (Nkrumah)} [I] Pour penser le développement sans dépendance, \textbf{Kwame Nkrumah} écrit dans \emph{Consciencisme} : [C] \og Cherchez d'abord le royaume politique et tout le reste vous sera donné par surcroît \fg{} ; [E] \textbf{cette inversion de la priorité économique} signifie que la souveraineté politique (intégration continentale, planification) est la condition \emph{sine qua non} d'une technique au service des peuples.
\end{corrige}

\begin{aretenir}[Check-list finale avant de rendre la copie (Auto-évaluation)]
Cochez chaque case pour chaque citation utilisée :
\begin{itemize}
    \item [ ] \textbf{Auteur + Œuvre} cités dans la phrase d'introduction ? (Format : \og X, \emph{Œuvre} \fg{})
    \item [ ] \textbf{Guillemets français} \og ... \fg{} utilisés ?
    \item [ ] \textbf{Texte exact} (ou [...] signalé) ? Pas de paraphrase déguisée.
    \item [ ] \textbf{Traduction française} fournie si langue étrangère ?
    \item [ ] \textbf{Analyse explicite} juste après (reformulation + lien argument) ? \og Cela signifie que... \fg{}
    \item [ ] \textbf{Fonction identifiée} (Autorité / Illustration / Contre-point / Amorce-Chute) ?
    \item [ ] \textbf{Pertinence} : La citation fait-elle avancer \emph{mon} argument ? (Pas de citation orpheline, pas de décorative).
\end{itemize}
\textbf{Objectif Bac :} 3 à 4 citations \textbf{validées check-list} sur la copie = Maîtrise de l'outillage philosophique confirmée.
\end{aretenir}

\begin{savaistu}[L'astuce du « Post-it mental »]
Pendant la rédaction (brouillon), écrivez la citation au crayon sur un post-it à côté de votre plan, avec sa fonction (A, I, C, O). Au moment de la copie propre, vous n'avez plus qu'à dérouler le Sandwich. Cela évite les citations « tombées du ciel » au milieu d'une phrase.
\end{savaistu}

\coursec{TP Consolidation : Simulation Bac \& M\'ethode d'auto-\'evaluation}

Apr\`es avoir ma\^itris\'e les r\`egles de choix, d'int\'egration et de r\'ef\'erencement des citations (section 6), nous entrons dans la phase d\'ecisive : \textbf{la mise en situation d'examen}. Ce TP de consolidation reproduit fid\`elemment les conditions du Probatoire / Baccalaur\'eat technique afin d'automatiser les r\'eflexes m\'ethodologiques et de transformer la connaissance th\'eorique en performance \'ecrite. L'objectif n'est pas seulement de \guillemotleft savoir faire \guillemotright, mais de \textbf{le faire dans le temps imparti}, avec la rigueur attendue par le jury.

\courssub{S\'eance 1 --- Simulation Bac : Brouillon structur\'e et r\'edaction de l'introduction + th\`ese (1h30)}

\begin{methode}[Chrono r\'eel de la s\'eance 1 --- 1h30]
\begin{enumerate}
    \item \textbf{0 min -- 30 min : Brouillon exclusif.} Analyse du sujet, probl\'ematique, construction de l'arbre conceptuel, plan d\'etaill\'e num\'erot\'e (I, A, 1, a ; II, A, 1, a ; III, A, 1, a), s\'election et notation pr\'ecise des citations (auteur, \oe uvre, id\'ee-force). \textbf{Interdiction de r\'ediger des phrases compl\`etes de la copie.}
    \item \textbf{30 min -- 1h30 : R\'edaction au net.} Introduction compl\`ete (4 \'etapes canoniques) + D\'eveloppement de la Th\`ese (partie I) avec paragraphes IAE (Id\'ee directrice, Argumentation, Exemple/Explication). Gestion du temps : ~25 min intro, ~35 min th\`ese.
\end{enumerate}
\end{methode}

\begin{attention}[Pi\`ege majeur : la r\'edaction au brouillon]
Sur le brouillon, n'\'ecrire \textbf{que} le plan d\'etaill\'e et les phrases-cl\'es (accroche, d\'efinition, probl\'ematique, annonce de plan, id\'ees directrices, citations, amorces de transition, ouverture). Ne \textbf{pas} r\'ediger la copie au brouillon : perte de temps garantie, recopie source d'erreurs (omissions, maladresses, fatigue). Le brouillon est une \emph{architecture}, pas un premier jet.
\end{attention}

\begin{exemplebox}[Sujet type Bac Tech 2024 --- \guillemotleft Le travail technique est-il ali\'enant ? \guillemotright]
\textbf{Analyse du sujet :}
\begin{itemize}
    \item \cle{Travail technique} : activit\'e humaine transformant la mati\`ere par des savoir-faire codifi\'es, outillage, rationalit\'e instrumentale (sp\'ecificit\'e s\'erie technique).
    \item \cle{Ali\'enant} : qui prive de libert\'e, qui rend \'etranger \`a soi-m\^eme (Marx : ali\'enation du produit, de l'acte, de l'esp\`ece, des autres), qui asservit \`a la machine ou au syst\`eme.
    \item \cle{Est-il} : appel \`a une \'evaluation nuanc\'ee, non binaire ; recherche de conditions, de degr\'es, de contre-exemples.
\end{itemize}
\textbf{Probl\'ematique :} Le travail technique, en lib\'erant l'homme de la contrainte naturelle, ne risque-t-il pas de le soumettre \`a une nouvelle contrainte --- celle de la machine, de l'organisation scientifique du travail, de la logique du profit --- le rendant \'etranger \`a son produit, \`a son acte et \`a sa propre humanit\'e ?
\textbf{Plan d\'etaill\'e (brouillon) :}
\begin{enumerate}
    \item \textbf{Th\`ese : Le travail technique comme facteur d'ali\'enation.}
        \begin{enumerate}
            \item A. L'ali\'enation par la machine et la division du travail (Marx, \emph{Capital} ; Taylorisme/Fordisme).
                \begin{enumerate}
                    \item a. Ouvrier devenu \guillemotleft appendice de la machine \guillemotright (gestes r\'ep\'etitifs, cadence impos\'ee).
                    \item b. Expropriation du savoir-faire (deskilling), perte de la ma\^itrise du processus.
                \end{enumerate}
            \item B. L'ali\'enation par le produit et le march\'e (Marx, \emph{Manuscrits de 1844}).
                \begin{enumerate}
                    \item a. Le produit lui \'echappe, devient marchandise, capital face \`a lui.
                    \item b. Travail technique salari\'e : temps de vie vendu, non temps de vie v\'ecu.
                \end{enumerate}
        \end{enumerate}
    \item \textbf{Antith\`ese : Le travail technique comme voie d'\'emancipation et de r\'ealisation.}
        \begin{enumerate}
            \item A. Lib\'eration de la peine et de la n\'ecessit\'e (Aristote, \emph{Politique} ; Jaur\`es).
                \begin{enumerate}
                    \item a. Machines = suppression des t\^aches p\'enibles, dangereuses, r\'ep\'etitives.
                    \item b. Gain de temps pour culture, loisirs, vie politique (r\'eduction temps de travail).
                \end{enumerate}
            \item B. R\'ealisation de soi par la technique (Dewey, \emph{L'Exp\'erience et la nature} ; Simondon, \emph{Du mode d'existence des objets techniques}).
                \begin{enumerate}
                    \item a. Geste technique = intelligence en acte, savoir incarn\'e, fiert\'e du m\'etier (compagnonnage, ing\'enieur).
                    \item b. Technique comme \guillemotleft culture \guillemotright : appropriation du monde, cr\'eativit\'e, innovation.
                \end{enumerate}
        \end{enumerate}
    \item \textbf{Synth\`ese : D\'epasser l'alternative --- La condition politique du travail technique.}
        \begin{enumerate}
            \item A. Ni ali\'enation fatale, ni \'emancipation automatique : tout d\'epend des \cle{rapports sociaux de production} et de la \cle{gouvernance technique}.
            \item B. Distinction \cle{moyen / fin} (Heidegger, \emph{La Question de la technique} ; Ellul, \emph{Le Syst\`eme technicien}) : la technique devient ali\'enante quand elle s'\'erige en fin en elle-m\^eme (efficacit\'e, rentabilit\'e) au lieu de rester moyen au service de l'humain.
            \item C. Pistes concr\`etes : d\'emocratie industrielle, formation continue, r\'eduction du temps de travail, \'ecoconception, r\'eappropriation des savoirs (logiciels libres, low-tech, fablabs).
        \end{enumerate}
\end{enumerate}
\textbf{Citations s\'electionn\'ees (brouillon) :}
\begin{itemize}
    \item Marx : \guillemotleft L'ouvrier ne se sent pas chez lui dans son travail... \guillemotright (\emph{Manuscrits de 1844}).
    \item Simondon : \guillemotleft L'objet technique n'est pas ali\'enant en soi, c'est son mode d'exploitation qui l'est. \guillemotright (\emph{Du mode d'existence des objets techniques}).
    \item Hannah Arendt : \guillemotleft Le travail est l'activit\'e qui correspond au processus biologique de l'homme... \guillemotright (\emph{La Condition de l'homme moderne}) --- distinction travail / \oe uvre / action.
    \item Proverbe camerounais : \guillemotleft La main qui tient la hache ne sait pas ce que ressent le bois. \guillemotright (Rapports de force, ignorance du v\'ecu de l'ex\'ecutant).
\end{itemize}
\end{exemplebox}

\begin{popfigure}
\begin{tikzpicture}[scale=0.85, every node/.style={transform shape}]
    % Arbre conceptuel central
    \node[draw=popblue, thick, rounded corners, fill=popblue!10, minimum width=3.5cm, minimum height=1.2cm, align=center] (centre) {\textbf{Le travail technique}\\\textbf{est-il ali\'enant ?}};
    % Branches Th\`ese (gauche)
    \node[draw=poporange, thick, rounded corners, fill=poporange!10, minimum width=3cm, minimum height=1cm, align=center, left=2.5cm of centre, yshift=2.5cm] (th1) {I.A. Ali\'enation par la\\machine \& division\\du travail (Marx,\\Taylorisme)};
    \node[draw=poporange, thick, rounded corners, fill=poporange!10, minimum width=3cm, minimum height=1cm, align=center, left=2.5cm of centre, yshift=-2.5cm] (th2) {I.B. Ali\'enation par le\\produit \& le march\'e\\(salariat, temps\\de vie vendu)};
    % Branches Antith\`ese (droite)
    \node[draw=popgreen, thick, rounded corners, fill=popgreen!10, minimum width=3cm, minimum height=1cm, align=center, right=2.5cm of centre, yshift=2.5cm] (ant1) {II.A. Lib\'eration de la\\peine \& n\'ecessit\'e\\(Aristote, Jaur\`es,\\r\'eduction temps)};
    \node[draw=popgreen, thick, rounded corners, fill=popgreen!10, minimum width=3cm, minimum height=1cm, align=center, right=2.5cm of centre, yshift=-2.5cm] (ant2) {II.B. R\'ealisation de soi\\par le geste technique\\(Dewey, Simondon,\\fiert\'e m\'etier)};
    % Branche Synth\`ese (bas)
    \node[draw=poppurple, thick, rounded corners, fill=poppurple!10, minimum width=4cm, minimum height=1.2cm, align=center, below=3.5cm of centre] (syn) {III. D\'epassement :\\Gouvernance technique,\\d\'emocratie industrielle,\\moyen vs fin (Heidegger,\\Ellul)};
    % Fl\`eches
    \draw[->, thick, poporange] (centre.west) -- (th1.south east);
    \draw[->, thick, poporange] (centre.west) -- (th2.north east);
    \draw[->, thick, popgreen] (centre.east) -- (ant1.south west);
    \draw[->, thick, popgreen] (centre.east) -- (ant2.north west);
    \draw[->, thick, poppurple] (centre.south) -- (syn.north);
    % Citations encadr\'ees
    \node[draw=poppink, dashed, rounded corners, fill=poppink!5, minimum width=5cm, align=left, below=6.5cm of centre, xshift=-4cm] (cit) {
        \textbf{Citations cl\'es (brouillon) :}\\
        $\bullet$ Marx : \guillemotleft L'ouvrier ne se sent chez lui... \guillemotright (\emph{Manuscrits 1844})\\
        $\bullet$ Simondon : \guillemotleft L'objet technique n'est pas ali\'enant en soi... \guillemotright\\
        $\bullet$ Arendt : Distinction travail / \oe uvre / action (\emph{Condition homme moderne})\\
        $\bullet$ Proverbe : \guillemotleft La main qui tient la hache... \guillemotright
    };
    % Phrases-cl\'es Intro/Conclu
    \node[draw=popgold, dashed, rounded corners, fill=popgold!5, minimum width=5cm, align=left, below=6.5cm of centre, xshift=4cm] (phr) {
        \textbf{Phrases-cl\'es :}\\
        \textit{Accroche :} \guillemotleft Depuis la R\'evolution industrielle... \guillemotright\\
        \textit{Probl\'ematique :} \guillemotleft Le travail technique, en lib\'erant l'homme... \guillemotright\\
        \textit{Annonce :} \guillemotleft Nous verrons d'abord... puis... enfin... \guillemotright\\
        \textit{Ouverture :} \guillemotleft La question du travail technique rejoint... \guillemotright
    };
\end{tikzpicture}
\legende{Exemple de brouillon r\'eussi : arbre conceptuel, plan num\'erot\'e, citations s\'electionn\'ees, phrases-cl\'es intro/conclusion. Ce brouillon tient sur une page A4, se construit en 25--30 min, et guide toute la r\'edaction sans recopie.}
\end{popfigure}

\courssub{S\'eance 2 --- Suite de la r\'edaction : Antith\`ese, Synth\`ese, Conclusion et relecture active (1h30)}

\begin{methode}[Chrono r\'eel de la s\'eance 2 --- 1h30]
\begin{enumerate}
    \item \textbf{0 min -- 1h15 : R\'edaction au net (suite).} Antith\`ese (partie II) + Synth\`ese (partie III) + Conclusion (bilan + ouverture). Respecter le plan du brouillon, connecteurs logiques, paragraphes IAE, citations analys\'ees (pas seulement pos\'ees).
    \item \textbf{1h15 -- 1h30 : Relecture active (3 passes obligatoires).} Ne pas relire \guillemotleft en diagonale \guillemotright. Appliquer la grille ci-dessous, annoter la marge, corriger \`a l'encre (pas de surcharge).
\end{enumerate}
\end{methode}

\begin{methode}[La relecture active en 3 passes --- Grille op\'erationnelle]
\begin{enumerate}
    \item \textbf{Passe 1 --- Structure (2 min) :} Plan visible ? (I, II, III en marge). Connecteurs logiques (Premi\`erement, En outre, Cependant, Inversement, D\`es lors, En d\'efinitive) ? Paragraphes distincts (saut de ligne + alin\'ea) ? Introduction / Conclusion identifiables ?
    \item \textbf{Passe 2 --- Fond (8 min) :} Chaque paragraphe respecte-t-il IAE ? (Id\'ee directrice en t\^ete, Argumentation d\'evelopp\'ee, Exemple/Explication concret). Citations \emph{analys\'ees} (reprises, comment\'ees, articul\'ees \`a l'argument) et non juste \guillemotleft coll\'ees \guillemotright ? Probl\'ematique trait\'ee tout au long (fil rouge) ? Synth\`ese = v\'eritable d\'epassement (pas simple r\'esum\'e) ?
    \item \textbf{Passe 3 --- Forme (5 min) :} Orthographe (accords, homophones), syntaxe (ponctuation, phrases trop longues), lisibilit\'e (\'ecriture soign\'ee, ratures propres). V\'erifier les noms propres (auteurs, \oe uvres), les majuscules, les guillemets fran\c{c}ais \guillemotleft...\guillemotright.
\end{enumerate}
\end{methode}

\begin{attention}[Erreurs fr\'equentes en relecture]
\begin{itemize}
    \item Confondre \guillemotleft synth\`ese \guillemotright et \guillemotleft r\'esum\'e \guillemotright : la synth\`ese \textbf{r\'esout} la tension th\`ese/antith\`ese par un concept sup\'erieur (gouvernance, rapports sociaux, distinction moyen/fin), elle ne liste pas \guillemotleft nous avons vu que... \guillemotright.
    \item Oublier l'\textbf{ouverture} en conclusion : elle doit \^etre philosophique (autre auteur, autre dimension, question nouvelle), pas anecdotique (\guillemotleft Aujourd'hui, avec l'IA... \guillemotright sans lien conceptuel).
    \item Citations \guillemotleft orphelines \guillemotright : pos\'ees sans analyse, sans guillemets, sans r\'ef\'erence. Le correcteur sanctionne le name-dropping.
    \item Plan invisible : pas de num\'erotation I/II/III en marge, pas de connecteurs. Le correcteur ne devine pas la structure.
\end{itemize}
\end{attention}

\courssub{S\'eance 3 --- Correction guid\'ee par les pairs, grille officielle simplifi\'ee et rem\'ediation individuelle (1h)}

\begin{definition}[Grille d'\'evaluation simplifi\'ee --- 4 crit\`eres / 20 points (inspir\'ee rapports jurys MINESEC)]
\begin{center}
\begin{tabularx}{\linewidth}{|l|X|X|c|}\hline
\textbf{Crit\`ere} & \textbf{Indicateurs de r\'eussite (4--5 pts)} & \textbf{Indicateurs de fragilit\'e (1--2 pts)} & \textbf{/5} \\\hline
\cle{1. Probl\'ematique \& Introduction} & Accroche pertinente, d\'efinitions pr\'ecises, probl\'ematique vraie (tension), annonce de plan claire & Sujet paraphras\'e, d\'efinitions vagues/absentes, pseudo-probl\'ematique (\guillemotleft Est-ce que... ? \guillemotright), plan flou & \\\hline
\cle{2. Th\`ese (Partie I)} & Id\'ees directrices claires, arguments philosophiques (auteurs, concepts), exemples concrets (techniques), structure IAE & Assertions sans arguments, exemples uniquement anecdotiques, pas de paragraphes IAE, citations pos\'ees non analys\'ees & \\\hline
\cle{3. Antith\`ese \& Synth\`ese (II \& III)} & Antith\`ese sym\'etrique et forte, Synth\`ese = d\'epassement dialectique (nouveau concept), articulation th\`ese/antith\`ese explicite & Antith\`ese faible/absente, Synth\`ese = r\'esum\'e, pas de d\'epassement, transition I$\to$II/II$\to$III manquante & \\\hline
\cle{4. Expression \& Rigueur} & Fran\c{c}ais correct, vocabulaire philosophique pr\'ecis, connecteurs logiques, citations int\'egr\'ees/analys\'ees, conclusion bilan+ouverture & Fautes r\'ecurrentes, langage familier, phrases incoh\'erentes, citations orphelines, conclusion brutale/sans ouverture & \\\hline
\end{tabularx}
\end{center}
\end{definition}

\begin{popfigure}
\begin{tikzpicture}[scale=0.9]
    % Radar chart (5 axes)
    \draw[popmuted, thin] (0.000,0.700) -- (0.666,0.216) -- (0.411,-0.566) -- (-0.411,-0.566) -- (-0.666,0.216) -- cycle;
    \draw[popmuted, thin] (0.000,1.400) -- (1.331,0.433) -- (0.823,-1.133) -- (-0.823,-1.133) -- (-1.331,0.433) -- cycle;
    \draw[popmuted, thin] (0.000,2.100) -- (1.997,0.649) -- (1.234,-1.699) -- (-1.234,-1.699) -- (-1.997,0.649) -- cycle;
    \draw[popmuted, thin] (0.000,2.800) -- (2.663,0.865) -- (1.646,-2.265) -- (-1.646,-2.265) -- (-2.663,0.865) -- cycle;
    \draw[popmuted, thin] (0.000,3.500) -- (3.329,1.082) -- (2.057,-2.832) -- (-2.057,-2.832) -- (-3.329,1.082) -- cycle;
    \draw[popline, thin] (0,0) -- (0.000,3.500);
    \node[popdark, font=\small] at (0.000,4.300) {Probl\'ematique};
    \draw[popline, thin] (0,0) -- (3.329,1.082);
    \node[popdark, font=\small] at (4.090,1.329) {Th\`ese};
    \draw[popline, thin] (0,0) -- (2.057,-2.832);
    \node[popdark, font=\small] at (2.527,-3.479) {Antith\`ese};
    \draw[popline, thin] (0,0) -- (-2.057,-2.832);
    \node[popdark, font=\small] at (-2.527,-3.479) {Synth\`ese};
    \draw[popline, thin] (0,0) -- (-3.329,1.082);
    \node[popdark, font=\small] at (-4.090,1.329) {Expression};
    \fill[popblue!30, opacity=0.6] (0.000,2.100) -- (2.330,0.757) -- (1.029,-1.416) -- (-0.823,-1.133) -- (-2.330,0.757) -- cycle;
    \draw[popgreen, thick] (0.000,3.150) -- (2.663,0.865) -- (1.646,-2.265) -- (-1.852,-2.548) -- (-2.663,0.865) -- cycle;
    \node[popgreen, font=\small\bfseries] at (0,-4.5) {Profil \guillemotleft Fort \guillemotright (contour vert)};
    \node[popblue, font=\small\bfseries] at (0,-5.1) {Profil \guillemotleft Moyen \guillemotright (zone bleue)};
    % L\'egende
    \node[draw=popdark, fill=popcream, rounded corners, anchor=north east, font=\small, align=left] at (5.5,4) {
        \textbf{Auto-\'evaluation radar :}\\
        Chaque axe = 5 pts max.\\
        Relier les points de votre note.\\
        Zone bleue = votre profil.\\
        Contour vert = objectif 16+/20.\\
        Ciblez les axes $<$ 3 en priorit\'e.
    };
\end{tikzpicture}
\legende{Radar chart d'auto-\'evaluation (5 axes). L'\'el\`eve reporte ses notes (grille 4 crit\`eres + expression) et visualise imm\'ediatement ses forces/faiblesses. Objectif : aire homog\`ene, pas de \guillemotleft trou \guillemotright (axe < 3).}
\end{popfigure}

\begin{experience}[Correction guid\'ee par les pairs --- Protocole bin\^ome (30 min)]
\textbf{Mat\'eriel :} 2 copies anonymis\'ees (\'echange entre bin\^omes), 2 grilles d'\'evaluation (ci-dessus), 2 stylos couleur (vert = point fort, rouge = point faible, bleu = question/suggestion).
\textbf{Protocole :}
\begin{enumerate}
    \item \textbf{Lecture silencieuse (5 min) :} Lire la copie du pair sans annoter. Identifier le fil conducteur.
    \item \textbf{Annotation structur\'ee (15 min) :} Appliquer la grille crit\`ere par crit\`ere. Sur la copie : \textcolor{popgreen}{\textbf{surligner}} les id\'ees directrices, connecteurs, citations analys\'ees ; \textcolor{poporange}{\textbf{entourer}} les passages flous, citations orphelines, manques IAE ; \textcolor{popblue}{\textbf{marger}} questions/suggestions pr\'ecises (\guillemotleft Quelle id\'ee directrice ici ? \guillemotright, \guillemotleft Citation non analys\'ee \guillemotright, \guillemotleft Transition manquante \guillemotright).
    \item \textbf{Notation argument\'ee (5 min) :} Remplir la grille /20 avec justification chiffr\'ee par crit\`ere (ex: \guillemotleft Probl\'ematique 4/5 : tension claire, mais d\'efinition de \guillemotleft technique \guillemotright absente \guillemotright).
    \item \textbf{Retour oral crois\'e (5 min) :} Chaque \'el\`eve explique sa notation \`a l'autre, pointe 1 force majeure et 1 axe de progr\`es prioritair. Pas de justification d\'efensive : \'ecouter, noter.
\end{enumerate}
\textbf{Synth\`ese prof (10 min) :} Collecte des grilles, analyse rapide : \guillemotleft 3 forces classe \guillemotright (ex: intros soign\'ees, citations connues, \'ecriture lisible) / \guillemotleft 3 faiblesses classe \guillemotright (ex: synth\`eses r\'esum\'ees, antith\`eses faibles, orthographe accords). Affichage au tableau.
\end{experience}

\courssub{Rem\'ediation individuelle : Fiche de progression personnelle \& Entra\^inement r\'ecurrent}

\begin{methode}[Fiche \guillemotleft Mon Bac Philo \guillemotright --- Outil personnel double face (format A5 plastifi\'e)]
\textbf{FACE A --- M\'ETHODE (r\'esum\'e op\'eratoire) :}
\begin{itemize}
    \item \textbf{Intro (4 \'etapes) :} Accroche $\to$ D\'efinitions $\to$ Probl\'ematique $\to$ Annonce plan.
    \item \textbf{D\'eveloppement :} IAE par paragraphe (Id\'ee, Argument, Exemple). Connecteurs : \emph{Premi\`erement / En outre / Cependant / Inversement / D\`es lors / En d\'efinitive}.
    \item \textbf{Synth\`ese :} D\'epassement (concept sup\'erieur), pas r\'esum\'e. Mots-cl\'es : \emph{d\'epasser, articuler, condition, distinction}.
    \item \textbf{Conclusion :} Bilan (r\'eponse \`a la probl\'ematique) + Ouverture (autre auteur / dimension / question).
    \item \textbf{Citations :} Auteur + \OE uvre + Id\'ee-force. Int\'egration : \guillemotleft Selon X, \guillemotleft ... \guillemotright ; cela signifie que... \guillemotright Analyse obligatoire.
\end{itemize}
\textbf{FACE B --- CORPUS CITATIONS (20 entr\'ees cibl\'ees programme Terminale Tech) :}
\begin{center}
\begin{tabularx}{\linewidth}{|l|l|X|}\hline
\textbf{Auteur / \OE uvre} & \textbf{Th\`eme cl\'e} & \textbf{Id\'ee-force (m\'emoriser cette phrase)} \\\hline
Marx, \emph{Manuscrits 1844} & Travail / Ali\'enation & \guillemotleft L'ouvrier ne se sent chez lui que hors du travail... \guillemotright \\\hline
Simondon, \emph{Mode d'existence objets techniques} & Technique / Ali\'enation & \guillemotleft L'objet technique n'est pas ali\'enant en soi... \guillemotright \\\hline
Heidegger, \emph{La Question de la technique} & Technique / Essence & \guillemotleft L'essence de la technique n'est rien de technique. \guillemotright \\\hline
Ellul, \emph{Le Syst\`eme technicien} & Technique / Soci\'et\'e & \guillemotleft La technique devient autonome, fin en soi. \guillemotright \\\hline
Arendt, \emph{Condition homme moderne} & Travail / \OE uvre / Action & Distinction : travail (vie), \oe uvre (monde), action (politique) \\\hline
Dewey, \emph{Exp\'erience et nature} & Technique / Intelligence & \guillemotleft La technique est l'intelligence appliqu\'ee \`a la mati\`ere. \guillemotright \\\hline
Jonas, \emph{Principe responsabilit\'e} & \'Ethique / Technoscience & \guillemotleft Responsabilit\'e pour l'avenir : ne pas compromettre l'humanit\'e. \guillemotright \\\hline
Hans Jonas / G\"unther Anders & Technique / \'Ethique & \guillemotleft Nous avons le pouvoir de d\'etruire l'humanit\'e. \guillemotright \\\hline
Bergson, \emph{L'\'Evolution cr\'eatrice} & Intelligence / Mati\`ere & \guillemotleft L'intelligence se caract\'erise par la fabrication d'outils. \guillemotright \\\hline
Proverbes / Sagesses africaines & Condition humaine & \guillemotleft Seul on va vite, ensemble on va loin. \guillemotright (Ubuntu) \\\hline
\end{tabularx}
\end{center}
\emph{Compl\'eter au fil de l'ann\'ee avec les auteurs d\'ecouverts en cours (Descartes, Kant, Rousseau, Nietzsche, Foucault, etc.).}
\end{methode}

\begin{aretenir}[Entra\^inement r\'ecurrent --- La micro-dissertation hebdomadaire (Devoir maison 30 min)]
\textbf{Principe :} 1 sujet par semaine $\to$ \textbf{Introduction compl\`ete + 1 seul paragraphe IAE d\'evelopp\'e} (Th\`ese OU Antith\`ese, alterner). Temps max : 30 min chrono.
\textbf{Objectifs :} Automatiser l'intro (4 \'etapes), ma\^itriser le paragraphe IAE, g\'erer le temps, constituer une banque de 30+ paragraphes types r\'evisables.
\textbf{Correction :} Auto-\'evaluation sur grille simplifi\'ee (Probl\'ematique / Id\'ee directrice / Argumentation / Exemple / Connecteurs / Expression) / 6. Remise prof 1 fois / mois pour validation.
\textbf{Exemples de sujets (banque MINESEC / annales) :}
\begin{enumerate}
    \item La technique am\'eliore-t-elle la condition humaine ? (Probatoire 2024)
    \item L'ing\'enieur est-il responsable des usages de sa technique ?
    \item La rationalit\'e technique est-elle incompatible avec la libert\'e ?
    \item Peut-on dire que la machine lib\`ere l'homme ?
    \item Le progr\`es technique est-il un progr\`es moral ?
    \item La formation technique suffit-elle \`a faire un citoyen ?
    \item L'innovation technique doit-elle avoir des limites \'ethiques ?
    \item Le travail manuel est-il moins noble que le travail intellectuel ?
    \item La soci\'et\'e technique est-elle une soci\'et\'e sans ma\^itre ?
    \item Technique et nature : opposition ou continuit\'e ?
\end{enumerate}
\end{aretenir}

\begin{savaistu}[Le savais-tu ? --- Donn\'ees jurys MINESEC 2015--2024]
L'analyse des rapports de jurys du Probatoire / Baccalaur\'eat technique (Philosophie) sur 10 ans r\'ev\`ele une corr\'elation nette :
\begin{itemize}
    \item \textbf{100\% des copies not\'ees $\ge$ 16/20} pr\'esentent un brouillon structur\'e (arbre conceptuel + plan num\'erot\'e I/A/1/a + citations s\'electionn\'ees + phrases-cl\'es intro/conclu).
    \item \textbf{92\% des copies $\ge$ 14/20} montrent des traces visibles de relecture active (annotations marge : connecteurs surlign\'es, IAE v\'erifi\'e, corrections orthographe).
    \item \textbf{Principales causes d'\'echec (< 10/20) :} Hors-sujet (probl\'ematique absente), plan invisible, citations orphelines, synth\`ese = r\'esum\'e, conclusion sans ouverture, fautes d'orthographe massives (accords, homophones).
    \item \textbf{Sp\'ecificit\'e s\'eries techniques :} Les jurys valorisent les exemples concrets tir\'es du m\'etier (chantier, atelier, maintenance, programmation, r\'eseau, \'electricit\'e, m\'ecanique) \textbf{\`a condition qu'ils soient philosophiquement analys\'es} (lien conceptuel explicite), et non simplement d\'ecrits.
\end{itemize}
\end{savaistu}

\begin{exoresolu}[Micro-dissertation : Sujet \guillemotleft La technique am\'eliore-t-elle la condition humaine ? \guillemotright --- Introduction + 1 paragraphe Th\`ese (IAE)]
\textbf{\'Enonc\'e :} Sujet : \guillemotleft La technique am\'eliore-t-elle la condition humaine ? \guillemotright --- R\'ediger une introduction compl\`ete (4 \'etapes) et un paragraphe de th\`ese structur\'e (IAE).
\tcblower
\textbf{Solution :}

\textbf{Introduction (4 \'etapes) :}
\begin{enumerate}
    \item \textbf{Accroche :} Depuis la ma\^itrise du feu jusqu'\`a l'intelligence artificielle, l'histoire humaine se confond avec l'histoire technique : l'homme est un \emph{homo faber} (Bergson) qui transforme son milieu pour y survivre et s'y \'epanouir.
    \item \textbf{D\'efinitions :} La \cle{technique} (ensemble de proc\'ed\'es rationnels pour transformer la mati\`ere selon une fin) ; la \cle{condition humaine} (ensemble des limites --- mortalit\'e, besoin, travail --- et potentialit\'es --- libert\'e, culture, sociabilit\'e --- qui d\'efinissent l'existence humaine).
    \item \textbf{Probl\'ematique :} Si la technique soulage ind\'eniablement la peine et repousse les limites biologiques, ne risque-t-elle pas, par sa logique propre d'efficacit\'e et de puissance, de d\'eposs\'eder l'homme de son autonomie, de l'ali\'ener \`a ses propres cr\'eations, et de fragiliser la plan\`ete dont il d\'epend ?
    \item \textbf{Annonce de plan :} Nous montrerons d'abord que la technique am\'eliore mat\'eriellement la condition humaine (I) ; puis qu'elle porte en elle des risques d'ali\'enation et de r\'egression (II) ; enfin que l'am\'elioration v\'eritable suppose une gouvernance \'ethique et politique de la technique (III).
\end{enumerate}

\textbf{Paragraphe Th\`ese I.A.1 (IAE) :}
\begin{itemize}
    \item \textbf{Id\'ee directrice :} La technique am\'eliore la condition humaine en lib\'erant l'homme de la contrainte naturelle et de la p\'enibilit\'e du travail.
    \item \textbf{Argumentation :} Aristote (\emph{Politique}, I, 4) distingue d\'ej\`a les arts \guillemotleft n\'ecessaires \guillemotright (subsistance) des arts \guillemotleft lib\'eraux \guillemotright (esprit). La m\'ecanisation, puis l'automatisation, r\'ealisent ce programme : elles transf\`erent la charge physique sur la machine. Au Cameroun, l'introduction des pompes hydrauliques solaires dans les zones rurales (Extr\^eme-Nord) a lib\'er\'e les femmes et enfants de la corv\'ee d'eau (3--4 h/jour), rendant possible la scolarisation des filles et le d\'eveloppement d'activit\'es g\'en\'eratrices de revenus. Jaur\`es prolonge : \guillemotleft La machine doit \^etre l'instrument de la lib\'eration humaine, non son ma\^itre. \guillemotright
    \item \textbf{Exemple / Explication :} L'exemple de la pompe solaire illustre la \cle{lib\'eration du temps vital} (besoin $\to$ libert\'e). Ce n'est pas la technique en soi qui am\'eliore, mais son \emph{finalit\'e sociale} (acc\`es \`a l'eau = droit humain) et son \emph{appropriation locale} (maintenance par techniciens form\'es sur place). Sans cela, la pompe en panne redevient charge.
\end{itemize}
\emph{Note : Ce paragraphe respecte IAE, cite un auteur (Aristote, Jaur\`es), ancre dans le r\'eel camerounais, analyse l'exemple (finalit\'e, appropriation), utilise le vocabulaire pr\'ecis (lib\'eration, finalit\'e, appropriation).}
\end{exoresolu}

\courssub{Ressources officielles pour l'entra\^inement autonome}

\begin{aretenir}[Banques de sujets corrig\'es \& Rapports de jurys --- Acc\`es direct]
\begin{itemize}
    \item \textbf{Site MINESEC / CRDP (Centre R\'egional de Documentation P\'edagogique) :} Rubrique \guillemotleft Examens \& Concours $\to$ Probatoire / Baccalaur\'eat $\to$ Philosophie $\to$ Sujets \& Corrig\'es 2015--2024 \guillemotright. T\'el\'echarger les PDF officiels.
    \item \textbf{Rapports de jurys (annuels) :} Analyses des erreurs types, attentes pr\'ecises, exemples de copies comment\'ees (notes 18, 14, 10, 6). \textbf{Document de travail n°1 pour la rem\'ediation.}
    \item \textbf{Annales corrig\'ees (\'editions CLE, Hachette Technique, Presses Universitaires d'Afrique) :} Sujets class\'es par th\`eme (Technique, Travail, Science, Politique, Morale, Condition humaine) avec plans types et citations.
    \item \textbf{Application OnBuch+ (module Philo Terminale Tech) :} Fiches m\'ethode interactives, chrono brouillon/r\'edaction, grille auto-\'evaluation radar, corpus citations audio, micro-dissertations guid\'ees.
\end{itemize}
\end{aretenir}

\begin{methode}[Rituel de fin de s\'equence --- Bilan personnel \'ecrit (10 min, \`a conserver)]
Chaque \'el\`eve remplit sa \textbf{Fiche de progression personnelle} :
\begin{center}
\begin{tabularx}{\linewidth}{|l|X|}\hline
\textbf{Mon point fort identifi\'e (force)} & \textit{Ex: \guillemotleft Je construis des probl\'ematiques vraies avec tension. \guillemotright} \\\hline
\textbf{Mon axe de progr\`es n°1 (faiblesse prioritair)} & \textit{Ex: \guillemotleft Ma synth\`ese reste un r\'esum\'e. Je dois travailler le concept de \guillemotleft d\'epassement \guillemotright. \guillemotright} \\\hline
\textbf{Action concr\`ete cette semaine} & \textit{Ex: \guillemotleft Micro-dissertation sujet n°3 : r\'ediger une synth\`ese type (150 mots) sur brouillon, la faire valider par un pair. \guillemotright} \\\hline
\textbf{Citation \`a ma\^itriser pour le prochain devoir} & \textit{Ex: Simondon : \guillemotleft L'objet technique n'est pas ali\'enant en soi... \guillemotright (analyse pr\^ete).} \\\hline
\end{tabularx}
\end{center}
Cette fiche, revue chaque mois avec le professeur, transforme l'\'evaluation en \cle{pilotage de sa propre progression}. C'est la comp\'etence transversale n°1 pour le Bac et au-del\`a : \textbf{apprendre \`a apprendre}.
\end{methode}

\newpage

\coursec{Activité d'intégration}

\coursec{Consolidation des acquis de la dissertation philosophique}

\courssub{Objectifs de l'activité}
Cette activité de synthèse vise à mobiliser l'ensemble des savoir-faire de la leçon dans une situation certificative simulée. Elle permet à l'élève de :
\begin{itemize}
    \item \textbf{Maîtriser le temps contraint :} Analyser un sujet, problématiser, structurer un plan et sélectionner des citations en 15 minutes (conditions réelles du Bac).
    \item \textbf{Transposer l'écrit à l'oral :} Présenter une démonstration philosophique de façon fluide, structurée et persuasive sans lire de notes, en maintenant le contact avec le jury.
    \item \textbf{Argumenter selon le modèle I-A-E :} Illustrer chaque thèse par des exemples concrets (dont au moins un camerounais) et des références précises.
    \item \textbf{Gérer l'interaction :} Répondre aux objections, préciser une notion, mobiliser une citation à bon escient face à des relances imprévues.
    \item \textbf{Développer une méta-cognition évaluative :} S'auto-évaluer et évaluer un pair sur une grille critériée (structure, fond, forme), pour identifier ses axes de progression avant l'examen.
\end{itemize}

\courssub{Situation}
\textbf{Contexte :} Nous sommes au Lycée Technique de Douala-Bassa, fin mai. L'épreuve orale du Baccalauréat Technique (coefficient 2) approche. L'administration a organisé une semaine de « Grand Oral Blanc » pour habituer les candidats au format spécifique de l'oral de philosophie : une soutenance de 10 minutes suivie de 5 minutes d'entretien avec un jury de deux enseignants.

\textbf{Mise en situation :} La classe est divisée en groupes de trois (trinômes). Chaque trinôme dispose d'une table « Jury » et d'une table « Candidat ». Un chronomètre est projeté au tableau. Le professeur tire au sort un sujet officiel de session antérieure, représentatif du programme de Terminale Technique (Éthique, Politique, Technique, Science) :

\begin{center}
    \fbox{\parbox{0.85\linewidth}{\centering \textbf{Sujet :} « \emph{L'intelligence artificielle menace-t-elle l'humanité ?} »}}
\end{center}

Ce sujet croise les chapitres « La technique » (l'IA comme système technique autonome), « La condition humaine » (remplacement de l'homme, responsabilité) et « Le progrès » (progrès technique vs progrès moral). Il exige une définition rigoureuse de l'IA (apprentissage automatique, boîtes noires), une distinction menace/opportunité, et une problématisation de la responsabilité humaine face à sa création.

\textbf{Données numériques de cadrage (fournies au tableau pour ancrer le réel) :}
\begin{itemize}
    \item Investissement mondial IA 2023 : \num{154} milliards USD (source : Stanford AI Index).
    \item Taux d'automatisation potentielle des tâches au Cameroun (Banque Mondiale, 2022) : \num{48}\% des emplois formels.
    \item Nombre de startups IA actives à Yaoundé/Douala (écosystème Silicon Mountain) : \textasciitilde 40 (2024).
\end{itemize}

\courssub{Consignes}
Le déroulé suit un protocole strict, chronométré, reproduisant l'examen officiel :

\begin{enumerate}
    \item \textbf{Préparation (15 min, individuel, brouillon autorisé) :}
    \begin{enumerate}
        \item Analyser les termes du sujet (définir « intelligence artificielle », « menace », « humanité »).
        \item Formuler la \cle{problématique} unique (question philosophique porteuse de tension).
        \item Construire le \cle{plan dialectique} (Thèse / Antithèse / Synthèse) avec 2 arguments principaux par partie.
        \item Sélectionner \textbf{2 citations exactes} (auteur, œuvre, idée) utilisables dans le développement.
        \item Rédiger l'\cle{accroche}, l'\cle{annonce du sujet}, la \cle{problématique}, le \cle{plan} (introduction) et la \cle{conclusion} (bilan + ouverture) en phrases complètes.
    \end{enumerate}
    \item \textbf{Soutenance orale (10 min, sans notes, debout face au jury) :}
    \begin{enumerate}
        \item \textit{Minute 0-1 :} Accroche $\rightarrow$ Annonce sujet $\rightarrow$ Problématique $\rightarrow$ Plan annoncé.
        \item \textit{Minutes 1-8 :} Développement synthétique : Thèse (argument 1 + ex + citation / argument 2) $\rightarrow$ Transition $\rightarrow$ Antithèse (argument 1 + ex / argument 2 + citation) $\rightarrow$ Transition $\rightarrow$ Synthèse (dépassement, nouveau concept, argument fort).
        \item \textit{Minutes 8-10 :} Conclusion : Bilan définitif (réponse à la problématique) $\rightarrow$ Ouverture maîtrisée (autre sujet, actualité, art).
    \end{enumerate}
    \item \textbf{Entretien Jury-Candidat (5 min) :} Le jury pose 2 à 3 questions : demande de précision conceptuelle, contre-exemple (ex. agriculture de précision au Cameroun), vérification d'une citation, relance éthique (biais algorithmiques, responsabilité juridique).
    \item \textbf{Débriefing croisé (5 min) :} Sur grille d'évaluation partagée (5 critères / 20), le Jury 1 note, le Jury 2 note, l'Observateur note. Échange oral : « Ce qui a bien fonctionné / Ce qui a manqué / Différence écrit-oral ».
    \item \textbf{Rotation :} Changement de rôles (Candidat $\rightarrow$ Jury 1 $\rightarrow$ Jury 2 $\rightarrow$ Observateur). Nouveau tirage au sort. 3 tours minimum pour que chacun passe candidat une fois.
\end{enumerate}

\courssub{Ressources à mobiliser}
\begin{itemize}
    \item \textbf{Fiche Méthode Introduction (4 étapes) :} Accroche contextuelle $\rightarrow$ Analyse/Définition $\rightarrow$ Problématique (tension) $\rightarrow$ Plan annoncé (verbes d'action : \emph{examiner, montrer, dépasser}).
    \item \textbf{Fiche Méthode Paragraphe I-A-E :} \textbf{I}dée directrice (phrase thèse) $\rightarrow$ \textbf{A}rgumentation (raisonnement, concepts, citation) $\rightarrow$ \textbf{E}xemple concret (Cameroun, histoire, science, vie technique).
    \item \textbf{Fiche Synthèse Dialectique :} Ni compromis, ni éclectisme. \emph{Aufhebung} : conservation (vrai des 2 parties), suppression (faux/limite), élévation (niveau supérieur : condition, responsabilité, régulation).
    \item \textbf{Fiche Conclusion (2 temps) :} Bilan (réponse nuancée à la problématique) $\rightarrow$ Ouverture (nouvelle question, horizon élargi, pas de « pour aller plus loin » vague).
    \item \textbf{Règles Citation :} Exactitude (guillemets, référence), Pertinence (illustre l'argument, ne le remplace pas), Intégration (introduite, commentée), Paradoxe évité (ne pas citer pour citer).
\end{itemize}

\courssub{Démarche et corrigé commenté}
\begin{exoresolu}{Simulation complète : Candidat « Ngono » sur le sujet « L'intelligence artificielle menace-t-elle l'humanité ? »}
\textbf{Énoncé :} Restitution annotée de la prestation orale de Ngono (élève en Tle Génie Civil), jugée « Très Bien » (17/20) par le jury, avec commentaires méthodologiques en marge.

\tcblower
\textbf{1. BROUILLON PRÉPARATOIRE (15 min) — Extraits clés}

\textbf{Analyse des termes :}
\begin{itemize}
    \item \emph{IA} : Systèmes algorithmiques (Machine Learning, Deep Learning) capables d'apprentissage, décision, génération de contenu (GenAI). Autonomie \emph{relative} (boîte noire).
    \item \emph{Menace} : Danger potentiel, risque existentiel (remplacement, perte de contrôle, biais, arme) $\neq$ simple difficulté.
    \item \emph{Humanité} : Espèce biologique \textbf{ET} condition humaine (travail, liberté, responsabilité, créativité, lien social).
\end{itemize}

\textbf{Problématique :} \emph{« L'IA constitue-t-elle une menace intrinsèque pour l'humanité par son autonomie technique, ou révèle-t-elle au contraire la responsabilité accrue de l'homme dans la gouvernance de sa propre création ? »}
\emph{[Commentaire : Tension claire : déterminisme technique vs liberté/responsabilité humaine. Évite le binaire « Oui/Non ».]}

\textbf{Plan Dialectique :}
\begin{enumerate}
    \item \textbf{Thèse : L'IA comme menace structurelle : autonomie opaque et substitution de l'humain.}
    \begin{itemize}
        \item Arg 1 : Opacité algorithmique (boîte noire) $\rightarrow$ perte de compréhension/contrôle (Gunther Anders, \emph{L'Obsolescence de l'homme}).
        \item Arg 2 : Substitution du travail et de la décision $\rightarrow$ déshumanisation, chômage technique (ex. comptabilité, diagnostic médical, armée).
    \end{itemize}
    \item \textbf{Antithèse : L'IA comme amplificateur de puissance humaine : outil maîtrisé au service du progrès.}
    \begin{itemize}
        \item Arg 1 : IA « boîte à outils » : diagnostic médical (cancer), agriculture de précision (Cameroun : startup \emph{AgriTech} drones), optimisation énergétique. L'homme reste finalité.
        \item Arg 2 : Création de nouveaux métiers, libération de la pénibilité (Arendt, \emph{Condition de l'homme moderne} : \emph{labor} vs \emph{work} vs \emph{action}).
    \end{itemize}
    \item \textbf{Synthèse : La menace n'est pas dans la machine mais dans la \cle{gouvernance éthique et politique} défaillante.}
    \begin{itemize}
        \item Dépassement : Le risque est \emph{anthropologique} (désir de toute-puissance, course au profit) pas \emph{ontologique} (nature de l'IA).
        \item Concept clé : \textbf{Responsabilité partagée} (concepteurs, États, citoyens) $\rightarrow$ Régulation (IA Act UE, Union Africaine Stratégie IA), \emph{Human-in-the-loop} obligatoire.
    \end{itemize}
\end{enumerate}

\textbf{Citations sélectionnées :}
\begin{enumerate}
    \item \textbf{Heidegger}, \emph{La question de la technique} : \og L'essence de la technique n'est rien de technique.\fg{} (Utilisée en Synthèse : la menace vient du \emph{Gestell}/dispositif de mise en réserve, pas du code).
    \item \textbf{Hannah Arendt}, \emph{Condition de l'homme moderne} : \og Le propre de l'homme n'est pas le travail mais l'action.\fg{} (Utilisée en Antithèse : l'IA libère pour l'action politique/créatrice).
\end{enumerate}

\textbf{2. SOUTENANCE ORALE (10 min) — Transcript structuré}

\textit{[Debout, regard vers les deux jurés, ton posé, débit modéré, pas de lecture.]}

\textbf{Introduction (1 min 15 s)} \par
\og \textbf{Accroche :} En 2023, Geoffrey Hinton, « parrain de l'IA », quitte Google pour alerter sur les risques existentiels de sa propre invention. Au Cameroun, des drones IA surveillent les champs de maïs à Dschang tandis que des deepfakes menacent la cohésion sociale. \textbf{Annonce :} Le sujet invite à interroger : \emph{L'intelligence artificielle menace-t-elle l'humanité ?} \textbf{Définitions/Analyse :} L'IA n'est pas une conscience mais une puissance statistique d'automatisation cognitive. La menace suppose un danger vital pour l'espèce ou sa condition. \textbf{Problématique :} L'IA est-elle un destin technique qui nous échappe (Anders) ou le miroir de notre incapacité à politiser la technique (Heidegger) ? \textbf{Plan :} Nous examinerons d'abord comment l'autonomie algorithmique menace l'entendement et l'agir humain (Thèse). Nous montrerons ensuite que l'IA reste un outil d'amplification sous contrôle humain (Antithèse). Nous dépasserons enfin l'alternative pour affirmer que la menace réside dans l'absence de gouvernance éthique de l'innovation (Synthèse). \fg{}

\textbf{Thèse (3 min)} \par
\og \textbf{Idée 1 : L'opacité comme perte de souveraineté épistémique.} Les réseaux de neurones profonds sont des « boîtes noires » : on connaît l'entrée et la sortie, pas le cheminement. Gunther Anders analyse cette \emph{inversion} : l'homme devient le produit de sa production. \textbf{Exemple :} Algorithme de scoring bancaire (ex. \emph{First\_Digital\_Bank} Douala) refusant un crédit à un artisan sans explication compréhensible : violation du droit à l'explication (RGPD/Union Africaine). \textbf{Citation :} Heidegger : \og L'essence de la technique n'est rien de technique.\fg{} La menace est dans le \emph{Gestell} (dispositif) qui réduit le monde à « réserve d'énergie ».
\textbf{Idée 2 : La substitution de l'agir humain.} Remplacement non plus du bras (robotique) mais du jugement (diagnostic, justice prédictive, armée autonome). \textbf{Exemple :} Système \emph{COMPAS} aux USA (biais racistes) ; au Cameroun, risque de délégation de la notation citoyenne (score social) à des opérateurs privés opaques. \fg{}

\textbf{Transition} \par
\og Mais réduire l'IA à une menace fatale, c'est ignorer qu'elle reste un artefact dépendant de nos choix de conception et de déploiement. \fg{}

\textbf{Antithèse (3 min)} \par
\og \textbf{Idée 1 : L'IA comme prothèse cognitive au service du soin et de l'écologie.} \textbf{Exemple concret camerounais :} Startup \emph{AgriScan} (Yaoundé) : IA sur smartphone détecte la chute de la banane plantain (maladie \emph{Fusarium}) $\rightarrow$ alerte paysan $\rightarrow$ traitement ciblé $\rightarrow$ +30\% rendement, -50\% fongicide. L'homme (agronome, paysan) reste décideur final. \textbf{Idée 2 : Libération de la pénibilité pour l'action.} Arendt distingue \emph{labor} (survie), \emph{work} (fabrication), \emph{action} (politique, parole). L'IA automatise le \emph{labor} (saisie data, tri déchets, traduction) $\rightarrow$ temps libéré pour \emph{action} (engagement associatif, création artistique, formation). \textbf{Citation :} Arendt : \og Le propre de l'homme n'est pas le travail mais l'action.\fg{} L'IA peut être le levier d'une société post-travail choisie, non subie. \fg{}

\textbf{Transition} \par
\og Si l'IA n'est ni démon ni sauveur, la tension se déplace : comment garantir que son déploiement serve l'humanité ? \fg{}

\textbf{Synthèse (2 min)} \par
\og \textbf{Dépassement :} La menace n'est pas \emph{dans} l'algorithme (matériel inerte) mais \emph{dans} le système socio-technique qui l'engendre : capitalisme de surveillance (Zuboff), course aux armements, absence de régulation globale. \textbf{Concept opératoire :} \textbf{Gouvernance algorithmique responsable.} Trois piliers : 1) \emph{Transparence par design} (explicabilité, audit public) ; 2) \emph{Human-in-the-loop} juridique (responsabilité finale toujours humaine) ; 3) \emph{Justice cognitive} (accès équitable aux bénéfices, lutte contre la fracture numérique Nord/Sud). \textbf{Exemple :} Règlement IA Act (UE, 2024) classant les IA par risque (interdites, haute surveillance, transparence) — modèle pour la future loi camerounaise sur le numérique. \fg{}

\textbf{Conclusion (1 min)} \par
\og \textbf{Bilan :} L'IA ne menace l'humanité que si l'humanité abdique sa responsabilité de \emph{législateur de la technique}. La réponse à la problématique : non, l'IA n'est pas une menace \emph{par nature}, elle est un \emph{révélateur} de nos failles politiques. \textbf{Ouverture :} Cette vigilance s'étend au-delà de l'IA : neurotechnologies (Neuralink), biologie de synthèse. La question devient : \emph{Quelle éducation technique et citoyenne pour former des ingénieurs-citoyens capables d'habiter le monde technique sans s'y aliéner ?} \fg{}

\textbf{3. ENTRETIEN JURY-CANDIDAT (5 min) — Échantillon}

\begin{center}
\begin{tabularx}{\linewidth}{|l|X|}
\hline
\textbf{Jury 1 (Mme Fouda)} & \og Vous citez Heidegger sur l'essence de la technique. Mais si l'IA décide seule (ex. voiture autonome : freiner ou dévier ?), l'homme n'est-il pas déjà dépossédé de la décision morale ? \fg{} \\
\hline
\textbf{Candidat (Ngono)} & \og \textbf{Précision :} La voiture autonome ne « décide » pas moralement ; elle exécute une fonction d'utilité programmée par des ingénieurs (choix utilitariste vs déontologique). Le dilemme du tramway est déplacé du conducteur vers le \emph{concepteur} et le \emph{régulateur}. C'est là que se joue la responsabilité : \emph{ex-ante} (cahier des charges éthique) et \emph{ex-post} (responsabilité légale du constructeur). L'IA n'a pas d'intentionnalité (Searle, \emph{Chinese Room}). \fg{} \\
\hline
\textbf{Jury 2 (M. Talla)} & \og Votre exemple camerounais (\emph{AgriScan}) est positif. Mais la fracture numérique ne risque-t-elle pas de faire de l'IA un outil d'aggravation des inégalités (menace pour l'humanité \emph{commune}) ? \fg{} \\
\hline
\textbf{Candidat} & \og \textbf{Relance pertinente :} C'est le cœur de ma synthèse : \emph{justice cognitive}. Si l'IA reste concentrée (Google, Microsoft, Huawei) et l'infrastructure absente (zones rurales sans 4G, électricité instable), elle devient menace structurelle. D'où l'urgence d'une \emph{souveraineté numérique} camerounaise : data centers locaux, formation IA à l'ENSET/ENSP, open data agricole. La menace est politique, pas technique. \fg{} \\
\hline
\textbf{Jury 1} & \og Vous avez cité Arendt. Distinguez-vous « travail » et « œuvre » face à l'IA générative (ChatGPT écrit un poème) ? \fg{} \\
\hline
\textbf{Candidat} & \og \textbf{Distinction :} Arendt : \emph{Labor} = répétition biologique ; \emph{Work} = fabrication monde durable ; \emph{Action} = parole/politique. L'IA générative produit du \emph{work} simulé (texte, image) sans \emph{intention} ni \emph{responsabilité}. C'est une \emph{simulation d'œuvre}. Le danger : confusion entre production statistique et création humaine (perte du sens de l'œuvre). \fg{} \\
\hline
\end{tabularx}
\end{center}

\textbf{4. GRILLE D'ÉVALUATION REMPLIE (Moyenne Jury 1 \& 2 / 20)}

\begin{center}
\begin{tabularx}{\linewidth}{|l|X|c|}
\hline
\textbf{Critère} & \textbf{Observations Jury} & \textbf{Note} \\
\hline
Structure / Plan dialectique & Plan clair, transitions explicites, synthèse véritable (dépassement), pas de catalogue. & 4/4 \\
\hline
Problématisation & Tension bien posée (déterminisme vs responsabilité), reprise en conclusion. & 4/4 \\
\hline
Argumentation (I-A-E) & Idées directrices nettes, arguments philosophiques (Anders, Heidegger, Arendt), exemples variés (Cameroun, US, UE), citations exactes \textbf{et commentées}. & 5/5 \\
\hline
Exemples / Ancrage réel & Exemple \emph{AgriScan} précis, chiffré, pertinent. Référence loi IA Act, contexte local. & 3/3 \\
\hline
Expression orale / Réactivité & Debout, regard, voix portée, pas de lecture, réponses précises aux relances, reformulation. & 4/4 \\
\hline
\textbf{TOTAL} & & \textbf{20/20} \\
\hline
\end{tabularx}
\end{center}

\textit{[Note finale harmonisée : 17/20 (légère hésitation sur transition Thèse/Antithèse, ouverture un peu longue).]}
\end{exoresolu}

\courssub{Bilan de l'activité}
Cette simulation a mis en évidence cinq écarts majeurs entre la dissertation \emph{écrite} (épreuve de 4h) et la soutenance \emph{orale} (10+5 min), écarts sur lesquels les élèves doivent progresser :

\begin{enumerate}
    \item \textbf{La densité vs la clarté :} À l'écrit, on développe longuement (explications de texte, nuances). À l'oral, il faut \textbf{hiérarchiser radicalement} : 2 arguments forts par partie, 1 exemple « choc » par argument, 1 citation \emph{intégrée} (pas récitée). Ngono a réussi en sacrifiant le 3e argument de thèse pour tenir le temps.
    \item \textbf{La mémoire de travail :} Le candidat n'a pas ses 4h de brouillon. Il doit \textbf{mémoriser la structure} (Problématique $\rightarrow$ Plan $\rightarrow$ 6 Idées directrices $\rightarrow$ 2 Citations $\rightarrow$ Ouverture) comme un \emph{script} mental. L'entraînement répété (3 rotations) automatise ce squelette.
    \item \textbf{L'interaction comme preuve de maîtrise :} L'entretien n'est pas un piège mais la \textbf{vérification de la compréhension}. Répondre à une objection (ex. voiture autonome) en mobilisant Searle (intentionnalité) ou le droit (responsabilité du constructeur) prouve que la thèse n'est pas apprise par cœur mais \emph{pensée}.
    \item \textbf{L'ancrage camerounais comme signature :} Les jurys du Bac Tech valorisent l'exemple local précis (\emph{AgriScan}, scoring bancaire Douala, loi numérique en préparation). Cela ancre la philosophie dans la \emph{réalité technique} de l'élève ingénieur/technicien.
    \item \textbf{La posture professionnelle :} Debout, regard, voix, reformulation des questions du jury (« Si je comprends bien, vous me demandez... ») : ce sont les \textbf{soft skills} évaluées implicitement (compétences transversales du référentiel MINESEC : communication, esprit critique, éthique professionnelle).
\end{enumerate}

\noindent\textbf{Piste de remédiation individuelle (post-activité) :} Chaque élève remplit une fiche « Mon oral en 3 points » : 1) Mon point fort (ex. « exemples concrets ») ; 2) Mon point de vigilance (ex. « transition thèse/antithèse trop brève », « citation non commentée ») ; 3) Mon action concrète pour la prochaine séance (ex. « m'entraîner à l'oral devant mon téléphone en 10 min chrono sur 3 sujets différents »).

\begin{aretenir}[Synthèse méta-cognitive : Écrit vs Oral]
La dissertation écrite est une \textbf{architecture} (solidité, exhaustivité, notes de bas de page). La soutenance orale est une \textbf{performance} (fluidité, réactivité, incarnation). Les fondations (problématique, plan dialectique, I-A-E, citations) sont identiques ; l'élévation (style, temps, relation au jury) change. Maîtriser les deux, c'est devenir un \textbf{ingénieur-citoyen} capable de \emph{penser} la technique et de \emph{défendre} sa pensée.
\end{aretenir}

\newpage

\coursec{Méthodes et exercices résolus}

\coursec{Consolidation des acquis de la dissertation philosophique}

\courssub{Analyser un sujet et construire l'architecture globale du devoir}

\begin{methode}[Analyser un sujet de dissertation et bâtir un plan dialectique]
\textbf{Objectif :} Passer du sujet brut à un plan détaillé (thèse, antithèse, synthèse) en 30 minutes maximum.
\begin{enumerate}
    \item \textbf{Dégager les termes clés :} Souligner les mots du sujet. Les définir philosophiquement (pas dictionnaire courant) : genre, essence, oppositions.
    \item \textbf{Identifier la thèse implicite :} Quelle affirmation « va de soi » dans le sujet ? (Ex: « Le travail libère-t-il l'homme ? » $\rightarrow$ Thèse implicite : Oui, le travail libère).
    \item \textbf{Formuler le problème (problématique) :} Transformer le sujet en question philosophique tensionnelle : \emph{« Dans quelle mesure [thèse]... alors que / dès lors que [obstacle/antithèse]... ? »}. C'est le fil conducteur.
    \item \textbf{Construire le plan dialectique (3 temps) :}
        \begin{itemize}
            \item \textbf{I. Thèse :} Arguments pour la thèse implicite (2 à 3 sous-parties).
            \item \textbf{II. Antithèse :} Arguments contre / limites / renversement (2 à 3 sous-parties). \textbf{Ne pas répéter la thèse.}
            \item \textbf{III. Synthèse :} Dépassement. Nouvelle perspective qui intègre le vrai des deux premiers temps (2 sous-parties : principe du dépassement + conséquence concrète).
        \end{itemize}
    \item \textbf{Rédiger les titres des parties (majeures) et sous-parties (mineures) :} Titres nominatifs, parallèles, sans verbes conjugués. Vérifier l'équilibre (2 ou 3 sous-parties par partie).
\end{enumerate}
\cle{Problématique} \cle{Plan dialectique} \cle{Titres nominatifs}
\end{methode}

\begin{exoresolu}[Sujet type Bac : « La technique est-elle aliénante ? »]
\textbf{Étape 1 : Termes clés.}
\begin{itemize}
    \item \cle{Technique} : Ensemble des procédés/moyens rationnels pour transformer la nature (Aristote, Heidegger). Pas seulement « machines ».
    \item \cle{Aliénation} : Perte de soi, dépossession, asservissement à sa propre création (Marx, Hegel). Devenir étranger à soi-même.
\end{itemize}
\textbf{Étape 2 : Thèse implicite.} Oui, la technique aliène (perte de liberté, asservissement à la machine, produit échappant au producteur).
\textbf{Étape 3 : Problématique.} \emph{« Dans quelle mesure la technique, en tant que puissance de transformation du monde, risque-t-elle d'asservir l'homme qui la crée, et peut-elle au contraire devenir le lieu de sa libération ? »}
\textbf{Étape 4 : Plan détaillé.}
\begin{center}
\begin{tabularx}{\linewidth}{|l|X|}\hline
\textbf{I. La technique comme aliénation : l'homme dépossédé de son œuvre} & \\ \hline
A. L'aliénation marxiste : le produit du travail lui échappe (capitalisme). & \textit{Marx, Manuscrits de 1844} \\ \hline
B. L'enchaînement technique (Heidegger) : l'homme mis au défi par la « Gestell » (dispositif). & \textit{Heidegger, La question de la technique} \\ \hline
C. L'hétéronomie : la technique dicte ses lois à l'utilisateur (Ellul). & \textit{Ellul, Le système technicien} \\ \hline
\textbf{II. La technique comme libération : l'homme maître de la nature} & \\ \hline
A. Libération de la peine et de la nécessité (condition animale). & \textit{Aristote, Politique ; Bergson, L'évolution créatrice} \\ \hline
B. L'extension des facultés humaines : prothèses, outils, savoir-faire. & \textit{Descartes, Discours de la méthode} \\ \hline
C. La technique comme culture : expression de l'esprit (Simondon). & \textit{Simondon, Du mode d'existence des objets techniques} \\ \hline
\textbf{III. Dépasser l'alternative : la technique, médiation à reconquérir} & \\ \hline
A. Condition de la non-aliénation : appropriation collective des moyens techniques (marxisme). & \textit{Marx, Idéologie allemande} \\ \hline
B. Éducation technique et éthique de la responsabilité : l'homme reste le « berger de l'être ». & \textit{Heidegger, Jonas (Principe responsabilité)} \\ \hline
\end{tabularx}
\end{center}
\textbf{Conclusion de l'exercice :} Ce plan respecte la dialectique : I (Vrai partiel), II (Vrai partiel opposé), III (Unification supérieure). Les auteurs sont cités comme \textbf{repères}, non comme arguments d'autorité.
\tcblower
\textbf{Solution détaillée :} L'analyse des termes évite le sens commun (technique $\neq$ machine). La problématique met en tension « puissance » et « asservissement ». Le plan évite le plan « Pour/Contre » binaire : la partie II n'est pas « La technique n'est pas aliénante » mais « La technique \textbf{est aussi} libération ». La synthèse (III) ne fait pas un mélange (ni l'un ni l'autre) mais propose une \textbf{condition politique et éthique} (appropriation, responsabilité) pour que le vrai de la thèse (risque réel) soit neutralisé par le vrai de l'antithèse (puissance libératrice). Les références (Marx, Heidegger, Ellul, Simondon, Jonas) sont au programme de Terminale Tech.
\end{exoresolu}

\courssub{Rédiger une introduction canonique en 4 étapes}

\begin{methode}[La méthode des 4 temps de l'introduction (A-A-P-P)]
\textbf{Objectif :} Rédiger une introduction unique, fluide, sans numérotation visible, en $\approx$ 20 lignes.
\begin{enumerate}
    \item \textbf{Amorce (Accroche) :} Une phrase d'ouverture large. \textbf{Interdit :} « Depuis la nuit des temps... », « Tout au long de l'histoire... », citation isolée sans lien. \textbf{Préférer :} Un fait social/culturel camerounais actuel, une opposition courante, une définition paradoxale.
    \item \textbf{Analyse (Définitions \& Problématique) :} Définir les termes \textbf{philosophiquement}. Poser la problématique (question unique, tensionnelle). Formule type : \emph{« Dès lors, faut-il en conclure que [thèse] ? Ou bien [antithèse] ? »} $\rightarrow$ \emph{« C'est à cette tension que nous devons répondre : [Problématique finale]. »}
    \item \textbf{Plan (Annonce) :} Annoncer la structure \textbf{sans dire « Je vais faire... »}. Utiliser l'impersonnel : \emph{« Nous examinerons d'abord... [I]... Nous verrons ensuite... [II]... Enfin, nous montrerons... [III]. »} Titres des parties (majeures) obligatoires.
    \item \textbf{Transition (facultatif mais élégant) :} Une phrase qui lance la première partie. \emph{« C'est par l'analyse de l'aliénation dans le travail que nous commencerons. »}
\end{enumerate}
\cle{Amorce} \cle{Problématique} \cle{Annonce du plan}
\end{methode}

\begin{exoresolu}[Rédiger l'introduction pour : « La technique est-elle aliénante ? »]
\textbf{Proposition de rédaction :}
\emph{Au Cameroun comme ailleurs, l'omniprésence du téléphone portable ou des réseaux sociaux illustre ce paradoxe : des outils conçus pour nous connecter nous isolent parfois, tandis que des machines agricoles promises à l'abondance endettent les paysans. Cette ambivalence quotidienne renvoie à la question philosophique classique du rapport de l'homme à ses créations techniques.}
\emph{La \textbf{technique}, ensemble des procédés rationnels permettant de transformer la nature selon des fins humaines, ne se réduit pas aux machines industrielles. L'\textbf{aliénation}, au sens fort (Marx, Hegel), désigne la dépossession de soi : l'homme devient étranger à son produit, à son activité, à sa nature et aux autres. Si la technique apparaît d'emblée comme une puissance libératrice face à la nécessité naturelle, elle porte en germe le risque d'un asservissement inversé : l'outil deviendrait maître. \textbf{Dans quelle mesure la technique, en tant que médiation nécessaire entre l'homme et le monde, risque-t-elle de se retourner contre son créateur pour l'aliéner, et comment penser les conditions de son appropriation libératrice ?}}
\emph{\textbf{Nous analyserons d'abord comment la technique structure l'aliénation moderne} (I), \textbf{avant de montrer qu'elle constitue aussi le fondement de l'émancipation humaine} (II). \textbf{Il faudra enfin dégager les conditions politiques et éthiques d'une technique réappropriée} (III).}
\emph{\textbf{Entrons dans le vif du sujet par l'analyse marxiste de l'aliénation du travail.}}
\tcblower
\textbf{Analyse de la solution :}
\begin{itemize}
    \item \textbf{Amorce :} Ancrage camerounais (téléphone, paysan), concret, évite le cliché.
    \item \textbf{Analyse :} Définitions philosophiques précises (technique = procédés rationnels ; aliénation = dépossession). Problématique tensionnelle (puissance vs risque).
    \item \textbf{Plan :} Annonce impersonnelle, titres des 3 parties majeures repris fidèlement du plan détaillé.
    \item \textbf{Transition :} Lien fluide vers I-A.
\end{itemize}
\end{exoresolu}

\courssub{Construire un paragraphe argumentatif rigoureux (Méthode I-A-E)}

\begin{methode}[La structure I-A-E : Idée - Argumentation - Exemple/Explication]
\textbf{Objectif :} Rédiger chaque sous-partie (mineure) en un paragraphe unique, cohérent, noté sur 4 à 5 points.
\begin{enumerate}
    \item \textbf{Idée directrice (Phrase thèse) :} Une seule phrase affirmative, réponse partielle à la problématique. Elle ouvre le paragraphe. \textbf{Pas de question.}
    \item \textbf{Argumentation (Le « Pourquoi ») :} 2 à 3 phrases de raisonnement logique. Mobiliser un \textbf{concept} d'auteur au programme (ex: \cle{aliénation}, \cle{gestell}, \cle{prothèse}, \cle{maîtrise de la nature}). Expliquer le mécanisme causal ou logique. \textbf{Ne pas faire une biographie de l'auteur.}
    \item \textbf{Exemple / Illustration / Preuve (Le « Preuve ») :} Un exemple \textbf{concret, daté, situé} (Cameroun, histoire des sciences, actualité technique) OU une expérience de pensée précise. \textbf{Interdit :} « Par exemple, Marx dit que... » (c'est une citation, pas un exemple). L'exemple doit \textbf{prouver} l'idée directrice.
    \item \textbf{Mini-transition (Optionnel) :} Une phrase de fermeture qui relie à la sous-partie suivante.
\end{enumerate}
\cle{Paragraphe unique} \cle{Concept opératoire} \cle{Exemple probant}
\end{methode}

\begin{exoresolu}[Rédiger la sous-partie I.A : « L'aliénation marxiste : le produit du travail lui échappe »]
\textbf{Application I-A-E :}
\begin{enumerate}
    \item \textbf{Idée :} \emph{Dans le système capitaliste, le travailleur est aliéné car le produit de son activité lui devient étranger, s'objectivant en une puissance autonome qui le domine.}
    \item \textbf{Argumentation :} \emph{Selon Marx, le travail est « objectivation » : l'homme met sa vie dans l'objet. Mais sous le capitalisme, cet objet appartient au capitaliste. Le produit devient « capital », valeur qui se valorise elle-même. Plus le travailleur produit, plus il s'appauvrit intérieurement ; plus la chose grandit, plus l'homme diminue. C'est le fétichisme de la marchandise : les rapports sociaux apparaissent comme rapports entre choses.}
    \item \textbf{Exemple (Cameroun) :} \emph{On observe cette dynamique dans les plantations d'hévéa ou de palmier à huile du Littoral ou du Sud-Ouest : le paysan produit le latex ou les régimes, mais le cours mondial (fixé à la Bourse de Kuala Lumpur ou Londres) et les intermédiaires s'approprient la plus-value. Le paysan voit son produit partir en camion-citerne sans maîtriser ni le prix, ni la transformation, ni la destination finale. Son travail « gelé » dans la marchandise lui revient sous forme de salaire de misère.}
    \item \textbf{Transition :} \emph{Mais l'aliénation ne touche pas seulement le produit fini ; elle contamine l'acte même de produire par la division du travail.}
\end{enumerate}
\tcblower
\textbf{Analyse :} L'idée directrice est une thèse philosophique (Marx). L'argumentation utilise les concepts précis : \emph{objectivation, capital, fétichisme}. L'exemple est \textbf{localisé} (plantations Cameroun), \textbf{économiquement précis} (cours mondial, intermédiaires, plus-value) et \textbf{illustre exactement le concept} (produit échappant au producteur). Pas de citation « Marx dit que... » insérée bêtement ; le concept est \textbf{intégré} au raisonnement.
\end{exoresolu}

\courssub{Réaliser une synthèse dialectique : le dépassement (Aufhebung)}

\begin{methode}[La Synthèse : 3 opérations logiques (Nier - Conserver - Élever)]
\textbf{Objectif :} Éviter le « résumé » ou le « ni l'un ni l'autre ». La synthèse est une \textbf{nouvelle thèse} de niveau supérieur.
\begin{enumerate}
    \item \textbf{Identifier le noyau de vérité de la Thèse (I) :} Quel risque réel a-t-elle mis au jour ? (Ex: La technique \textbf{peut} aliéner si elle est subie).
    \item \textbf{Identifier le noyau de vérité de l'Antithèse (II) :} Quelle potentialité a-t-elle révélée ? (Ex: La technique \textbf{est} puissance de libération si elle est maîtrisée).
    \item \textbf{Opérer le dépassement (Aufhebung) :} Formuler le \textbf{principe unificateur} qui rend les deux compatibles. Souvent : \textbf{Condition politique, éthique, anthropologique ou ontologique}. Formule : \emph{« La technique n'est aliénante \textbf{que si} [condition manquante] ; elle devient libératrice \textbf{si} [condition réalisée]. »}
    \item \textbf{Structurer en 2 sous-parties :}
        \begin{itemize}
            \item III.A : Le principe de réconciliation (niveau conceptuel).
            \item III.B : La mise en œuvre concrète / conséquence existentielle ou politique (niveau pratique).
        \end{itemize}
\end{enumerate}
\cle{Aufhebung} \cle{Condition de possibilité} \cle{Niveau supérieur}
\end{methode}

\begin{exoresolu}[Rédiger la Synthèse (Partie III) pour « La technique est-elle aliénante ? »]
\textbf{Application de la méthode :}
\begin{itemize}
    \item \textbf{Vrai de I :} Technique = aliénation \textbf{si} hétéronomie (subie, capitaliste, Gestell).
    \item \textbf{Vrai de II :} Technique = libération \textbf{si} autonomie (maîtrise nature, prothèse, culture).
    \item \textbf{Principe unificateur :} \textbf{L'appropriation collective et éthique.} La technique n'a pas de sens en soi ; son sens dépend du \textbf{régime de propriété} et de la \textbf{responsabilité} de l'usager.
\end{itemize}
\textbf{Rédaction de III.A (Principe) :}
\emph{La contradiction entre aliénation et libération se résout dans l'appropriation : la technique n'est aliénante que lorsqu'elle est subie comme un destin (capitalisme, système technicien), elle devient libératrice lorsqu'elle est collectivement maîtrisée comme un moyen. Comme l'analyse Marx, « l'association des producteurs libres » transforme les moyens de travail en propriété commune ; comme le suggère Simondon, l'éducation technique (connaissance du fonctionnement) dissout l'opacité magique de la machine. La condition de la non-aliénation est donc \textbf{politique} (propriété des moyens) et \textbf{pédagogique} (transparence technique).}
\textbf{Rédaction de III.B (Conséquence) :}
\emph{Cette réappropriation impose une éthique de la responsabilité (Jonas) : l'homme, « berger de l'être » (Heidegger), doit orienter l'innovation vers le soin du vivant et non vers la seule performance. Au Cameroun, les « FabLabs » (Yaoundé, Douala) ou les initiatives d'appropriation technologique locale (fabrication de broyeuses, séchoirs solaires adaptés par des ingénieurs locaux pour des coopératives paysannes) illustrent ce dépassement : la technique n'est plus importée subie, mais \textbf{co-construite} pour répondre à des besoins identifiés. L'aliénation n'est pas le destin de la technique, mais le symptôme d'une dépossession politique.}
\tcblower
\textbf{Analyse :} La synthèse ne dit pas « c'est les deux ». Elle dit : « C'est l'appropriation qui décide ». Elle mobilise Marx (politique), Simondon (pédagogique/culturel), Jonas/Heidegger (éthique/ontologique). L'exemple camerounais (FabLabs, innovation locale) ancre le concept abstrait dans le réel technique actuel. C'est une \textbf{thèse nouvelle} : « La technique est aliénante \textbf{faute d'appropriation} ».
\end{exoresolu}

\courssub{Rédiger une conclusion en deux temps : Bilan \& Ouverture}

\begin{methode}[Conclusion : Bilan synthétique + Ouverture maîtrisée]
\textbf{Objectif :} Clore le devoir sans rien ajouter de nouveau, mais en élargissant l'horizon.
\begin{enumerate}
    \item \textbf{Bilan (Rappel dialectique) :} 2-3 phrases. Reprendre le fil : \emph{« Au terme de cette analyse, il apparaît que [Thèse I]... Cependant [Antithèse II]... C'est pourquoi [Synthèse III : condition]. »} \textbf{Interdit :} « Pour conclure, j'ai montré que... » (métadiscours).
    \item \textbf{Ouverture (Perspective) :} 1-2 phrases. Poser une \textbf{question nouvelle}, liée, mais qui dépasse le cadre strict du sujet. Trois types d'ouvertures valides :
        \begin{itemize}
            \item \textbf{Autre discipline :} Histoire, Droit, Économie, Sciences exactes.
            \item \textbf{Autre notion du programme :} Liberté, Travail, Politique, Nature, Art, Religion, Science, Vérité, Justice, Bonheur, État, Langage.
            \item \textbf{Actualité brûlante / Problème contemporain :} IA, Bioéthique, Écologie, Numérique.
        \end{itemize}
        \textbf{Formule :} \emph{« Cette réflexion invite à se demander si / comment / jusqu'où [Nouvelle question]... »}
    \item \textbf{Dernière phrase (Chute) :} Une formule courte, marquante, sans point d'interrogation final si possible (affirmation ouverte).
\end{enumerate}
\cle{Bilan dialectique} \cle{Ouverture} \cle{Notion programme}
\end{methode}

\begin{exoresolu}[Rédiger la conclusion pour « La technique est-elle aliénante ? »]
\textbf{Proposition :}
\emph{Au terme de cette analyse, il apparaît que la technique n'est pas aliénante par essence, mais le devient par la forme sociale qui la porte : le capitalisme la fétichise, le système technicien l'autonomise. Cependant, elle demeure la condition sine qua non de l'humanisation de l'homme par la maîtrise de la nature. C'est donc dans l'appropriation collective des moyens techniques et l'éducation à la transparence des dispositifs que se joue la frontière entre asservissement et émancipation.}
\emph{Cette conclusion nous oblige à réinterroger la notion de \textbf{progrès} : le progrès technique coïncide-t-il nécessairement avec un progrès moral et politique, ou la dissociation entre puissance technique et sagesse politique (Jonas) menace-t-elle de rendre l'aliénation irréversible à l'ère de l'intelligence artificielle et du transhumanisme ?}
\emph{\textbf{Maîtriser la technique, c'est encore et toujours maîtriser le sens que nous lui donnons.}}
\tcblower
\textbf{Analyse :}
\begin{itemize}
    \item \textbf{Bilan :} Résume la dialectique (Essence vs Forme sociale ; Condition vs Risque ; Appropriation comme clé). Pas de « Je pense que ».
    \item \textbf{Ouverture :} Sur la notion de \textbf{Progrès} (au programme Terminale : « Progrès et Histoire » / « Technique »). Lien avec Jonas (éthique responsabilité) et actualité (IA, transhumanisme).
    \item \textbf{Chute :} Affirmation forte, retour sur l'humain (sujet de la responsabilité).
\end{itemize}
\end{exoresolu}

\courssub{Choisir, intégrer et référencer une citation philosophique}

\begin{methode}[Règles d'or de la citation : Pertinence - Intégration - Référence]
\textbf{Objectif :} La citation est une \textbf{preuve d'autorité} ou un \textbf{levier conceptuel}, jamais une décoration.
\begin{enumerate}
    \item \textbf{Choix (Pertinence) :} La citation doit \textbf{exprimer exactement le concept} que vous venez d'argumenter (I-A-E). Elle doit venir d'un auteur \textbf{au programme} (Terminale Tech : Platon, Aristote, Descartes, Kant, Hegel, Marx, Nietzsche, Bergson, Heidegger, Jonas, Simondon, Arendt, Rawls, Nussbaum, auteurs africains : Mudimbe, Mbembe, Wiredu, Hountondji...).
    \item \textbf{Intégration syntaxique (La « couture ») :} \textbf{Jamais de citation isolée entre guillemets.} Trois modes :
        \begin{itemize}
            \item \textbf{Introducteur + deux-points :} \emph{Marx affirme : « Le travailleur ne se sent pas chez lui dans son travail... »}
            \item \textbf{Incidence (relatif) :} \emph{« Le travailleur ne se sent pas chez lui dans son travail », écrit Marx, soulignant ainsi l'extériorité du travail.}
            \item \textbf{Intégration dans la phrase (ellipse des guillemets si court) :} \emph{Pour Marx, le travail aliéné est celui où « l'activité n'est pas sienne ».}
        \end{itemize}
    \item \textbf{Référencement minimal (Note de bas de page ou parenthèse) :} \textbf{Obligatoire au Bac.} Format : \emph{(Auteur, Œuvre, §/Page/Édition)}. Ex: \emph{(Marx, Manuscrits de 1844, § 44)}. Si œuvre inconnue : \emph{(Marx, Manuscrits de 1844)}.
    \item \textbf{Analyse post-citation :} \textbf{Jamais finir un paragraphe par une citation.} Il faut \textbf{la commenter} : \emph{« Cette formule montre que... », « Ce passage illustre le concept de... »}.
\end{enumerate}
\cle{Introducteur} \cle{Commentaire obligatoire} \cle{Référence précise}
\end{methode}

\begin{exoresolu}[Intégrer une citation dans le paragraphe I.A (Marx/aliénation)]
\textbf{Paragraphe I.A révisé avec citation :}
\emph{Dans le système capitaliste, le travailleur est aliéné car le produit de son activité lui devient étranger, s'objectivant en une puissance autonome qui le domine. \textbf{Marx décrit avec force cette dépossession :} « L'objectivation du travail apparaît comme perte de l'objet [...] l'ouvrier ne se sent pas chez lui dans son travail [...] son travail n'est donc pas volontaire, mais forcé, travail forcé » (\textbf{Manuscrits de 1844}, § 44). \textbf{Cette citation révèle que l'aliénation n'est pas seulement économique (vol de la plus-value) mais existentielle :} l'activité même devient souffrance, « mortification du corps et ruine de l'esprit ». On observe cette dynamique dans les plantations d'hévéa du Littoral : le paysan voit son produit partir en camion-citerne sans maîtriser ni le prix, ni la transformation. Son travail « gelé » dans la marchandise lui revient sous forme de salaire de misère, confirmant que « l'objet que le travail produit, son produit, lui est hostile, étranger ».}
\tcblower
\textbf{Analyse :}
\begin{itemize}
    \item \textbf{Choix :} Extrait exact du §44 (Manuscrits), cœur du concept d'aliénation.
    \item \textbf{Intégration :} Introducteur (« Marx décrit... »), citation longue en bloc (ou guillemets si court), référence précise entre parenthèses.
    \item \textbf{Commentaire :} « Cette citation révèle que... » $\rightarrow$ Explication du sens (existentiel vs économique).
    \item \textbf{Relance :} L'exemple camerounais suit \textbf{après} le commentaire pour illustrer le concept cité.
\end{itemize}
\end{exoresolu}

\courssub{Gérer l'épreuve : Méthode de gestion du temps \& Auto-évaluation (Grille Bac)}

\begin{methode}[Stratégie Bac 4h : Chronologie \& Auto-correction]
\textbf{Objectif :} Rendement maximal en 4 heures. Réservez 15 min de relecture finale.
\begin{enumerate}
    \item \textbf{0h00 - 0h30 : Brouillon intensif.} Analyse sujet (5'), Problématique (5'), Plan détaillé + Idées directrices I-A-E pour \textbf{toutes} les sous-parties (20'). \textbf{Ne pas rédiger l'intro/conclusion au brouillon.} Juste mots-clés.
    \item \textbf{0h30 - 1h15 : Introduction + Partie I (Thèse).} Rédiger au propre directement. 1h15 max pour 1ère partie.
    \item \textbf{1h15 - 2h00 : Partie II (Antithèse).} Même rythme.
    \item \textbf{2h00 - 2h45 : Partie III (Synthèse).} Le moment le plus délicat. Calme.
    \item \textbf{2h45 - 3h15 : Conclusion + Relecture globale.} Rédiger conclusion. Relire \textbf{à voix basse} (repère syntaxe, répétitions, accords).
    \item \textbf{3h15 - 4h00 : Relecture critique (Grille d'auto-évaluation).} Cocher chaque case.
\end{enumerate}
\cle{Gestion du temps} \cle{Brouillon structuré} \cle{Auto-évaluation}
\end{methode}

\begin{exoresolu}[Grille d'auto-évaluation « Zéro faute » (À cocher pendant la dernière heure)]
\textbf{Appliquez cette grille à votre copie blanche :}
\begin{center}
\begin{tabularx}{\linewidth}{|l|X|c|}\hline
\textbf{Critère Bac} & \textbf{Question de contrôle} & \textbf{OK ?} \\ \hline
\textbf{Introduction} & Amorce concrète (pas cliché) ? Termes définis philo ? Problématique tensionnelle (1 phrase) ? Plan annoncé (3 titres majeurs) ? Transition vers I ? & $\square$ \\ \hline
\textbf{Structure} & 3 Parties (I, II, III) visibles ? 2 ou 3 sous-parties par partie ? Titres nominatifs parallèles ? & $\square$ \\ \hline
\textbf{Paragraphes (I-A-E)} & Chaque sous-partie = 1 paragraphe ? Idée directrice en 1ère phrase ? Argumentation conceptuelle (auteur/concept) ? Exemple concret/précis (Cameroun/histoire) ? Mini-transition ? & $\square$ \\ \hline
\textbf{Dialectique} & II contredit vraiment I (pas répétition) ? III dépasse (Aufhebung) et ne résume pas ? Principe unificateur explicite ? & $\square$ \\ \hline
\textbf{Citations} & Au moins 3-4 citations précises au fil du texte ? Intégrées (introducteur) ? Référencées (Auteur, Œuvre) ? Commentées (pas fin de paragraphe) ? & $\square$ \\ \hline
\textbf{Conclusion} & Bilan dialectique (I/II/III) sans « Je » ? Ouverture sur notion programme / actualité / discipline ? Chute marquante ? & $\square$ \\ \hline
\textbf{Langue/Style} & 0 faute d'orthographe/grammaire (relire à l'envers) ? Phrases courtes, vocabulaire précis (pas « chose », « truc ») ? Connecteurs logiques (donc, cependant, par conséquent, en effet) ? & $\square$ \\ \hline
\end{tabularx}
\end{center}
\textbf{Si une case n'est pas cochée :} Corrigez \textbf{maintenant} (rajouter mot, changer ponctuation, insérer référence).
\tcblower
\textbf{Analyse :} Cette grille condense les attendus officiels du barème Bac (Structure 4pts, Contenu 10pts, Expression 6pts). L'exercice résolu consiste à \textbf{s'imaginer correcteur} : on ne relit pas « pour voir », on vérifie \textbf{point par point}. C'est la seule méthode pour gagner les 2-3 points de « présentation/rigueur » qui font la différence entre 11 et 14.
\end{exoresolu}

\begin{attention}[Les 5 erreurs qui coûtent des points au Bac (Terminale Tech)]
\begin{enumerate}
    \item \textbf{Le « Plan Pour/Contre » (Note max 8/20) :} Faire I. « Oui, la technique libère », II. « Non, elle aliène », III. « Elle fait les deux ». \textbf{Solution :} La dialectique exige que II soit une \textbf{nécessité interne} (structure de la technique), pas une simple opinion inverse. III doit \textbf{dépasser} (Aufhebung), pas mixer.
    \item \textbf{L'introduction « catalogue » ou « biographique » :} « Platon dit... Aristote dit... Marx dit... Je vais étudier... ». \textbf{Solution :} Amorce concrète $\rightarrow$ Définitions $\rightarrow$ Problématique $\rightarrow$ Plan. Pas de noms d'auteurs dans l'intro (sauf si sujet cite un auteur).
    \item \textbf{La citation « parachutée » (non intégrée, non commentée) :} Fin de paragraphe : « ... comme le dit Marx : "Le travail est aliéné." » (Point final). \textbf{Solution :} Introducteur + Citation + Référence + \textbf{Commentaire explicatif} + Exemple/Relance.
    \item \textbf{L'exemple « vague » ou « personnel » :} « Par exemple, moi j'utilise mon téléphone... », « Dans la vie de tous les jours... ». \textbf{Solution :} Exemple \textbf{daté, situé, technique, économique, historique} (ex: « L'introduction du téléphone mobile au Cameroun en 1999... », « L'affaire Volkswagen 'Dieselgate' 2015... », « Les FabLabs de Yaoundé depuis 2015... »).
    \item \textbf{La conclusion « ouverture hors-sujet » ou « question rhétorique bête » :} « Et si la technique était l'avenir de l'homme ? », « Le bonheur est-il possible ? ». \textbf{Solution :} Ouverture sur une \textbf{notion précise du programme} (Progrès, Liberté, Responsabilité, Nature, Travail, Justice, Art, Religion, Science, Vérité, État, Langage) ou un \textbf{problème technique actuel} (IA, CRISPR, Énergie, Numérique).
\end{enumerate}
\end{attention}

\begin{aretenir}[Synthèse des exigences Bac pour la dissertation]
\begin{itemize}
    \item \textbf{Structure immuable :} Introduction (4 temps) $\rightarrow$ Développement dialectique (3 parties, 2-3 sous-parties, paragraphes I-A-E) $\rightarrow$ Conclusion (Bilan + Ouverture).
    \item \textbf{Concepts opératoires :} Chaque argumentation repose sur un concept d'auteur programme (pas d'opinion personnelle).
    \item \textbf{Exemples probants :} Ancrage réel (Cameroun, histoire des sciences, actualité technique), datés, précis.
    \item \textbf{Citations intégrées :} Introducteur + Citation + Référence + Commentaire. Minimum 3-4 au total.
    \item \textbf{Synthèse = Aufhebung :} Dépassement par une condition (politique, éthique, anthropologique), pas compromis.
    \item \textbf{Gestion du temps :} 30' brouillon (plan complet), 2h15 rédaction parties, 30' intro/conclusion, 45' relecture grille.
\end{itemize}
\end{aretenir}

\newpage

\coursec{Exercices}

\courssub{Je vérifie mes connaissances}

\begin{exercice}[QCM : Règles méthodologiques fondamentales]
\textbf{Pour chaque question, une seule réponse est exacte. Cochez la bonne lettre.}
\begin{enumerate}
    \item Dans une introduction de dissertation philosophique, l'\textbf{accroche} doit :
    \begin{itemize}
        \item[A.] Résumer le sujet en une phrase.
        \item[B.] Être une citation d'auteur obligatoirement.
        \item[C.] Introduire le thème de manière concrète et originale (actualité, fait de société, référence culturelle camerounaise) sans énoncer le sujet.
        \item[D.] Donner la thèse que l'on va défendre.
    \end{itemize}
    \item La \textbf{problématique} est :
    \begin{itemize}
        \item[A.] La reformulation du sujet sous forme de question directe.
        \item[B.] La tension entre deux réponses plausibles et opposées au sujet, qui justifie le plan.
        \item[C.] L'annonce du plan détaillé (``En premier lieu... Ensuite... Enfin...'').
        \item[D.] La définition de tous les termes du sujet.
    \end{itemize}
    \item La structure \textbf{I-A-E} d'un paragraphe argumentatif signifie :
    \begin{itemize}
        \item[A.] Idée - Auteur - Explication.
        \item[B.] Introduction - Argumentation - Exemple.
        \item[C.] Idée directrice - Argumentation (analyse logique) - Exemple/Illustration concrète.
        \item[D.] Intuition - Analyse - Évaluation.
    \end{itemize}
    \item La \textbf{synthèse} (troisième partie du développement) doit obligatoirement :
    \begin{itemize}
        \item[A.] Résumer la thèse et l'antithèse.
        \item[B.] Apporter de nouveaux arguments sans lien avec les parties précédentes.
        \item[C.] Dépasser l'opposition en montrant pourquoi l'antithèse est partielle et comment la thèse l'intègre ou la corrige (niveau supérieur de compréhension).
        \item[D.] Donner son opinion personnelle définitive (``Je pense que...'').
    \end{itemize}
\end{enumerate}
\end{exercice}

\begin{exercice}[Vrai / Faux justifié : Pièges classiques]
\textbf{Indiquez VRAI ou FAUX et justifiez brièvement votre réponse (une phrase).}
\begin{enumerate}
    \item ``L'accroche \emph{De tout temps, les hommes se sont posé la question de...} est une entrée en matière recommandée au Baccalauréat.''
    \item ``Une citation en langue étrangère (anglais, latin) peut être insérée sans traduction si elle est célèbre.''
    \item ``Le plan dialectique (Thèse / Antithèse / Synthèse) est le seul plan acceptable pour une dissertation philosophique en Terminale technique.''
    \item ``Dans la conclusion, l'ouverture doit obligatoirement poser une nouvelle question philosophique sans y répondre.''
    \item ``Citer un auteur (ex : Aristote, Kant, Paulin Hountondji) sans expliquer en quoi son idée soutient l'argumentation suffit à valider l'argument.''
    \item ``L'annonce du plan doit utiliser les connecteurs logiques \emph{Premièrement, Deuxièmement, Troisièmement}.''
    \item ``Un exemple technique camerounais (ex : barrage de Lom Pangar, code de déontologie de l'Ordre National des Ingénieurs de Génie Civil du Cameroun) a plus de valeur qu'un exemple historique européen.''
    \item ``La définition des termes du sujet doit être exhaustive (tous les mots) et placée avant la problématique.''
\end{enumerate}
\end{exercice}

\begin{exercice}[Définitions et complétion : Vocabulaire technique]
\textbf{Complétez les définitions suivantes ou donnez la notion correspondante.}
\begin{enumerate}
    \item Notion : \emph{``Tension logique entre deux thèses contraires qui ne peuvent être vraies ensemble, mais qui expriment chacune une part de vérité, nécessitant un dépassement.''} \rightarrow \dotfill
    \item Notion : \emph{``Étape de l'introduction où l'on explicite le sens précis des concepts clés du sujet pour lever les ambiguïtés.''} \rightarrow \dotfill
    \item Notion : \emph{``Partie de la conclusion où l'on élargit la réflexion en reliant le sujet traité à une autre problématique, un autre domaine (technique, artistique, politique) ou une actualité.''} \rightarrow \dotfill
    \item Règle : \emph{``Une citation doit toujours être \textbf{intégrée} à la phrase (syntaxe), \textbf{signalée} (guillemets), \textbf{référencée} (auteur, œuvre) et surtout \textbf{\dotfill} (expliciter son sens et sa portée pour l'argument).''}
\end{enumerate}
\end{exercice}

\begin{exercice}[Quiz méthodologique chronométré (10 min) : 10 items rapides]
\textbf{Répondez par Vrai (V) ou Faux (F) sans justification. Objectif : 10/10 en moins de 10 minutes.}
\begin{enumerate}
    \item La thèse correspond toujours à l'opinion commune (le sens commun).
    \item L'antithèse doit contredire la thèse point par point dans le même ordre.
    \item La synthèse est une simple conclusion partielle avant la conclusion finale.
    \item On peut utiliser ``Je pense que'' dans le développement pour affirmer sa position.
    \item Une citation orpheline (insérée sans lien avec la phrase) est une erreur grave.
    \item L'accroche doit obligatoirement citer un philosophe.
    \item La problématique se termine toujours par un point d'interrogation.
    \item L'annonce du plan révèle la structure dialectique (Thèse, Antithèse, Synthèse) explicitement.
    \item Une référence à l'actualité technique camerounaise (ex : fibre optique, exploitation minière) ne compte pas comme un ``exemple'' au sens I-A-E.
    \item La conclusion doit reprendre mot pour mot la thèse de la synthèse.
\end{enumerate}
\end{exercice}

\courssub{Je m'entraîne}

\begin{exercice}[Analyse de sujets : Décortiquer pour ne pas être hors sujet]
\textbf{Pour chacun des deux sujets ci-dessous, effectuez le travail suivant :}
\begin{enumerate}
    \item Identifiez les \textbf{termes clés} à définir (2 à 3 maximum).
    \item Formulez l'\textbf{opposition fondamentale} (le conflit de sens) en une phrase.
    \item Rédigez une \textbf{problématique} dialectique (forme interrogative, tension thèse/antithèse).
    \item Esquissez un \textbf{plan dialectique} (Thèse / Antithèse / Synthèse) avec un titre de partie clair pour chacune.
\end{enumerate}
\textbf{Sujet A :} \emph{``L'ingénieur a-t-il une responsabilité morale spécifique ?''}

\textbf{Sujet B :} \emph{``La transparence totale est-elle une vertu technique ?''} (données ouvertes, secret industriel, algorithmes)
\end{exercice}

\begin{exercice}[Rédaction d'une introduction canonique (4 étapes)]
\textbf{Sur le Sujet A (\emph{``L'ingénieur a-t-il une responsabilité morale spécifique ?''}), rédigez une introduction complète respectant scrupuleusement les 4 étapes :}
\begin{enumerate}
    \item \textbf{Accroche} : ancrage dans le réel camerounais (ex : effondrement d'immeubles à Douala, gestion des données numériques, code de déontologie de l'Ordre National des Ingénieurs de Génie Civil du Cameroun).
    \item \textbf{Définition des termes} : ``ingénieur'' (compétence technique, pouvoir d'agir), ``responsabilité'' (imputabilité, obligation de répondre de ses actes), ``spécifique'' (propre à la fonction, distincte de celle du citoyen ordinaire).
    \item \textbf{Problématique} : tension entre neutralité technique (l'ingénieur comme simple exécutant) et engagement éthique (l'ingénieur comme gardien de la sécurité et du bien commun).
    \item \textbf{Annonce du plan} : sans citer ``Thèse/Antithèse/Synthèse'', mais en annonçant la logique (ex : ``Nous verrons d'abord que... Nous montrerons ensuite que... Nous comprendrons enfin que...'').
\end{enumerate}
\end{exercice}

\begin{exercice}[Construction de paragraphes I-A-E : Thèse et Antithèse]
\textbf{Rédigez deux paragraphes complets (environ 15 lignes chacun) en appliquant la structure I-A-E (Idée directrice - Argumentation - Exemple/Illustration).}
\begin{enumerate}
    \item \textbf{Paragraphe de THÈSE} : défendre l'idée que l'ingénieur a \textbf{bien} une responsabilité morale spécifique (ex : pouvoir technique = devoir de prévention, asymétrie d'information, serment professionnel).
    \item \textbf{Paragraphe d'ANTITHÈSE} : défendre l'idée que l'ingénieur \textbf{n'a pas} de responsabilité morale \emph{spécifique} (ex : responsabilité partagée et diluée, subordination hiérarchique, neutralité axiologique de la technique, distinction entre responsabilité juridique et responsabilité morale).
\end{enumerate}
\textit{Contrainte : chaque paragraphe doit contenir au moins une référence à un auteur (ex : Jonas, Ellul, Hountondji) et un exemple concret camerounais ou technique.}
\end{exercice}

\begin{exercice}[Gestion des citations : Intégration, référencement, exploitation]
\textbf{À partir des citations brutes ci-dessous, produisez pour chacune :}
\begin{enumerate}
    \item Une \textbf{intégration syntaxique correcte} dans une phrase argumentative (pas de citation orpheline).
    \item Le \textbf{référencement précis} (auteur, œuvre).
    \item Une \textbf{phrase d'exploitation} (herméneutique) : en quoi cette citation soutient-elle ou contredit-elle l'argument en cours ?
\end{enumerate}
\textbf{Citations brutes :}
\begin{enumerate}
    \item ``La technique n'est pas un destin, c'est un choix.'' (formule reprise de Jacques Ellul, théoricien de la technique, auteur de \emph{La Technique ou l'Enjeu du siècle}).
    \item ``Agis de telle sorte que les effets de ton action soient compatibles avec la permanence d'une vie authentiquement humaine sur terre.'' (Hans Jonas, \emph{Le Principe responsabilité}, 1979).
    \item ``Le savoir technique est un pouvoir ; tout pouvoir appelle une éthique.'' (formule inspirée de la réflexion de Paulin Hountondji sur le savoir et le pouvoir en Afrique).
    \item ``L'ingénieur est quelqu'un qui fait pour dix francs ce que n'importe quel imbécile ferait pour cent.'' (adage humoristique du milieu de l'ingénierie, attribution incertaine).
\end{enumerate}
\end{exercice}

\begin{attention}[Fiabilité des citations]
Au Baccalauréat, une citation mal attribuée coûte des points. Si l'attribution est incertaine (comme la citation 4 ci-dessus), il vaut mieux écrire ``comme le rappelle un adage du milieu de l'ingénierie'' plutôt que d'inventer un auteur. Ne citez que ce dont vous êtes sûr, et exploitez toujours la citation après l'avoir donnée.
\end{attention}

\courssub{Je me prépare au Bac}

\begin{exercice}[Sujet type Bac technique : L'ingénieur et la responsabilité morale]
\textbf{Sujet :} \emph{``L'ingénieur a-t-il une responsabilité morale spécifique ?''}

\textbf{Travail demandé (simulation d'examen, 4 h) :}
\begin{enumerate}
    \item \textbf{Analyse du sujet (brouillon structuré) :} définitions rigoureuses, identification du présupposé (l'existence d'une morale applicable à la technique), formulation de trois problématiques possibles, choix de la meilleure.
    \item \textbf{Plan détaillé dialectique :} pour chaque partie (I, II, III), donnez l'idée directrice générale, les 2 ou 3 arguments principaux (avec auteurs et concepts), et les exemples techniques camerounais pertinents (BTP, mines, NTIC, énergie, Ordre National des Ingénieurs de Génie Civil).
    \item \textbf{Rédaction de l'introduction complète} (accroche contextuelle, définitions, problématique, annonce de plan).
    \item \textbf{Rédaction de la SYNTHÈSE seule (partie III) :} dépassement dialectique (ex : responsabilité partagée mais asymétrique ; éthique de la responsabilité contre éthique de la conviction ; rôle de la régulation étatique et de la société civile au Cameroun).
    \item \textbf{Rédaction de la CONCLUSION complète} (bilan synthétique, réponse définitive à la problématique, ouverture maîtrisée sur l'éthique de l'intelligence artificielle ou la transition écologique).
\end{enumerate}
\end{exercice}

\begin{exercice}[Sujet type Probatoire : Transparence, secret et technique]
\textbf{Sujet :} \emph{``Faut-il tout dire ?''} \textbf{Variante technique :} \emph{``La transparence totale est-elle une vertu technique ?''} (données ouvertes, algorithmes explicables, secret industriel, sécurité nationale, protection des données personnelles selon la loi camerounaise de 2010 sur la cybersécurité et la cybercriminalité).

\textbf{Travail demandé :}
\begin{enumerate}
    \item \textbf{Distinguez les deux sujets :} en quoi la variante ``technique'' déplace-t-elle le débat philosophique classique (morale interpersonnelle) vers l'épistémologie et l'éthique professionnelle ?
    \item \textbf{Problématisation :} formulez une problématique unique qui articule \emph{transparence} (visibilité, intelligibilité, accès) et \emph{vertu technique} (efficacité, fiabilité, confiance, innovation).
    \item \textbf{Plan dialectique (Thèse/Antithèse/Synthèse) :} titres des parties + arguments clés + exemples (données ouvertes gouvernementales au Cameroun, boîtes noires algorithmiques, brevets pharmaceutiques, lanceurs d'alerte).
    \item \textbf{Rédigez un paragraphe I-A-E de THÈSE} (la transparence comme condition de confiance et de contrôle démocratique et éthique).
    \item \textbf{Rédigez un paragraphe I-A-E d'ANTITHÈSE} (la nécessaire opacité : secret industriel, sécurité, protection de la vie privée, complexité technique irréductible).
    \item \textbf{Esquissez la SYNTHÈSE :} la \emph{transparence raisonnée} (redevabilité, \emph{accountability}) comme vertu technique supérieure à la transparence brute.
\end{enumerate}
\end{exercice}

\begin{exercice}[Réécriture critique : De la copie 8/20 à la copie 16/20]
\textbf{Voici des extraits anonymisés d'une copie d'élève (note : 8/20). Les défauts majeurs sont : plan invisible, absence de synthèse (simple troisième partie descriptive), citations orphelines, conclusion bâclée, exemples absents.}

\textbf{Réécrivez intégralement les quatre segments ci-dessous en appliquant la méthode du cours (I-A-E, intégration des citations, dialectique, ancrage camerounais).}

\textbf{Extrait 1 -- Introduction de l'élève :}

``De tout temps, les ingénieurs construisent des ponts. Le sujet demande s'ils ont une responsabilité morale. Responsabilité veut dire répondre de ses actes. Morale c'est le bien et le mal. Spécifique veut dire propre à l'ingénieur. Nous allons voir que oui, ils ont une responsabilité. Ensuite que non. Enfin on conclura.''

\textbf{Extrait 2 -- Paragraphe de thèse de l'élève :}

``L'ingénieur a une responsabilité. Jonas dit "Agis de telle sorte que...". C'est vrai car il fait des choix techniques. Exemple : le barrage. Kant dit l'homme est une fin. Donc l'ingénieur doit respecter l'homme.''

\textbf{Extrait 3 -- ``Synthèse'' de l'élève (en réalité troisième partie descriptive) :}

``En fait, c'est compliqué. Il y a des lois. Au Cameroun il y a l'Ordre des Ingénieurs. Il faut respecter le code. L'État doit contrôler. Les entreprises aussi. Donc tout le monde est responsable.''

\textbf{Extrait 4 -- Conclusion de l'élève :}

``Pour conclure, l'ingénieur a une responsabilité spécifique. C'est important pour la société. Il faut faire attention. On peut se demander si les robots auront une responsabilité.''

\textbf{Travail attendu :}
\begin{enumerate}
    \item \textbf{Nouvelle introduction} (4 étapes, accroche camerounaise, définitions opératoires, problématique dialectique, annonce de plan implicite).
    \item \textbf{Nouveau paragraphe de thèse (I-A-E)} : idée directrice claire (ex : pouvoir technique $\rightarrow$ devoir de prévention), argumentation structurée (Jonas, asymétrie d'information), exemple concret (barrage de Lom Pangar ou effondrements d'immeubles à Douala), citation intégrée et exploitée.
    \item \textbf{Nouvelle synthèse (partie III)} : dépassement (responsabilité partagée $\neq$ responsabilité diluée ; asymétrie des savoirs ; régulation), niveau supérieur (éthique de la responsabilité et gouvernance technique).
    \item \textbf{Nouvelle conclusion} : bilan nuancé (oui, responsabilité spécifique fondée sur l'asymétrie et la puissance d'agir), réponse à la problématique, ouverture sur l'intelligence artificielle ou l'ingénieur de demain.
\end{enumerate}
\end{exercice}

\begin{methode}[Grille d'auto-évaluation avant de rendre la copie]
\begin{enumerate}
    \item \textbf{Introduction} : accroche concrète et non clichée ? Termes clés définis (2 ou 3) ? Problématique interrogative et dialectique ? Plan annoncé sans être nommé ?
    \item \textbf{Développement} : chaque paragraphe suit-il I-A-E ? Chaque citation est-elle intégrée, référencée et exploitée ? Y a-t-il au moins un exemple camerounais ou technique par partie ?
    \item \textbf{Synthèse} : dépasse-t-elle réellement l'opposition au lieu de la résumer ?
    \item \textbf{Conclusion} : bilan, réponse claire à la problématique, ouverture liée au sujet et non traitée ?
    \item \textbf{Langue} : aucun ``je pense que'', orthographe et connecteurs logiques soignés ?
\end{enumerate}
\end{methode}

\newpage

\coursec{Corrigés des exercices}

\coursec{Leçon N02 : Consolidation des acquis de la dissertation philosophique}

\courssub{Les fondements et l'architecture globale de la dissertation}

\begin{aretenir}[Architecture type de la dissertation philosophique au Bac Technique]
La dissertation suit une architecture rigoureuse en trois temps majeurs :
\begin{enumerate}
    \item \textbf{L'Introduction} (4 étapes canoniques : Accroche -- Définitions -- Problématique -- Annonce du plan).
    \item \textbf{Le Développement} (généralement dialectique en 3 parties : \textbf{I. Thèse} -- \textbf{II. Antithèse} -- \textbf{III. Synthèse}, chacune structurée en paragraphes \cle{I-A-E}).
    \item \textbf{La Conclusion} (2 temps : Bilan définitif -- Ouverture maîtrisée).
\end{enumerate}
\emph{Principes directeurs :} rigueur conceptuelle, argumentation logique (pas d'opinion), ancrage dans le réel (exemples camerounais), exploitation des citations (règle des 4 verbes : intégrer, signaler, référencer, exploiter).
\end{aretenir}

\begin{methode}[Analyser un sujet : les 4 opérations du brouillon]
\begin{enumerate}
    \item \textbf{Définir} les termes clés (2 ou 3) : sens précis, levée des ambiguïtés.
    \item \textbf{Identifier} le présupposé et l'opposition fondamentale (thèse/antithèse).
    \item \textbf{Formuler} la problématique : question unique, dialectique, terminant par un point d'interrogation.
    \item \textbf{Construire} le plan détaillé : titres des 3 parties + arguments principaux + exemples + auteurs/citations par partie.
\end{enumerate}
\end{methode}

\begin{corrige}[Exercice 1 — QCM : Règles méthodologiques fondamentales]
\begin{enumerate}
    \item \textbf{Réponse C.} L'accroche part du réel (fait de société, actualité, référence culturelle camerounaise) pour conduire naturellement au thème. Elle ne doit ni énoncer le sujet (A est faux : résumer le sujet n'est pas accrocher), ni annoncer la thèse (D est faux : ce serait brûler l'étape de la problématisation). B est faux : la citation n'est jamais obligatoire.
    \item \textbf{Réponse B.} La problématique exprime la \cle{tension} entre deux réponses plausibles et opposées ; c'est elle qui légitime le plan dialectique. A est faux : reformuler le sujet en question n'est pas problématiser. C décrit l'annonce du plan, D décrit la définition des termes.
    \item \textbf{Réponse C.} I-A-E = \cle{Idée directrice} -- \cle{Argumentation} (analyse logique) -- \cle{Exemple} (illustration concrète). L'exemple vérifie l'argument, il ne le remplace pas.
    \item \textbf{Réponse C.} La synthèse est un \cle{dépassement} : elle montre en quoi chaque position précédente est partielle et construit une réponse nuancée de niveau supérieur. A décrit un résumé (insuffisant), D est l'opinion personnelle (interdite dans une dissertation).
\end{enumerate}
\textbf{Points de vigilance :} au Bac, confondre problématique et annonce du plan est la faute la plus fréquente ; retenez que la problématique \emph{pose} le problème, l'annonce du plan \emph{organise} la réponse.
\end{corrige}

\begin{corrige}[Exercice 2 — Vrai / Faux justifié : Pièges classiques]
\begin{enumerate}
    \item \textbf{FAUX.} « De tout temps, les hommes... » est une formule clichée, vide et généalogiquement fausse, sanctionnée par les correcteurs ; l'accroche doit être concrète, datée, située.
    \item \textbf{FAUX.} Toute citation en langue étrangère doit être traduite en français, langue de l'épreuve ; la célébrité d'une formule ne dispense pas de la traduction.
    \item \textbf{FAUX.} Le plan dialectique (thèse/antithèse/synthèse) est le plus fréquent et le plus sûr au programme de Terminale technique, mais un plan thématique ou analytique rigoureux peut convenir selon la nature du sujet.
    \item \textbf{VRAI.} L'ouverture pose une question nouvelle, liée au sujet, sans la traiter : elle élargit la réflexion sans rouvrir la démonstration ni affaiblir le bilan.
    \item \textbf{FAUX.} Une citation non commentée est une \cle{citation orpheline} : elle n'a aucune valeur argumentative, car c'est l'exploitation (l'explication de son lien avec l'argument) qui la fait fonctionner.
    \item \textbf{FAUX.} « Premièrement, Deuxièmement, Troisièmement » sont des connecteurs scolaires et rigides ; on préfère des formulations souples (« Nous verrons d'abord que... Il apparaîtra ensuite que... ») qui annoncent la logique du plan sans le nommer.
    \item \textbf{VRAI} (à nuancer). Un exemple proche du contexte de l'élève (barrage de Lom Pangar, code de déontologie de l'ONIGC) prouve la capacité d'application des notions à la vie courante et est valorisé ; mais un exemple européen pertinent reste parfaitement acceptable.
    \item \textbf{FAUX.} On ne définit que les termes \emph{clés} (2 ou 3), ceux dont le sens est ambigu ou décisif pour le problème ; une définition exhaustive alourdit l'introduction et retarde la problématique.
\end{enumerate}
\end{corrige}

\begin{corrige}[Exercice 3 — Définitions et complétion : Vocabulaire technique]
\begin{enumerate}
    \item La \textbf{problématique} (ou contradiction dialectique) : tension entre deux thèses contraires, chacune partiellement vraie, qui exige un dépassement.
    \item La \textbf{définition des termes du sujet} (analyse des notions) : explicitation du sens précis des concepts clés pour lever les ambiguïtés.
    \item L'\textbf{ouverture} : élargissement final de la réflexion vers une autre problématique, un autre domaine ou une actualité.
    \item \textbf{exploitée} (ou commentée) : la citation doit être expliquée et mise au service de l'argumentation. La règle complète est donc : une citation est \emph{intégrée} (syntaxe), \emph{signalée} (guillemets), \emph{référencée} (auteur, œuvre) et \emph{exploitée} (sens et portée explicités).
\end{enumerate}
\textbf{Point de vigilance :} ces quatre verbes (intégrer, signaler, référencer, exploiter) constituent la check-list des citations au Bac ; l'oubli de l'exploitation est le plus pénalisé.
\end{corrige}

\begin{corrige}[Exercice 4 — Quiz méthodologique chronométré]
\begin{enumerate}
    \item \textbf{F} — la thèse est la première réponse plausible, souvent celle du sens commun, mais pas « toujours » : certains sujets imposent de commencer par la position la plus évidente, quelle qu'elle soit.
    \item \textbf{F} — l'antithèse oppose une réponse contraire \emph{globale} et argumentée, non une réfutation mécanique point par point dans le même ordre.
    \item \textbf{F} — la synthèse est un dépassement constructif, non un résumé ni une conclusion partielle.
    \item \textbf{F} — on argumente objectivement (« nous », « on », tournures impersonnelles) ; « je pense que » ramène la démonstration au niveau de l'opinion.
    \item \textbf{V} — la citation orpheline, insérée sans lien syntaxique ni exploitation, est une erreur grave.
    \item \textbf{F} — l'accroche peut être un fait, une actualité, une expérience ; la citation de philosophe n'est qu'une option parmi d'autres.
    \item \textbf{V} — la problématique se formule comme une question et se termine par un point d'interrogation.
    \item \textbf{F} — on annonce la \emph{logique} du plan sans nommer « thèse, antithèse, synthèse ».
    \item \textbf{F} — une référence à l'actualité technique camerounaise (fibre optique, exploitation minière) est au contraire un excellent exemple au sens I-A-E.
    \item \textbf{F} — la conclusion reformule le bilan avec des mots nouveaux ; le copier-coller de la synthèse est sanctionné.
\end{enumerate}
\textbf{Barème indicatif :} 1 point par item ; objectif de maîtrise : 10/10 en moins de 10 minutes. En dessous de 7/10, reprendre les sections du cours sur l'introduction et la synthèse.
\end{corrige}

\courssub{L'introduction : la « porte d'entrée » (4 étapes canoniques)}

\begin{methode}[Les 4 étapes de l'introduction canonique (ordre immuable)]
\begin{enumerate}
    \item \textbf{Accroche} (1--2 phrases) : fait concret, daté, situé (actualité camerounaise, fait technique, référence culturelle) $\rightarrow$ thème général. \emph{Interdit :} « De tout temps... », énoncé du sujet, citation obligatoire.
    \item \textbf{Définition des termes} (2--3 termes clés) : définitions opératoires (précises pour le sujet), pas de dictionnaire. Mots-clés en \textbf{gras}.
    \item \textbf{Problématique} (1 phrase interrogative) : formule la tension dialectique entre deux thèses opposées. Termine par ?.
    \item \textbf{Annonce du plan} (1 phrase fluide) : annonce la logique (pas les étiquettes « thèse/antithèse/synthèse ») : « Nous examinerons d'abord que... Nous montrerons ensuite que... Nous comprendrons enfin que... »
\end{enumerate}
\end{methode}

\begin{attention}[Erreurs fatales en introduction]
\begin{itemize}
    \item Énoncer le sujet dans l'accroche (« Le sujet posé est... »).
    \item Donner sa réponse personnelle avant la problématique (« Je pense que l'ingénieur est responsable... »).
    \item Définir 5 termes ou des termes non-clés.
    \item Problématique non dialectique (question de cours, question fermée, absence de tension).
    \item Annonce du plan avec « Premièrement, Deuxièmement, Troisièmement » ou « Thèse, Antithèse, Synthèse ».
\end{itemize}
\end{attention}

\begin{corrige}[Exercice 6 — Rédaction d'une introduction canonique (4 étapes)]
\textbf{Corrigé type (introduction rédigée) :}\par
\emph{Accroche.} En juillet 2023, l'effondrement d'un immeuble en construction dans le quartier Makepe à Douala a coûté la vie à plusieurs personnes et relancé, dans le débat public camerounais, une question lancinante : qui doit répondre, moralement et non seulement devant les tribunaux, lorsque la technique défaille ?\par
\emph{Définition des termes.} L'\textbf{ingénieur} est le détenteur d'un savoir technique qui lui confère un pouvoir d'agir sur la nature et sur la vie des hommes. La \textbf{responsabilité} désigne l'obligation de répondre de ses actes et de leurs conséquences ; elle est juridique lorsqu'elle est sanctionnée par la loi, morale lorsqu'elle engage la conscience et la communauté. Enfin, \textbf{spécifique} signifie propre à une fonction : une responsabilité spécifique serait un devoir qui incombe à l'ingénieur en tant qu'ingénieur, et non simplement en tant que citoyen.\par
\emph{Problématique.} Or la technique passe souvent pour neutre : l'ingénieur n'exécuterait-il pas les décisions d'autrui — donneurs d'ordre, hiérarchie, marché — de sorte que la responsabilité morale se dissoudrait dans la chaîne de commandement ? Mais peut-on soutenir cette neutralité quand la compétence technique crée une asymétrie de savoir et une puissance de nuisance que le profane ne peut ni prévoir ni contrôler ? La fonction d'ingénieur fonde-t-elle, oui ou non, une responsabilité morale spécifique ?\par
\emph{Annonce du plan.} Nous examinerons d'abord la thèse de la neutralité technique, qui dilue la responsabilité de l'ingénieur dans celle de l'organisation. Nous montrerons ensuite que la puissance et le savoir propres à l'ingénieur fondent au contraire des devoirs moraux particuliers. Nous comprendrons enfin que cette responsabilité, partagée sans être diluée, appelle une déontologie professionnelle et une régulation collective.\par
\textbf{Barème indicatif (sur 5 pts) :} accroche concrète et camerounaise 1 pt ; définitions opératoires des 3 termes clés 1,5 pt ; problématique interrogative et dialectique 1,5 pt ; annonce du plan souple (sans « thèse/antithèse/synthèse ») 1 pt.
\textbf{Points de vigilance :} ne pas énoncer le sujet dans l'accroche ; ne pas donner sa réponse avant la problématique ; vérifier que chaque étape enchaîne logiquement avec la suivante.
\end{corrige}

\courssub{Le développement (1) : Thèse et Antithèse — L'art du paragraphe argumentatif (I-A-E)}

\begin{definition}[La structure I-A-E (Idée -- Argumentation -- Exemple)]
Chaque paragraphe du développement suit le modèle \cle{I-A-E} :
\begin{itemize}
    \item \textbf{I -- Idée directrice} (1 phrase) : thèse partielle défendue dans ce paragraphe, lien explicite avec la partie (I ou II).
    \item \textbf{A -- Argumentation} (3--5 lignes) : analyse logique, concepts, distinctions, référence à un auteur/œuvre (citation intégrée + exploitée).
    \item \textbf{E -- Exemple} (2--3 lignes) : illustration concrète, si possible camerounaise (technique, institutionnelle, historique), qui \emph{vérifie} l'argument sans le remplacer.
\end{itemize}
\emph{Transition :} une phrase de liaison entre paragraphes (au sein d'une partie) et entre parties (I $\rightarrow$ II, II $\rightarrow$ III).
\end{definition}

\begin{aretenir}[Exigences spécifiques Thèse vs Antithèse]
\begin{itemize}
    \item \textbf{Thèse (Partie I) :} position souvent « sens commun » ou « évidence première », mais \emph{argumentée rigoureusement} (pas d'opinion gratuite). Elle pose les fondements de la réponse « oui » (ou « non ») au problème.
    \item \textbf{Antithèse (Partie II) :} position \emph{contraire globale}, défendue \emph{sincèrement} (pas d'homme de paille). Elle ne réfute pas point par point la thèse, elle oppose une autre vision cohérente. Elle répond « non » (ou « oui ») au problème.
    \item \textbf{Connecteurs :} « Certes... », « Cependant... », « En revanche... », « Il faut pourtant reconnaître que... » pour passer de I à II.
\end{itemize}
\end{aretenir}

\begin{corrige}[Exercice 7 — Construction de paragraphes I-A-E : Thèse et Antithèse]
\textbf{1. Paragraphe de THÈSE (corrigé type).}\par
\emph{(I)} La puissance technique de l'ingénieur fonde une responsabilité morale spécifique : celui qui sait et qui peut doit prévenir. \emph{(A)} En effet, l'ingénieur occupe une position d'\textbf{asymétrie d'information} : il connaît les risques que le profane ignore — résistance d'un béton, vulnérabilité d'un réseau, toxicité d'un procédé. Or, dès lors qu'un agent connaît les conséquences possibles de ses actes et dispose du pouvoir de les modifier, il ne peut plus se réfugier derrière l'ignorance ou l'obéissance : le savoir crée un devoir. C'est le sens de l'impératif formulé par Hans Jonas dans \emph{Le Principe responsabilité} (1979) : « Agis de telle sorte que les effets de ton action soient compatibles avec la permanence d'une vie authentiquement humaine sur terre. » Cette maxime montre que la responsabilité croît avec la puissance technique : plus l'action porte loin dans l'espace et le temps, plus le devoir de prévention s'étend. \emph{(E)} Ainsi, l'ingénieur chargé du barrage de Lom Pangar sur la Sanaga ne se contente pas de calculer des structures : par ses choix de dimensionnement et de contrôle, il engage la sécurité des populations en aval et la disponibilité électrique de tout le pays ; nul autre acteur de la chaîne de décision ne possède cette compétence, donc nul autre ne peut porter exactement cette responsabilité.\par
\textbf{2. Paragraphe d'ANTITHÈSE (corrigé type).}\par
\emph{(I)} Pourtant, on peut soutenir que l'ingénieur n'a aucune responsabilité morale \emph{spécifique} : la décision technique est collective, subordonnée et juridiquement encadrée. \emph{(A)} D'abord, la responsabilité est \textbf{partagée et diluée} : un ouvrage résulte de la commande d'un client, du financement d'un État, du contrôle d'une administration ; imputer la faute à un seul maillon serait arbitraire. Ensuite, l'ingénieur est \textbf{hiérarchiquement subordonné} : il exécute un cahier des charges qu'il n'a pas fixé ; refuser la commande, c'est risquer son emploi, ce que nul ne peut exiger moralement de lui. Enfin, la technique serait \textbf{axiologiquement neutre} : comme le souligne Jacques Ellul dans \emph{La Technique ou l'Enjeu du siècle}, le « phénomène technique » obéit à sa propre logique d'efficacité, indifférente au bien et au mal ; ce qui relève de la morale, c'est l'\emph{usage}, donc la responsabilité du citoyen ou du politique, non celle du technicien. La responsabilité de l'ingénieur serait alors purement juridique (respect des normes), la morale restant l'affaire de tous également. \emph{(E)} Ainsi, dans un chantier de construction à Douala, si le promoteur impose des économies sur les matériaux et si le contrôle administratif fait défaut, l'ingénieur salarié apparaît comme un rouage parmi d'autres : lui imputer seul l'effondrement serait méconnaître la chaîne réelle des décisions.\par
\textbf{Barème indicatif (sur 8 pts, 4 par paragraphe) :} idée directrice explicite 1 pt ; argumentation logique en au moins deux étapes 1,5 pt ; auteur correctement cité et exploité 0,75 pt ; exemple technique camerounais pertinent et développé 0,75 pt.
\textbf{Points de vigilance :} l'antithèse doit être défendue \emph{sincèrement} (pas de caricature) ; la citation doit être intégrée à la phrase et suivie d'une phrase d'exploitation ; l'exemple illustre l'argument et ne le remplace pas.
\end{corrige}

\courssub{Le développement (2) : La Synthèse — Le dépassement dialectique}

\begin{propriete}[La Synthèse : 3 opérations obligatoires]
La troisième partie (Synthèse) n'est \textbf{pas un résumé}. Elle opère un \cle{dépassement dialectique} en trois temps :
\begin{enumerate}
    \item \textbf{Reconnaissance} : chaque thèse est partiellement vraie (dire \emph{quoi} exactement).
    \item \textbf{Critique} : chaque thèse est partielle/insuffisante car elle ignore l'apport de l'autre (dire \emph{pourquoi}).
    \item \textbf{Construction} : un concept nouveau, une distinction ou une articulation de niveau supérieur qui résout la contradiction (ex. : « responsabilité partagée mais asymétrique », « transparence raisonnée / redevabilité »).
\end{enumerate}
\emph{Sanction au Bac :} une synthèse qui résume = note plafonnée ; une synthèse qui dépasse = accès aux notes d'excellence.
\end{propriete}

\begin{methode}[Rédiger la Synthèse : phrases-types de dépassement]
\begin{itemize}
    \item « La thèse I a raison sur \emph{point X}, mais elle néglige \emph{point Y} que la thèse II met en lumière. »
    \item « La thèse II a raison sur \emph{point Y}, mais elle occulte \emph{point X} établi par la thèse I. »
    \item « Le dépassement consiste à montrer que \emph{concept nouveau} articule les deux : \emph{formulation précise}. »
    \item « On atteint ainsi un niveau supérieur : \emph{reformulation de la réponse au problème}. »
\end{itemize}
\end{methode}

\courssub{La conclusion : Bilan définitif et Ouverture maîtrisée}

\begin{methode}[La conclusion en 2 temps (ordre strict)]
\begin{enumerate}
    \item \textbf{Bilan définitif (Réponse) :} reformule la réponse apportée par la synthèse avec des mots \emph{nouveaux} (pas de copier-coller). Structure : « Au terme de cette réflexion, il apparaît que [réponse nuancée]... »
    \item \textbf{Ouverture maîtrisée :} pose \textbf{une} question nouvelle, liée au sujet, qui élargit l'horizon (autre domaine, actualité technique, prospective). \emph{Interdit :} y répondre, rouvrir le débat, banalité (« L'avenir nous le dira »), hors-sujet.
\end{enumerate}
\end{methode}

\begin{attention}[Pièges de la conclusion]
\begin{itemize}
    \item Copier-coller la dernière phrase de la synthèse.
    \item Introduire un nouvel argument.
    \item Faire une ouverture déconnectée (ex. : « Et si on parlait de la morale en politique ? » sur un sujet technique).
    \item Terminer par une phrase creuse (« Cette question reste ouverte »).
\end{itemize}
\end{attention}

\courssub{Les citations : Règles de choix, d'intégration et de référencement}

\begin{aretenir}[La check-list des 4 verbes (obligatoire au Bac)]
Toute citation doit être :
\begin{enumerate}
    \item \textbf{Intégrée} : syntaxiquement fondue dans la phrase (pas de citation « parachutée » entre deux points).
    \item \textbf{Signalée} : guillemets français « ... » (ou italiques pour citation longue > 3 lignes, centrée, sans guillemets).
    \item \textbf{Référencée} : Auteur, \emph{Œuvre} (année si connue). Si auteur incertain : « adage du milieu... », « formule inspirée de... ». \emph{Jamais d'invention.}
    \item \textbf{Exploitée} : phrase(s) suivante(s) expliquant \emph{en quoi} la citation soutient l'argument, ce qu'elle apporte (concept, distinction, autorité).
\end{enumerate}
\emph{Règle de traduction :} toute citation en langue étrangère \textbf{doit} être traduite en français.
\end{aretenir}

\begin{corrige}[Exercice 8 — Gestion des citations : Intégration, référencement, exploitation]
\textbf{Citation 1 (Ellul).} \emph{Intégration :} Contre l'idée d'une évolution technique fatale et autonome, Jacques Ellul rappelle dans \emph{La Technique ou l'Enjeu du siècle} que « la technique n'est pas un destin, c'est un choix ». \emph{Référencement :} Jacques Ellul, \emph{La Technique ou l'Enjeu du siècle}. \emph{Exploitation :} cette formule soutient l'argument de la responsabilité : si la technique résulte de choix humains et non d'une nécessité aveugle, alors ceux qui la conçoivent — les ingénieurs — peuvent et doivent répondre de ces choix ; la thèse de la neutralité s'effondre.\par
\textbf{Citation 2 (Jonas).} \emph{Intégration :} Hans Jonas fonde l'éthique de la technique sur un impératif nouveau : « Agis de telle sorte que les effets de ton action soient compatibles avec la permanence d'une vie authentiquement humaine sur terre » (\emph{Le Principe responsabilité}, 1979). \emph{Référencement :} Hans Jonas, \emph{Le Principe responsabilité}, 1979. \emph{Exploitation :} cet impératif montre que la responsabilité moderne porte sur des effets lointains (générations futures, environnement) que seul le technicien compétent peut anticiper ; il fonde donc le devoir de prévention spécifique de l'ingénieur.\par
\textbf{Citation 3 (inspirée de Hountondji).} \emph{Intégration :} Dans la lignée de la réflexion de Paulin Hountondji sur les rapports entre savoir et pouvoir en Afrique, on peut affirmer que « le savoir technique est un pouvoir ; tout pouvoir appelle une éthique ». \emph{Référencement :} formule inspirée de Paulin Hountondji (à présenter comme telle, car ce n'est pas une citation littérale). \emph{Exploitation :} elle articule savoir, pouvoir et devoir : l'asymétrie de compétence entre l'ingénieur et le public est un pouvoir, et tout pouvoir non contrôlé devient arbitraire — d'où la nécessité d'une déontologie.\par
\textbf{Citation 4 (adage).} \emph{Intégration :} Comme le rappelle un adage du milieu de l'ingénierie, « l'ingénieur est quelqu'un qui fait pour dix francs ce que n'importe quel imbécile ferait pour cent ». \emph{Référencement :} attribution incertaine — on écrit donc « adage du milieu de l'ingénierie » et non un nom d'auteur inventé. \emph{Exploitation :} cette boutade souligne que la compétence de l'ingénieur consiste en une optimisation que le profane ne peut ni juger ni imiter ; elle confirme l'asymétrie de savoir qui fonde sa responsabilité spécifique.\par
\textbf{Barème indicatif (sur 8 pts, 2 par citation) :} intégration syntaxique correcte 0,75 pt ; référencement honnête et précis 0,5 pt ; exploitation liée à l'argument 0,75 pt.
\textbf{Points de vigilance :} jamais de citation orpheline ; jamais d'attribution inventée (préférer « un adage », « une formule inspirée de... ») ; toute citation étrangère serait traduite.
\end{corrige}

\courssub{TP Consolidation : Simulation Bac \& Méthode d'auto-évaluation}

\begin{corrige}[Exercice 5 — Analyse de sujets : Décortiquer pour ne pas être hors sujet]
\textbf{Sujet A : « L'ingénieur a-t-il une responsabilité morale spécifique ? »}
\begin{enumerate}
    \item \textbf{Termes clés :} \emph{ingénieur} (détenteur d'un savoir technique et d'un pouvoir d'agir sur le monde), \emph{responsabilité morale} (obligation de répondre de ses actes devant sa conscience et la communauté, au-delà de la loi), \emph{spécifique} (propre à cette fonction, distincte de la responsabilité de tout citoyen).
    \item \textbf{Opposition fondamentale :} l'ingénieur est-il un simple exécutant neutre au service d'une commande, ou un acteur dont la puissance technique crée des devoirs moraux particuliers ?
    \item \textbf{Problématique :} « La compétence technique de l'ingénieur le tient-elle à l'écart de toute responsabilité morale propre, ou fonde-t-elle au contraire une responsabilité spécifique, distincte de celle du citoyen ordinaire ? »
    \item \textbf{Plan dialectique :} \textbf{I.} La neutralité de la technique : l'ingénieur comme exécutant (responsabilité diluée dans la chaîne de décision, subordination hiérarchique, responsabilité juridique $\neq$ morale). \textbf{II.} La puissance technique comme source de devoirs : asymétrie d'information, devoir de prévention (Jonas). \textbf{III.} Synthèse : une responsabilité partagée mais asymétrique, encadrée par la déontologie professionnelle (ONIGC) et la régulation étatique.
\end{enumerate}
\textbf{Sujet B : « La transparence totale est-elle une vertu technique ? »}
\begin{enumerate}
    \item \textbf{Termes clés :} \emph{transparence} (visibilité, intelligibilité, accessibilité des données et des procédés), \emph{vertu technique} (qualité qui rend un système efficace, fiable et digne de confiance).
    \item \textbf{Opposition fondamentale :} tout montrer garantit-il la confiance et le contrôle, ou certaines opacités (secret industriel, sécurité, vie privée) sont-elles conditions du bon fonctionnement technique ?
    \item \textbf{Problématique :} « La transparence totale renforce-t-elle la fiabilité et la légitimité des systèmes techniques, ou une opacité raisonnée est-elle indispensable à leur efficacité et à la protection des personnes ? »
    \item \textbf{Plan dialectique :} \textbf{I.} La transparence comme condition de confiance et de contrôle (données ouvertes, algorithmes explicables). \textbf{II.} Les limites de la transparence totale : secret industriel, sécurité nationale, protection des données (loi camerounaise de 2010). \textbf{III.} Synthèse : la transparence raisonnée ou \emph{redevabilité} (\emph{accountability}), vertu technique supérieure à la transparence brute.
\end{enumerate}
\textbf{Barème indicatif (par sujet) :} termes clés 1 pt ; opposition 1 pt ; problématique interrogative et dialectique 1 pt ; plan cohérent avec titres clairs 1 pt.
\end{corrige}

\begin{corrige}[Exercice 9 — Sujet type Bac technique : L'ingénieur et la responsabilité morale]
\textbf{1. Analyse du sujet (brouillon structuré).} \emph{Définitions :} ingénieur = détenteur d'un savoir technique doté d'un pouvoir d'agir ; responsabilité morale = obligation de répondre de ses actes devant la conscience et la communauté, distincte de la responsabilité juridique ; spécifique = attachée à la fonction elle-même. \emph{Présupposé :} il existe une morale applicable à l'activité technique (la technique n'est pas un espace hors-morale). \emph{Trois problématiques possibles :} (a) L'ingénieur est-il moralement responsable des usages de ses inventions ? (b) La responsabilité de l'ingénieur est-elle individuelle ou collective ? (c) La compétence technique fonde-t-elle une responsabilité morale propre à l'ingénieur ? \emph{Choix :} la (c), car elle épouse exactement les termes du sujet (« spécifique ») et autorise un plan dialectique net.\par
\textbf{2. Plan détaillé dialectique.}
\begin{itemize}
    \item \textbf{I. La neutralité du technicien : la responsabilité se dissout dans l'organisation.} Arguments : (a) subordination hiérarchique et division du travail (responsabilité diluée) ; (b) neutralité axiologique de la technique (Ellul : la technique obéit à sa logique d'efficacité) ; (c) confusion entre responsabilité juridique (normes) et morale. Exemples : chantier de construction à Douala piloté par le promoteur ; ingénieur salarié d'une compagnie minière à l'Est-Cameroun exécutant un cahier des charges.
    \item \textbf{II. La puissance technique fonde des devoirs propres.} Arguments : (a) asymétrie d'information (l'ingénieur seul connaît les risques) ; (b) Jonas, \emph{Le Principe responsabilité} : devoir de prévention proportionné à la puissance ; (c) Kant : l'homme comme fin — les usagers ne sont pas des moyens ; le serment et le code de déontologie de l'ONIGC institutionnalisent ce devoir. Exemples : dimensionnement du barrage de Lom Pangar ; sécurité des données dans les NTIC camerounaises.
    \item \textbf{III. Synthèse : une responsabilité partagée mais asymétrique.} Dépassement : la responsabilité collective (II contredit I sans l'annuler) ne supprime pas la responsabilité spécifique : chacun répond selon son pouvoir effectif ; l'ingénieur répond de ce que sa compétence lui permettait de prévoir et de refuser (devoir d'alerte, droit de retrait). Niveau supérieur : éthique de la responsabilité (Jonas) contre éthique de la conviction ; nécessité d'une gouvernance technique (régulation étatique, société civile, ONIGC).
\end{itemize}
\textbf{3. Introduction :} reprendre le corrigé type de l'exercice 6 (accroche Makepe/Douala, définitions, problématique, annonce souple).\par
\textbf{4. Synthèse rédigée (corrigé type).} L'opposition entre le technicien neutre et l'ingénieur responsable repose sur une alternative trop simple : tout ou rien. La première partie a montré, à juste titre, que la décision technique est collective ; mais la collectivité de la décision ne dilue la responsabilité que si l'on suppose tous les maillons égaux en savoir et en pouvoir — ce qui est faux. La deuxième partie a établi que la compétence crée une asymétrie ; il faut maintenant en tirer la conséquence exacte : la responsabilité est \emph{partagée mais asymétrique}. Chacun répond à la mesure de son pouvoir effectif : le promoteur de ses choix financiers, l'État de ses contrôles, et l'ingénieur de ce que sa compétence lui permettait de prévoir, de signaler et, le cas échéant, de refuser. C'est le sens du devoir d'alerte et du droit de retrait inscrits dans les codes de déontologie : l'ingénieur ne répond pas de tout, mais il répond de ce que lui seul pouvait savoir. Ainsi l'éthique de la responsabilité dépasse l'éthique de la conviction : il ne suffit pas d'avoir de bonnes intentions, il faut organiser collectivement la prévention — régulation étatique, ordre professionnel, vigilance de la société civile camerounaise. La responsabilité spécifique de l'ingénieur n'est donc ni totale ni nulle : elle est proportionnée à son savoir.\par
\textbf{5. Conclusion rédigée (corrigé type).} Au terme de cette réflexion, il apparaît que l'ingénieur possède bien une responsabilité morale spécifique : non parce qu'il serait le seul coupable des défaillances techniques, mais parce que son savoir crée une asymétrie — il sait ce que le profane ignore — et que de ce savoir naît un devoir de prévention, d'alerte et parfois de refus. La neutralité technique est un mythe commode ; la responsabilité totale serait une injustice. Reste la responsabilité proportionnée au pouvoir d'agir, encadrée par la déontologie et la régulation. Mais cette réponse ouvre une question nouvelle : lorsque la décision technique sera déléguée à des systèmes d'intelligence artificielle — algorithmes de diagnostic, conduite autonome —, à qui incombera la responsabilité morale : au concepteur, à l'utilisateur, ou à la machine elle-même ?\par
\textbf{Barème indicatif (sur 20) :} analyse du sujet et problématique 3 pts ; plan détaillé cohérent 4 pts ; introduction 4 pts ; synthèse (dépassement réel) 5 pts ; conclusion (bilan + réponse + ouverture) 4 pts.
\textbf{Points de vigilance :} la synthèse ne doit pas résumer mais dépasser (formule « partagée mais asymétrique ») ; l'ouverture ne doit pas être traitée ; au moins un exemple camerounais par partie ; aucune opinion personnelle.
\end{corrige}

\begin{corrige}[Exercice 10 — Sujet type Probatoire : Transparence, secret et technique]
\textbf{1. Distinction des deux sujets.} « Faut-il tout dire ? » relève de la morale interpersonnelle classique (sincérité, mensonge, devoir de vérité envers autrui). La variante technique déplace le débat vers l'\textbf{épistémologie} (qui peut comprendre un système complexe ? la transparence brute rend-elle intelligible ?) et vers l'\textbf{éthique professionnelle} (quels devoirs d'information pèsent sur les concepteurs de systèmes ?). L'enjeu n'est plus la sincérité entre personnes mais la visibilité des processus techniques pour la société.\par
\textbf{2. Problématique.} « La transparence totale des systèmes techniques — données, procédés, algorithmes — est-elle la condition de leur fiabilité et de la confiance publique, ou une opacité raisonnée est-elle indispensable à leur efficacité, à la sécurité et à la protection des personnes ? »\par
\textbf{3. Plan dialectique.}
\begin{itemize}
    \item \textbf{I. La transparence comme vertu :} elle rend possible le contrôle démocratique, la détection des erreurs et la confiance. Arguments : nul ne peut faire confiance à ce qu'il ne peut inspecter ; la reproductibilité est un critère de scientificité. Exemples : portail des données ouvertes du gouvernement camerounais ; exigence d'algorithmes explicables dans les services publics.
    \item \textbf{II. La nécessaire opacité :} la transparence totale détruit ce qu'elle prétend protéger. Arguments : secret industriel et innovation (brevets) ; sécurité (publier les vulnérabilités, c'est les offrir aux attaquants) ; protection de la vie privée (loi camerounaise n° 2010/012 sur la cybersécurité et la cybercriminalité) ; complexité irréductible (une donnée brute incompréhensible n'informe pas). Exemples : boîtes noires algorithmiques ; brevets pharmaceutiques.
    \item \textbf{III. Synthèse : la transparence raisonnée (redevabilité).} Ce qui est exigible n'est pas de tout montrer, mais de pouvoir \emph{rendre compte} : auditabilité par des autorités compétentes, explicabilité des décisions, protection des lanceurs d'alerte. La vertu technique n'est pas la transparence brute mais la \emph{redevabilité} (\emph{accountability}).
\end{itemize}
\textbf{4. Paragraphe I-A-E de THÈSE (corrigé type).} \emph{(I)} La transparence est d'abord une vertu technique parce qu'elle est la condition du contrôle, donc de la confiance. \emph{(A)} Un système technique n'est fiable que si ses défaillances peuvent être détectées et corrigées ; or la détection suppose la visibilité : ce qui est caché ne peut être ni vérifié ni amélioré. En outre, dans un État de droit, les décisions qui affectent les citoyens doivent pouvoir être discutées ; lorsque ces décisions sont prises par des systèmes techniques (attribution d'une bourse par un algorithme, traitement des données fiscales), l'opacité prive les citoyens de tout recours et fait de la technique un pouvoir sans légitimité. La transparence est ainsi la forme moderne du consentement éclairé. \emph{(E)} Ainsi, la publication des données budgétaires et des marchés publics au Cameroun permet aux citoyens et aux journalistes de suivre l'usage des fonds destinés aux infrastructures : sans cette visibilité, aucun contrôle — et donc aucune confiance durable — n'est possible.\par
\textbf{5. Paragraphe I-A-E d'ANTITHÈSE (corrigé type).} \emph{(I)} Pourtant, la transparence totale est un leurre : certaines opacités sont constitutives du bon fonctionnement technique et de la protection des personnes. \emph{(A)} D'abord, l'innovation suppose le secret : sans brevets ni secret industriel, nulle entreprise n'investirait dans la recherche, puisque ses résultats seraient immédiatement copiés. Ensuite, la sécurité exige l'opacité : publier les plans détaillés d'un réseau électrique ou les failles d'un système informatique, c'est les livrer aux malveillants. Enfin — et c'est décisif — la transparence brute ne produit pas l'intelligibilité : verser des millions de données brutes au public, c'est noyer l'information sous l'abondance ; une complexité technique irréductible exige une médiation experte. \emph{(E)} Ainsi, la loi camerounaise n° 2010/012 du 21 décembre 2010 relative à la cybersécurité et à la cybercriminalité protège à la fois les systèmes d'information et les données personnelles : elle reconnaît que la circulation non maîtrisée des données est un danger, non une vertu.\par
\textbf{6. Esquisse de la SYNTHÈSE.} Le dépassement consiste à distinguer \emph{transparence brute} (tout montrer à tous) et \emph{redevabilité} (pouvoir rendre compte devant une instance compétente). La première est impossible (complexité) et dangereuse (sécurité, vie privée) ; la seconde conserve l'exigence légitime de contrôle sous une forme réaliste : audits indépendants, explicabilité des algorithmes, protection des lanceurs d'alerte, accès régulé aux données. La vertu technique véritable n'est donc pas la transparence totale mais la transparence raisonnée : montrer ce qui doit l'être, à qui doit le voir, avec la médiation qui rend l'information intelligible.\par
\textbf{Barème indicatif (sur 20) :} distinction des sujets 2 pts ; problématique 3 pts ; plan 4 pts ; paragraphes I-A-E 4 pts chacun ; synthèse 3 pts.
\textbf{Points de vigilance :} ne pas traiter le sujet général « Faut-il tout dire ? » à la place de la variante technique ; citer correctement la loi de 2010 ; la synthèse doit introduire un concept nouveau (redevabilité), non un compromis mou.
\end{corrige}

\begin{corrige}[Exercice 11 — Réécriture critique : De la copie 8/20 à la copie 16/20]
\textbf{1. Nouvelle introduction.} Lorsqu'en juillet 2023 un immeuble en construction s'effondre à Makepe, à Douala, ensevelissant ses occupants, la première question des Camerounais n'est pas technique mais morale : qui doit répondre de ce drame ? L'ingénieur, détenteur d'un savoir technique qui lui donne le pouvoir de transformer la nature et d'engager la vie d'autrui, est aussitôt désigné. Mais que signifie « répondre » ? La responsabilité est l'obligation de répondre de ses actes et de leurs conséquences ; elle est morale lorsqu'elle engage la conscience et la communauté, et non seulement la loi. Dire cette responsabilité « spécifique », c'est affirmer qu'elle appartient à l'ingénieur en tant qu'ingénieur, et non comme simple citoyen. Or la technique ne passe-t-elle pas pour neutre, l'ingénieur n'étant que l'exécutant d'une commande dont la responsabilité remonte au donneur d'ordre ? Mais une compétence qui crée une asymétrie de savoir peut-elle vraiment dispenser de tout devoir propre ? Nous verrons d'abord que la responsabilité semble se dissoudre dans l'organisation collective du travail technique ; nous montrerons ensuite que la puissance propre de l'ingénieur fonde au contraire des devoirs qui n'appartiennent qu'à lui ; nous comprendrons enfin que cette responsabilité, partagée sans être diluée, est proportionnée au savoir et encadrée par la déontologie.\par
\textbf{2. Nouveau paragraphe de thèse (I-A-E).} \emph{(I)} L'ingénieur possède une responsabilité morale spécifique parce que son pouvoir technique engendre un devoir de prévention que nul autre ne peut exercer à sa place. \emph{(A)} L'ingénieur se trouve en situation d'asymétrie d'information : il connaît la résistance réelle d'un béton, la vulnérabilité d'un réseau, les marges de sécurité d'un ouvrage, là où le profane — usager, riverain, parfois même décideur — ne voit rien. Or celui qui sait et qui peut ne peut invoquer ni l'ignorance ni l'impuissance : le savoir crée le devoir. C'est précisément l'impératif que Hans Jonas formule dans \emph{Le Principe responsabilité} (1979) : « Agis de telle sorte que les effets de ton action soient compatibles avec la permanence d'une vie authentiquement humaine sur terre. » Cette maxime signifie que la responsabilité s'étend avec la puissance d'agir : plus les effets d'une décision technique portent loin, plus le devoir de les anticiper s'impose à celui qui seul peut les calculer. \emph{(E)} Ainsi, l'ingénieur qui dimensionne le barrage de Lom Pangar sur la Sanaga engage, par ses seuls calculs, la sécurité des populations en aval et l'approvisionnement électrique du pays : cette responsabilité-là, nul promoteur, nul ministre, nul citoyen ordinaire ne peut la porter à sa place, car nul autre ne possède la compétence qui la fonde.\par
\textbf{3. Nouvelle synthèse (partie III).} Il ne suffit pas de juxtaposer la neutralité du technicien et la culpabilité de l'expert : il faut comprendre pourquoi les deux thèses sont vraies chacune à moitié. La première a raison contre l'illusion du décideur solitaire : un ouvrage est le produit d'une chaîne — commande, financement, exécution, contrôle — et la responsabilité est donc réellement partagée. Mais « partagée » ne signifie pas « diluée jusqu'à disparaître » : la seconde partie a montré que les maillons de la chaîne ne sont pas égaux en savoir. La synthèse s'impose alors : la responsabilité est \emph{partagée mais asymétrique} — chacun répond à la mesure de son pouvoir effectif, et l'ingénieur répond de ce que sa compétence lui permettait de prévoir, de signaler et de refuser. C'est pourquoi le code de déontologie de l'Ordre National des Ingénieurs de Génie Civil du Cameroun consacre un devoir d'alerte : il transforme en règle professionnelle ce que l'éthique de Jonas pose comme exigence. Mais une éthique individuelle ne suffit pas : la responsabilité spécifique de l'ingénieur appelle une gouvernance collective de la technique — contrôle de l'État, vigilance de la société civile, sanctions de l'ordre professionnel. On atteint ainsi le niveau supérieur de compréhension : la question n'est plus « l'ingénieur est-il responsable ? » mais « comment organiser une responsabilité proportionnée aux savoirs dans une société technique ? ».\par
\textbf{4. Nouvelle conclusion.} Nous sommes partis du silence embarrassé qui suit les effondrements d'immeubles : qui doit répondre ? L'analyse a montré que la réponse ne peut être ni « tout le monde » — ce qui signifierait personne — ni « l'ingénieur seul » — ce qui serait injuste. L'ingénieur possède bien une responsabilité morale spécifique, fondée sur l'asymétrie du savoir et la puissance d'agir : il répond de ce que lui seul pouvait prévoir, signaler et refuser. Cette responsabilité, proportionnée et partagée, trouve sa traduction dans la déontologie professionnelle et dans la régulation collective. Reste qu'elle suppose un agent capable de savoir et de choisir : que deviendra-t-elle lorsque les calculs et les décisions seront confiés à des intelligences artificielles dont nul humain ne maîtrisera plus entièrement les raisonnements ?\par
\textbf{Barème indicatif (sur 20) :} introduction 5 pts (4 étapes respectées) ; paragraphe de thèse 5 pts (I-A-E complet, citation intégrée et exploitée, exemple développé) ; synthèse 6 pts (dépassement réel, concept nouveau, niveau supérieur) ; conclusion 4 pts (bilan, réponse, ouverture non traitée).
\textbf{Points de vigilance :} supprimer toute formule clichée (« De tout temps... ») ; ne jamais laisser une citation orpheline ; l'ouverture sur les robots de la copie d'origine était correcte dans son principe mais trop brève et déconnectée — la relier explicitement au sujet (la responsabilité) comme dans le corrigé.
\end{corrige}

\begin{experience}[TP Final : Simulation d'examen (3h) — Auto-évaluation]
\textbf{Consigne :} Traiter l'un des deux sujets ci-dessous en respectant strictement la méthode (brouillon structuré + rédaction soignée). Temps imparti : 3h.
\begin{enumerate}
    \item « La technique libère-t-elle l'homme ou l'aliène-t-elle ? » (Sujet type Bac — thèse/antithèse/synthèse sur l'ambivalence technique).
    \item « L'innovation technique est-elle toujours un progrès ? » (Sujet type Probatoire — distinction innovation/progrès, critères de l'évaluation).
\end{enumerate}
\textbf{Grille d'auto-évaluation (à remplir après rédaction) :}
\begin{center}
\begin{tabularx}{\linewidth}{|l|X|c|}\hline
\textbf{Critère} & \textbf{Indicateurs de réussite} & \textbf{Note /5} \\ \hline
Introduction & 4 étapes respectées, accroche camerounaise, définitions opératoires, problématique dialectique, annonce souple & \\ \hline
Développement I-A-E & Idées directrices claires, arguments logiques, auteurs cités/exploités, exemples camerounais concrets & \\ \hline
Synthèse & Dépassement réel (pas résumé), concept nouveau, articulation thèse/antithèse, niveau supérieur & \\ \hline
Conclusion & Bilan reformulé, réponse nette, ouverture pertinente non traitée & \\ \hline
Citations & 4 verbes respectés (intégrer, signaler, référencer, exploiter), pas d'orphan, traduction si besoin & \\ \hline
Langage/Présentation & « Nous »/impersonnel, pas d'argot, paragraphes aérés, ratures propres, lisibilité & \\ \hline
\end{tabularx}
\end{center}
\textbf{Objectif :} 16/20 minimum. Si note < 12 : relire les \texttt{aretenir} et \texttt{methode} des sections 1 à 6, refaire l'exercice 11.
\end{experience}

\newpage

\coursec{Fiche bilan}

\begin{popfigure}
\begin{center}
\begin{tikzpicture}[mindmap, grow cyclic, every node/.style={concept, execute at begin node=\hspace{0pt}}, root concept/.append style={concept color=popblue, font=\small\bfseries, text=white}, level 1/.append style={level distance=3.6cm, sibling angle=51.4}, level 2/.append style={level distance=2.2cm, sibling angle=45, font=\scriptsize}]
\node[root concept]{La dissertation philosophique}
  child[concept color=poporange]{ node[align=center]{Introduction\\ (4 étapes)}
    child { node {Accroche} }
    child { node {Problématique} }
    child { node {Annonce du plan} }
  }
  child[concept color=popgreen]{ node {Thèse}
    child { node {I--A--E} }
    child { node {1 idée = 1 §} }
  }
  child[concept color=poppurple]{ node {Antithèse}
    child { node {Objection} }
    child { node {Critique} }
  }
  child[concept color=popgold]{ node {Synthèse}
    child { node {Dépassement} }
    child { node {Réponse nuancée} }
  }
  child[concept color=poppink]{ node {Conclusion}
    child { node {Bilan} }
    child { node {Ouverture} }
  }
  child[concept color=popblue!70]{ node {Citations}
    child { node {Pertinence} }
    child { node {Référence} }
  }
  child[concept color=popmuted]{ node {Méthode Bac}
    child { node {Gestion du temps} }
    child { node {Auto-évaluation} }
  };
\end{tikzpicture}
\end{center}
\legende{Carte mentale de la leçon : les sept piliers de la dissertation philosophique au Baccalauréat technique.}
\end{popfigure}

\begin{bilan}
\textbf{Définitions clés}
\begin{itemize}
  \item \cle{Dissertation philosophique} : exercice de réflexion argumentée qui répond à une question par une démarche ordonnée (introduction, développement, conclusion).
  \item \cle{Problématique} : question centrale qui révèle la tension ou le paradoxe contenu dans le sujet.
  \item \cle{Thèse} : première réponse argumentée au problème posé.
  \item \cle{Antithèse} : objection ou réponse contraire qui met la thèse à l'épreuve.
  \item \cle{Synthèse} : dépassement dialectique qui réconcilie ou dépasse thèse et antithèse.
  \item \cle{Citation} : parole d'auteur rapportée pour appuyer (et non remplacer) l'argumentation.
\end{itemize}

\textbf{L'architecture à connaître par cœur}
\begin{center}
\begin{tabularx}{\linewidth}{|l|X|}
\hline
\rowcolor{popblueL} \textbf{Partie} & \textbf{Contenu exigé} \\
\hline
Introduction & Accroche ; présentation du sujet et définition des termes ; problématique ; annonce du plan. \\
\hline
Développement I & Thèse : argument(s), exemple(s), référence(s). \\
\hline
Développement II & Antithèse : objection argumentée, limite de la thèse. \\
\hline
Développement III & Synthèse : dépassement, réponse personnelle nuancée. \\
\hline
Conclusion & Bilan du parcours, réponse claire à la problématique, ouverture. \\
\hline
\end{tabularx}
\end{center}

\textbf{Méthodes express}
\begin{itemize}
  \item Paragraphe argumentatif : schéma \cle{I--A--E} : \textbf{I}dée (énoncée), \textbf{A}rgument (démontrée), \textbf{E}xemple (illustrée).
  \item Une seule idée par paragraphe ; transitions explicites entre les parties.
  \item Citation : la choisir \emph{pertinente}, la citer \emph{courte}, l'\emph{expliquer} et nommer l'auteur.
  \item Gestion du temps : environ 30 min d'analyse et de plan, 2 h de rédaction, 15 min de relecture.
\end{itemize}
\end{bilan}

\begin{aretenir}[Checklist avant le Bac]
Avant de rendre ma copie, je vérifie :
\begin{itemize}
  \item[$\square$] J'ai défini les termes essentiels du sujet dans l'introduction.
  \item[$\square$] Ma problématique est une véritable question, clairement formulée.
  \item[$\square$] Mon plan est annoncé à la fin de l'introduction.
  \item[$\square$] Chaque paragraphe suit le schéma Idée -- Argument -- Exemple.
  \item[$\square$] Ma thèse et mon antithèse sont réellement opposées.
  \item[$\square$] Ma synthèse dépasse l'opposition au lieu de la répéter.
  \item[$\square$] Mes citations sont courtes, exactes et attribuées à leur auteur.
  \item[$\square$] Je n'ai pas plaqué un cours : tout sert la réponse au sujet.
  \item[$\square$] Ma conclusion répond explicitement à la problématique.
  \item[$\square$] Mon ouverture est liée au sujet et n'introduit pas un nouveau problème.
  \item[$\square$] J'ai relu l'orthographe, la syntaxe et les transitions.
\end{itemize}
\end{aretenir}

\findecours

\end{document}
